% concatenated from MAIN_FILE.tex
\documentclass[12pt]{amsart}
\usepackage{amscd,amsmath,amsthm,amssymb}
\usepackage{mathtools}
\usepackage[margin=2cm]{geometry}
\usepackage{epsfig}
\usepackage{rawfonts}
\usepackage{enumerate}
\usepackage{graphics}
\usepackage{multirow}
\usepackage{xspace}
\usepackage{graphicx}
\usepackage{mathrsfs}
\usepackage{amsmath}
\usepackage{amsfonts}
\usepackage{amssymb}
\usepackage{amsthm}
\usepackage{graphicx}
\usepackage{booktabs}
\usepackage{caption}
\usepackage{listings}
\usepackage{setspace}
\usepackage[mathscr]{eucal}
\usepackage{pgfplots}
\usepackage{hyperref}
\usepackage{wrapfig}
\usepackage{floatflt,epsfig}
\usepackage{ dsfont }
\usepackage{amscd}
\usepackage{tikz-cd}
\usepackage{fancyhdr}
\usepackage[all]{xy}
\usepackage{latexsym}
\usepackage{amscd}
\usepackage{pifont}
\usepackage{subfig}
\usepackage{easyReview}
\usepackage{subfig}
\usepackage{pstricks-add}
\usepackage{pgf,tikz,pgfplots}
\pgfplotsset{compat=1.15}
\usepackage{mathrsfs}
\usetikzlibrary{arrows}
\usetikzlibrary[patterns]
%\usepackage{subcaption}
 \usepackage[normalem]{ulem}
 \usepackage{multicol}

 
%
%------    GENERAL MACROS    -----
%
% Standard rings and fields, affine and projective space
%
\def\NZQ{\Bbb}               % the font for N,Z,Q,R,C
\def\NN{{\NZQ N}}
\def\QQ{{\NZQ Q}}
\def\ZZ{{\NZQ Z}}
\def\RR{{\NZQ R}}
\def\CC{{\NZQ C}}
\def\AA{{\NZQ A}}
\def\PP{{\NZQ P}}
\def\DD{{\NZQ D}}
\def\FF{{\NZQ F}}
\def\GG{{\NZQ G}}
\def\HH{{\NZQ H}}
\def\EE{{\NZQ E}}
\def\KK{{\NZQ K}}
\def\MM{{\NZQ M}}
%
%------------------------------------------------
% Symbols in "Fraktur"
%
\def\frk{\frak}               % font for "Fraktur"
\def\aa{{\frk a}}
\def\pp{{\frk p}}
\def\Pp{{\frk P}}
\def\qq{{\frk q}}
\def\Qq{{\frk Q}}
\def\mm{{\frk m}}
\def\Mm{{\frk M}}
\def\nn{{\frk n}}
%\def\Phi{{\frk N}}
%
%------------------------------------------------
\newcommand{\Ac}{\mathcal{A}}
\newcommand{\Bc}{\mathcal{B}}
\newcommand{\Cc}{\mathcal{C}}
\newcommand{\Dc}{\mathcal{D}}

\newcommand{\cA}{\mathcal{A}}
\newcommand{\cB}{\mathcal{B}}
\newcommand{\cC}{\mathcal{C}}
\newcommand{\cD}{\mathcal{D}}
\newcommand{\Ec}{\mathcal{E}}
\newcommand{\Ob}{\mathcal{O}}
\newcommand{\cP}{\mathcal{P}}
\newcommand{\Pc}{\mathcal{P}}
\newcommand{\Qc}{\mathcal{Q}}
\newcommand{\cH}{\mathcal{H}}
\newcommand{\cI}{\mathcal{I}}
\newcommand{\cK}{\mathcal{K}}
\newcommand{\Lc}{\mathcal{L}}
\newcommand{\cM}{\mathcal{M}}
\newcommand{\cF}{\mathcal{F}}
\newcommand{\cU}{\mathcal{U}}
\newcommand{\cQ}{\mathcal{Q}}
\newcommand{\cW}{\mathcal{W}}
\newcommand{\cR}{\mathcal{R}}
\newcommand{\Rc}{\mathcal{R}}
\newcommand{\cS}{\mathcal{S}}
\newcommand{\cT}{\mathcal{T}}
\newcommand{\cL}{\mathcal{L}}
\newcommand{\cV}{\mathcal{V}}
\newcommand{\cG}{\mathcal{G}}
\newcommand{\rHP}{\mathrm{HP}}
\newcommand{\sF}{\mathscr{F}}
\newcommand{\sH}{\mathscr{H}}
\newcommand{\fR}{\mathfrak{R}}

\renewcommand{\qedsymbol}{$\square$}
\newcommand{\rev}{\mathrm{rev}}


% Small letters in bold
%
\def\ab{{\bold a}}
\def\bb{{\bold b}}
\def\xb{{\bold x}}
\def\yb{{\bold y}}
\def\zb{{\bold z}}
\def\cb{{\bold c}}
\def\db{{\bold d}}
\def\eb{{\bold e}}
%
\def\opn#1#2{\def#1{\operatorname{#2}}} % to make operators
%------------------------------------------------
% Numerical invariants of rings, ideals, and modules
%
\opn\chara{char} \opn\length{\ell} \opn\pd{pd} \opn\rk{rk}
\opn\projdim{proj\,dim} \opn\injdim{inj\,dim} \opn\rank{rank}
\opn\depth{depth} \opn\grade{grade} \opn\height{height}
\opn\embdim{emb\,dim} \opn\codim{codim}
\def\OO{{\mathcal O}}
\opn\Tr{Tr} \opn\bigrank{big\,rank}
\opn\superheight{superheight}\opn\lcm{lcm}
\opn\trdeg{tr\,deg}%\emph{
	\opn\reg{reg} \opn\lreg{lreg} \opn\ini{in} \opn\lpd{lpd}
	\opn\size{size} \opn\sdepth{sdepth}
	\opn\link{link}\opn\fdepth{fdepth}\opn\lex{lex}\opn\dist{dist}
	%------------------------------------------------
	% Divisors
	%
	\opn\div{div} \opn\Div{Div} \opn\cl{cl} \opn\Cl{Cl}
	%
	%------------------------------------------------
	% Subsets of the spectrum of a ring
	%
	\opn\Spec{Spec} \opn\Supp{Supp} \opn\supp{supp} \opn\Sing{Sing}
	\opn\Ass{Ass} \opn\Min{Min}\opn\Mon{Mon}
	%
	%------------------------------------------------
	% Standard operations on ideals and modules
	%
	\opn\Ann{Ann} \opn\Rad{Rad} \opn\Soc{Soc}
	%
	%------------------------------------------------
	% Linear algebra and homology, endo- and automorphisms
	%
	\opn\Im{Im} \opn\Ker{Ker} \opn\Coker{Coker} \opn\Am{Am}
	\opn\Hom{Hom} \opn\Tor{Tor} \opn\Ext{Ext} \opn\End{End}
	\opn\Aut{Aut} \opn\id{id}
	\def\Frob{{\mathcal F}}
	\opn\nat{nat}
	\opn\pff{pf}%   \pf exists already
	\opn\Pf{Pf} \opn\GL{GL} \opn\SL{SL} \opn\mod{mod} \opn\ord{ord}
	\opn\Gin{Gin} \opn\Hilb{Hilb}\opn\sort{sort}
	%
	%------------------------------------------------
	% Convexity
	%
	\opn\aff{aff} \opn
	\con{conv} \opn\relint{relint} \opn\st{st}
	\opn\lk{lk} \opn\cn{cn} \opn\core{core} \opn\vol{vol}
	\opn\link{link} \opn\star{star}\opn\lex{lex}\opn\set{set}
	%------------------------------------------------
	% Graded rings and Rees algebras
	\opn\gr{gr}
	\def\Rees{{\mathcal R}}
	%
	%------------------------------------------------
	% Polynomials and power series
	%
	\def\poly#1#2#3{#1[#2_1,\dots,#2_{#3}]}
	\def\pot#1#2{#1[\kern-0.28ex[#2]\kern-0.28ex]}
	\def\Pot#1#2#3{\pot{#1}{#2_1,\dots,#2_{#3}}}
	\def\konv#1#2{#1\langle #2\rangle}
	\def\Konv#1#2#3{\konv{#1}{#2_1,\dots,#2_{#3}}}
	%
	%------------------------------------------------
	% Direct and inverse limits
	%
	\opn\dirlim{\underrightarrow{\lim}}
	\opn\inivlim{\underleftarrow{\lim}}
	%
	%
	% Names with a meaning
	%
	\let\union=\cup
	\let\sect=\cap
	\let\dirsum=\oplus
	\let\tensor=\otimes
	\let\iso=\cong
	\let\Union=\bigcup
	\let\Sect=\bigcap
	\let\Dirsum=\bigoplus
	\let\Tensor=\bigotimes
	%
	%------------------------------------------------
	%
	\let\to=\rightarrow
	\let\To=\longrightarrow
	\def\Implies{\ifmmode\Longrightarrow \else
		\unskip${}\Longrightarrow{}$\ignorespaces\fi}
	\def\implies{\ifmmode\Rightarrow \else
		\unskip${}\Rightarrow{}$\ignorespaces\fi}
	\def\iff{\ifmmode\Longleftrightarrow \else
		\unskip${}\Longleftrightarrow{}$\ignorespaces\fi}
	\let\gets=\leftarrow
	\let\Gets=\longleftarrow
	\let\followsfrom=\Leftarrow
	\let\Followsfrom=\Longleftarrow
	\let\:=\colon
	%
	%
	%
	%We like the var forms of some greek letters (as taught in German schools)
	\let\epsilon\varepsilon
	%\let\phi=\varphi
	\let\kappa=\varkappa
 %
	% %           We print on A4 paper
	% %
	% \textwidth=15cm \textheight=22cm \topmargin=0.5cm
	% \oddsidemargin=0.5cm \evensidemargin=0.5cm \pagestyle{plain}
        % %
	% %           The pf environment of AMSART needs a little help
	% %
	\def\qed{\ifhmode\textqed\fi
		\ifmmode\ifinner\quad\qedsymbol\else\dispqed\fi\fi}
	\def\textqed{\unskip\nobreak\penalty50
		\hskip2em\hbox{}\nobreak\hfil\qedsymbol
		\parfillskip=0pt \finalhyphendemerits=0}
	\def\dispqed{\rlap{\qquad\qedsymbol}}
	\def\noqed{\def\qed{\relax}}
	%
	% ------    END OF GENERAL MACROS    -------
	\opn\dis{dis}
	\def\pnt{{\raise0.5mm\hbox{\large\bf.}}}
	\def\lpnt{{\hbox{\large\bf.}}}
	\opn\Lex{Lex}
	
	%-- macro for local cohomology-----------------------------
	\def\Coh#1#2{H_{\mm}^{#1}(#2)}
	%-- macro for a complicated condition for the extended
	%-- Hochster's formula
	\def\hchst#1{for all $u\in G(I)$ there  exists $i\notin {#1}$
		such that $\nu_i(u) > a_i \geq 0$}	



    %%%%%%%%%%%%% Enumeration Theorems, Propositions, and so on %%%%%%%%%%%%%%%%%%%%%%%%%%
        
        \newtheorem{Theorem}{Theorem}[section]
	\newtheorem{Lemma}[Theorem]{Lemma}
	\newtheorem{Corollary}[Theorem]{Corollary}
	\newtheorem{Proposition}[Theorem]{Proposition}
	\newtheorem{Remark}[Theorem]{Remark}
	\newtheorem{Remarks}[Theorem]{Remarks}
	\newtheorem{Example}[Theorem]{Example}
	\newtheorem{Examples}[Theorem]{Examples}
	\newtheorem{Definition}[Theorem]{Definition}
	\newtheorem{Problem}[Theorem]{}
	\newtheorem{Conjecture}[Theorem]{Conjecture}
	\newtheorem{Question}[Theorem]{Question}
        \newtheorem{Discussion}[Theorem]{Discussion}
        \newtheorem{Notation}[Theorem]{Notation}
                \newtheorem{Assumption}[Theorem]{Assumption}
       \newtheorem{Setup}[Theorem]{Setup}
       \newtheorem{Construction}[Theorem]{Construction}


        
\begin{document}

    \title[TBA]{TBA}


\title[On the Linearity of Squarefree Powers of Edge Ideals]{On the Linearity of Squarefree Powers of Edge Ideals}

%\title[Linear Resolutions of Squarefree Powers of Edge Ideals]{Linear Resolutions of Squarefree Powers of Edge Ideals}

%\title[Linearly Related Squarefree Powers of Edge Ideals and Stanley--Reisner Ideals]{Linearly Related Squarefree Powers of Edge Ideals and Stanley--Reisner Ideals}

%\title[Squarefree Powers of Edge Ideals: Linearity and Combinatorial Obstructions]{Squarefree Powers of Edge Ideals: Linearity and Combinatorial Obstructions}

%\title[On the Linearity of Squarefree Powers of Edge Ideals]{On the Linearity of Squarefree Powers of Edge Ideals}

%\title[Squarefree Powers of Stanley--Reisner Ideals of Flag Complexes]{Squarefree Powers of Stanley--Reisner Ideals of Flag Complexes}

%\title[Combinatorial Characterizations of Linear Squarefree Powers of Edge Ideals]{Combinatorial Characterizations of Linear Squarefree Powers of Edge Ideals}

%\title[Linear Syzygies of Squarefree Powers of Edge Ideals]{Linear Syzygies of Squarefree Powers of Edge Ideals}

%\title[Squarefree Powers of Edge Ideals and Induced Bipartite Subgraphs]{Squarefree Powers of Edge Ideals and Induced Bipartite Subgraphs}

%\title[Linearity Properties of Squarefree Powers of Edge Ideals]{Linearity Properties of Squarefree Powers of Edge Ideals}
	
    

\author[F. Navarra]{Francesco Navarra}      
\address[Francesco Navarra]{Sabanci University, Faculty of Engineering and Natural Sciences, Orta Mahalle, Tuzla 34956, Istanbul, Turkey}	
\email{francesco.navarra@sabanciuniv.edu}

\author[A. A. Qureshi]{Ayesha Asloob Qureshi}      
\address[Ayesha Asloob Qureshi]{Sabanci University, Faculty of Engineering and Natural Sciences, Orta Mahalle, Tuzla 34956, Istanbul, Turkey}	
\email{aqureshi@sabanciuniv.edu, ayesha.asloob@sabanciuniv.edu}

\author[N. Terai]{Naoki Terai}      
\address[Naoki Terai]{Department of Mathematics, Okayama University, 3-1-1 Tsushima-naka, Kita-ku, Okayama 700-8530, Japan}	
\email{terai@okayama-u.ac.jp}
        
\keywords{Simplicial complex, Stanley-Reisner ideal, square-free power, linearly related, linear resolution.}
	
\subjclass[2020]{13D02, 13F55, 05E40}
    
\maketitle
      
\begin{abstract}
Let $G$ be a graph and $I(G)$ its edge ideal. The $p$-th squarefree power $I(G)^{[p]}$ is the monomial ideal generated by squarefree monomials corresponding to the matchings of size $p$ of $G$. In this paper, we provide a combinatorial characterization of when $I(G)^{[p]}$ is linearly related, i.e., when its first syzygy module is generated by linear forms. Moreover, for a $1$-dimensional flag simplicial complex $\Delta$ and its Stanley-Reisner ideal $I_{\Delta}$, which arises as the edge ideal of the complement graph of $\Delta$, we describe the shape of the Betti table of $I_{\Delta}^{[p]}$ and we give a combinatorial characterization of when $I_{\Delta}^{[p]}$ has a linear resolution.
\end{abstract}
%A fundamental aspect of the subject is the deep interplay between the algebraic and homological properties of squarefree monomial ideal and the combinatorial structure of the object from which it arises.
%In particular, squarefree monomial ideals admit natural interpretations in terms of simplicial complexes and graphs, thereby allowing one to investigate their algebraic invariants through combinatorial and topological techniques.
%More precisely, to every squarefree monomial ideal $I$ one can associate both a Stanley--Reisner simplicial complex, whose faces correspond to the squarefree monomials not belonging to $I$, and a facet complex, whose facets correspond to the minimal monomial generators of $I$. %Equivalently, the facet complex may be regarded as a hypergraph whose edge ideal is precisely $I$.
%These correspondences have proved particularly fruitful, since several important algebraic properties of monomial ideals admit elegant combinatorial interpretations in terms of the associated simplicial complexes. 
% Recall that a homogeneous ideal $I$ generated in degree $d$ is said to have a linear resolution if its minimal graded free resolution has the form
% $$
% 0 \to \bigoplus S(-(d+p))^{\beta_p}
% \to \cdots \to
% \bigoplus S(-(d+1))^{\beta_1}
% \to S(-d)^{\beta_0}
% \to I \to 0,
% $$
% that is, if all syzygies occur in the smallest possible degrees. 
%In general, however, it is unrealistic to expect that the existence of a linear resolution can always be characterized purely in combinatorial terms, since this property may depend on the characteristic of the base field. A classical example is provided by triangulations of the real projective plane. Nevertheless, one of the remarkable aspects of the results obtained in the present paper is that they are completely characteristic-free.


\section*{Introduction}

The study of the algebraic and homological properties of squarefree monomial ideals through combinatorial and topological methods has been a central theme in Combinatorial Commutative Algebra over the last few decades. Some of the most standard and widely read references on this subject include \cite{BH, kozlov2007combinatorial, miller2004combinatorial, Villarreal}. 

In this context, Villarreal introduced the \textit{edge ideal} of a graph $G$ at the end of the 1980s \cite{Villarreal1990}. This ideal, denoted by $I(G)$, is defined as the squarefree monomial ideal generated by the monomials corresponding to the edges of $G$. Since then, many algebraic properties of $I(G)$ have been studied in terms of the structure of the graph $G$. One of the foundational results in this direction is due to Fr\"oberg~\cite{Froberg1990}, who proved that the edge ideal of a graph $G$ has a linear resolution if and only if the complement of $G$ is chordal.
%Subsequently, Eagon and Reiner~\cite{EAGON1998265} established a duality between the existence of a linear resolution and the Cohen--Macaulay property. More precisely, they showed that a squarefree monomial ideal has a linear resolution if and only if the Stanley--Reisner ring associated with its Alexander dual is Cohen--Macaulay.

In this paper, we investigate the minimal graded free resolution of squarefree powers of edge ideals. Given a monomial ideal $I$, its $k$-th squarefree power, denoted by $I^{[k]}$, is the ideal generated by the squarefree monomials among the generators of $I^k$. The systematic study of squarefree powers was initiated in~\cite{BHZ} in the context of edge ideals of graphs, although such ideals had already appeared implicitly in earlier works, as in the proof of Proposition~2.10 in~\cite{JZ}. Since then, squarefree powers have been extensively investigated; see \cite{CFL, das2026regularitysquarefreepowersedge, EHHM2, EH, Kamalesh1, KAMBERI2026240, tyoungcomplexessquarefreepowers, SEYEDFAKHARI2024107488}, to name just a few.

The study of squarefree powers $I^{[k]}$ is motivated both by their algebraic relation to ordinary powers and by their combinatorial properties. The multigraded minimal free resolution of $I^{[k]}$ appears as a subcomplex of the resolution of $I^k$ (see \cite[Lemma~4.4]{HerzogHibiZheng}), which implies that $\reg(I^{[k]}) \leq \reg(I^k)$. Hence, if $I^{[k]}$ does not admit a linear resolution, neither does $I^k$. Furthermore, squarefree powers naturally encode the matching theory of simplicial complexes. If $\Delta$ is a simplicial complex and $I(\Delta)$ is its facet ideal, the minimal generators of $I(\Delta)^{[k]}$ correspond precisely to the matchings of $\Delta$ of size $k$. In particular, $I(\Delta)^{[k]} \neq 0$ if and only if $1 \leq k \leq \nu(\Delta)$, where $\nu(\Delta)$ is the matching number of $\Delta$.

A fundamental result due to Bigdeli, Herzog, and Zaare-Nahandi~\cite[Theorem~5.1]{BHZ} stated that the highest non-vanishing squarefree power of the edge ideal of a graph has linear quotients and therefore admits a linear resolution. Motivated by this phenomenon, the authors conjectured in~\cite[Conjecture~5.3]{BHZ} that, for a graph $G$, the squarefree power $I(G)^{[p]}$ is linearly related if and only if $p \geq \nu_0(G),$ where $\nu_0(G)$ denotes the restricted matching number of $G$. %Recall that a matching is said to be restricted if one of its edges forms a gap with every other edge of the matching. 
However, a counterexample to this conjecture was later provided in~\cite[page~12]{EHHM}. This naturally motivates the following problem.

\medskip

\noindent
\textbf{Problem.}
Characterize the graphs $G$ for which the squarefree power $I(G)^{[p]}$ is linearly related.

\medskip

One of the principal goals of this paper is to provide a complete answer to this question. Our first main result is the following theorem.

\begin{Theorem}[Theorem~\ref{Thm:linrelated1}]
Let $G$ be a graph and let $2 \leq p \leq \nu(G)$. Then $I(G)^{[p]}$ is linearly related if and only if there does not exist an induced subgraph $H$ of $G$ on $2p+2$ vertices such that $H$ is disconnected and $\nu(H)=p+1$.
\end{Theorem}

%The previous theorem completely settles the problem arising from~\cite[Conjecture~5.3]{BHZ}.% It also reveals that the obstruction to the linearity of the first syzygy module is governed by the existence of very precise induced disconnected subgraphs.

The problem of characterizing when squarefree powers of edge ideals admit a linear resolution is significantly more challenging. In this paper, we address this question for edge ideals of complements of triangle-free graphs (see Theorem~\ref{Thm: characterization linear resolution graph}) or, equivalently, for Stanley--Reisner ideals of $1$-dimensional flag simplicial complexes. Our main result is formulated as follows.

\begin{Theorem}
Let $\Delta$ be a $1$-dimensional flag simplicial complex, $I_{\Delta}$ be its Stanley-Reisner ideal and $2 \leq p \leq \nu(\overline{\Delta})$. Then: 

\begin{itemize}
\item[(A)] \textup{[Corollary~\ref{coro: lin related}]}
$I_{\Delta}^{[p]}$ is linearly related if and only if $\Delta$ does not contain an induced subgraph isomorphic to the complete bipartite graph $K_{2+i,\,2p-i}$, for any even integer $i$ with $0 \leq i \leq p$.

\item[(B)] \textup{[Theorem~\ref{Thm: characterization linear resolution}]}
$I_{\Delta}^{[p]}$ has a linear resolution if and only if $\Delta$ contains neither an induced subgraph isomorphic to $K_{2+i,\,2p-i}$ for any even integer $i$ with $0 \leq i \leq p$, nor an induced subgraph isomorphic to the crown graph $Cr(2p+1)$ on $4p+2$ vertices.
\end{itemize}
\end{Theorem}

The theorem above shows that two different kinds of combinatorial configurations prevent $I_{\Delta}^{[p]}$ from having a linear resolution: the complete bipartite graphs in (A) already prevent the ideal from being linearly related, while the crown graph in (B) disrupts the linearity of the resolution at higher syzygy levels.

Moreover, when $I_{\Delta}^{[p]}$ does not admit a linear resolution, we describe the exact shape of its Betti table, whose description is summarized in Discussion \ref{Disc. summarize}.

%%%%%%%%%%%%%%%%%%%%%%%%%%%%%%%%%%%%%%%%%%%%%%%%%%%%%%%%%%%%%%%%%%%%%%%%%%%%%%

The paper is organized as follows. Section~\ref{Sec1: Preliminaries} provides the necessary background on simplicial complexes, edge ideals, and graded free resolutions, along with the main properties of squarefree powers and Hochster's formula.

Section~\ref{sec: 1-dim flag simplicial complexes} focuses on the squarefree powers $I_{\Delta}^{[p]}$ of the Stanley--Reisner ideal of a $1$-dimensional flag simplicial complex $\Delta$, and describes its associated Stanley--Reisner complex, denoted by $\Delta^{[p]}$. We first show in Proposition~\ref{Thm: dimension} that its dimension can only take the values; in particular: $\dim \Delta^{[p]}=2p-2$ if and only if $I_{\Delta}^{[p]}$ is a squarefree Veronese ideal (Proposition~\ref{prop: veronese}), and $\dim \Delta^{[p]}=2p-1$ if and only if $\Delta$ contains an induced subgraph isomorphic to $K_{i,\,2p-i}$, for some odd integer $1 \leq i \leq p$ (Corollary~\ref{coro: facet}). The latter is a consequence of Theorem \ref{Thm: Facet, pure, CM}, where we also completely describe the $(2p-1)$-dimensional facets of $\Delta^{[p]}$, and we characterize when $\Delta^{[p]}$ is pure and Cohen--Macaulay. We finally conclude the section by combinatorial describe the projective dimension of $I_{\Delta}^{[p]}$ (Theorem~\ref{Thm: pd}).
%%%%%%%%%%%%%%%%%%%%%%%%%%%%%%%%%%%%%%%%%%%%%%%%%%%%%%%%%%%%%%%%%%%%%%%%%%%%%%%%

 In Section~\ref{Sec: linearly related}, we investigate when the edge ideal $I(G)^{[p]}$ of an arbitrary graph $G$ is linearly related. We obtain a complete characterization in Theorem~\ref{Thm:linrelated1}, which constitutes one of the principal results of the paper. Its proof relies on two main ingredients. The first is the study of the connectivity of the graph $G_I^{(u,v)}$ associated with a monomial ideal, introduced in~\cite[Corollary~2.2]{BHZ}. The second is a theorem of Gasharov, Peeva, and Welker~\cite{Gasharov1999} concerning the lcm-lattice of a monomial ideal, used in Proposition~\ref{thm:betti} to prove that the graded Betti numbers $\beta_{1,j}(I(G)^{[p]})$ can be non-zero only in the two degrees $j=2p+1$ and $j=2p+2.$ 

It is worth mentioning that the Stanley--Reisner ideal of a $1$-dimensional flag simplicial complex can be regarded as the edge ideal of the complement of a triangle-free graph. Hence, as a consequence of Theorem~\ref{Thm:linrelated1}, we are able to characterize when $I_{\Delta}^{[p]}$ is linearly related in Corollary~\ref{coro: lin related}. The section concludes with the following result, describing the Betti number responsible for the failure of the linearly related property of $I_{\Delta}^{[p]}$. 

\begin{Theorem}[Theorem~\ref{Thm: betti 2,2p+2}]
Let $\Delta$ be a $1$-dimensional flag simplicial complex on $[n]$, and let $2 \leq p \leq \nu(\overline{\Delta})$. Then
$$
\beta_{1,2p+2}(I_{\Delta}^{[p]})
=
\sum_{\substack{0 \leq i \leq p \\ i \text{ even}}}
\Big|
\big\{
\text{induced subgraphs of } \Delta
\text{ isomorphic to } K_{2+i,\,2p-i}
\big\}
\Big|.
$$
\end{Theorem}

The proof is strongly based on methods from Combinatorial Topology, in particular on the description of the homological cycles in top degree together with Hochster's formula.

In Section~\ref{sec: linearity simpl complex}, we go beyond the study of the first syzygy module of $I_{\Delta}^{[p]}$ and investigate the linearity of the entire graded free resolution. Here, we establish our second main result, namely Theorem~\ref{Thm: characterization linear resolution}, in which we completely characterize when $I_{\Delta}^{[p]}$ has a linear resolution. As in the proof of Theorem~\ref{Thm: betti 2,2p+2}, our approach combines arguments from Combinatorial Topology, in particular the study of simplicial homology in top degree, and pure Combinatorics (see Propositions~\ref{Prop: Crown implies no linear} and~\ref{Prop: other implciation}). We conclude the section with two consequences: one shows that if $I_{\Delta}^{[p]}$ has a linear resolution, then so does $I_{\Delta}^{[q]}$ for any $q\geq p$ (see Corollary~\ref{Coro: lin res for p then lin res for p+1}), and the second shows that the Betti number $\beta_{2p+1,\,4p+2}(I_{\Delta}^{[p]})$ counts the induced crown graphs $Cr(2p+1)$ contained in $\Delta$ (see Corollary~\ref{coro: betti for crown}).

Section~\ref{Sec: final} is the final section of this work. Here, we determine the shape of the Betti table of $I_{\Delta}^{[p]}$; in particular, we provide the non-vanishing Betti numbers with internal degree $2p+1$ by using the Mayer--Vietoris sequence (see Proposition~\ref{Prop: non van betti}). In Discussion~\ref{Disc. summarize}, we summarize all the results on the graded minimal free resolution of $I_{\Delta}^{[p]}$. We conclude with a short subsection, where we provide some open questions for future works.

%%%%%%%%%%%%%%%%%%% OLD %%%%%%%%%%%%%%%%%%%%%%%%%%%%%%%%%%%

% The investigation of algebraic invariants of monomial ideals by means of combinatorial and topological techniques constitutes one of the fundamental directions in combinatorial commutative algebra. A central aspect of the subject is the deep interaction between the algebraic and homological properties of a monomial ideal and the combinatorial structure of the objects from which it arises. One approach to studying algebraic properties of square-free monomial ideals is to examine the
% combinatorics of their associated simplicial complexes or hypergraphs. To any square-free monomial
% ideal I one can associate both a Stanley–Reisner complex whose faces are given by the square-free
% monomials not in I and a facet complex whose facets correspond to the minimal generators of I.
% One can also consider the facet complex as a hypergraph whose edge ideal is I. It turns out that
% specific algebraic characteristics often correspond to complexes or hypergraphs with a well-defined
% combinatorial structure. In particular, squarefree monomial ideals are naturally associated with combinatorial objects such as simplicial complexes and graphs, allowing one to study their algebraic invariants through combinatorial and topological methods. 

% Among the various homological properties that have been extensively investigated, the existence of a linear resolution plays a distinguished role. Roughly speaking, an homogeneous ideal admits a linear resolution if its minimal graded free resolution is generated in the smallest possible degrees, so that means that it looks like fancy. 

% One of the foundational results in this direction is due to Fr\"oberg~\cite{Froberg1990}, who proved that the edge ideal of a graph has a linear resolution if and only if the graph is co-chordal. In \cite{EAGON1998265}, Eagon and Reiner showed that the property of having a linear resolution is dual to the Cohen–Macaulay property. 
% Unfortunately, it is too much to expect that the property of having a linear resolution can be
% described purely through a combinatorial property of an associated combinatorial structure as the existence of such a resolution does depend on the field over which the polynomial ring is defined. The triangulation of the real projective plane is a typical example. However the results that we obtain in this paper are surprisinlgy characeristic free. 



% In this paper, we study  squarefree powers of edge ideals of graphs and its linearity.  In the last decade squarefree-powers have been attracted a lot of attention. 
% Given a monomial ideal $I$, its $k$-th \emph{squarefree power}, denoted by $I^{[k]}$, is generated by the squarefree monomials among the generators of~$I^k$.  The study of squarefree powers was formally initiated by \cite{BHZ} in the context of edge ideals of graphs, although such ideals had implicitly appeared earlier (for instance, in the proof of Proposition 2.10 in \cite{JZ}). Since then, numerous works have appeared, such as \cite{CFL, EF, EHHM, EHHM2, EH, SEYEDFAKHARI2024107488, Kamalesh1, KAMBERI2026240, S}. Squarefree powers provide a combinatorial shadow of the ordinary powers of~$I$, and they are often easier to analyze while still encoding important homological information. Indeed, by~\cite[Lemma~4.4]{HerzogHibiZheng}, the multigraded minimal free resolution of $I^{[k]}$ is a subcomplex of that of~$I^k$, implying $\reg I^{[k]} \leq \reg I^k$.  Moreover, if $I^{[k]}$ fails to have a linear resolution, then neither does~$I^k$. 
% Another interesting aspect of the study of squarefree powers arises from their deep connection with the matching theory of simplicial complexes. Let $\Delta$ be a simplicial complex, and let $I(\Delta)$ denote its facet ideal. A \emph{matching} of $\Delta$ is a collection of pairwise disjoint facets of $\Delta$. The minimal generators of $I(\Delta)^{[k]}$ bijectively correspond to the matchings of $\Delta$ of size $k$. In particular, $I(\Delta)^{[k]} \neq 0$ if and only if $1 \leq k \leq \nu(\Delta)$, where $\nu(\Delta)$ denotes the maximum size of a matching of $\Delta$.  

%  Bigdeli, Herzog, and Zaare-Nahandi \cite[Theorem
% 5.1]{BHZ} showed that the highest non-vanishing squarefree power of the edge ideal of a graph has
% linear quotients and so linear resolution. MOreover, in \cite[Conjecture 5.3]{BHZ}, it was conjectured that $I(G)^{[p]}$ is linearly related, i.e., its first syzygy module is generated by linear forms, if and only if $p \geq \nu_0(G)$, where $\nu_0(G)$ denotes the restricted matching number of $G$. Recall that a matching is called \emph{restricted} if one of its edges forms a gap with every other edge in the matching. However, a counterexample to this conjecture was later given in \cite[page 12]{EHHM}. This leaves open the problem of determining, for a given graph $G$, the precise conditions on $G$ under which $I(G)^{[p]}$ is linearly related. One of the main result of this paper is to address this question. We prove in fact the following result.

% \begin{Theorem}[Theorem \ref{Thm:linrelated1}]
% Let $G$ be a graph and $2 \leq p \leq \nu(G)$. Then $I(G)^{[p]}$ is linearly related if and only if there does not exist an induced subgraph $H$ of $G$ on $2p+2$ vertices such that $H$ is disconnected and $\nu(H)=p+1$.
% \end{Theorem}

% This leave pose that next question.

% \textbf{Question.} For a given graph $G$, the precise conditions on $G$ under which $I(G)^{[p]}$ has a linear resolution?

% Answer this question looks very hard. In any case, we treat some case, in particular when $G$ is the complement of a triangle-free graph. This case is particular interisting, because this ideals appears as the Stanley-Reisner ideal of $1$-dimensional flag simplcial complexes. 

% \begin{Theorem}
% Let $\Delta$ be a $1$-dimensional flag simplicial complex on $[n]$, and let $2 \leq p \leq \nu(\overline{\Delta})$. 

% A - Corollary \ref{coro: lin related}) $I_{\Delta}^{[p]}$ is linearly related if and only $\Delta$ does not contain an induced subgraph isomorphic to $K_{2+i,\,2p-i}$ for some even integer $i$ with $0 \leq i \leq p$.

% B - Theorem \ref{Thm: characterization linear resolution} ) The following conditions are equivalent:
% \begin{enumerate}
% \item $I_{\Delta}^{[p]}$ does not have a linear resolution;
% \item one of the following conditions holds:
% \begin{enumerate}
%     \item $\Delta$ contains an induced subgraph isomorphic to $K_{2+i,\,2p-i}$ for some even integer $i$ with $0 \leq i \leq p$;
%     \item $\Delta$ contains an induced subgraph isomorphic to the crown graph $Cr(2p+1)$.
% \end{enumerate}
% \end{enumerate}
% \end{Theorem}

% Condition~(2) in the theorem above shows that there exist two distinct types of obstructions to linearity. The first one, described in~(a), already obstructs the property of being linearly related, as shown in Corollary~\ref{coro: lin related}. The second one, described in~(b), arises at higher syzygy levels and obstructs the existence of a linear resolution.

% %%%%%%%%%%%%%%%%%%%%%%%%%%%%%%%%%%%%%%%%%%%%%%%%%%%%%%%%%%%%%%%%%%%%%%%%%%%%%%

% The paper is organized as follows. Section~\ref{Sec1: Preliminaries} contains the preliminary material and is divided into two subsections for the reader’s convenience. In Subsection~\ref{Subsec1}, we recall the basic definitions and notation from simplicial complex theory and graph theory, together with the associated $K$-algebras, such as Stanley--Reisner rings. In Subsection~\ref{SubSec2}, we review some background from commutative and homological algebra, including graded free resolutions, linearity, and Castelnuovo--Mumford regularity. We also recall Hochster’s formula for computing the graded Betti numbers of squarefree monomial ideals and introduce the notion of squarefree powers of monomial ideals, following \cite{BHZ,EHHM}.

% In Section~\ref{sec: 1-dim flag simplicial complexes}, we initiate the study of the $p$-th squarefree powers $I_{\Delta}^{[p]}$ of the Stanley--Reisner ideal associated with a $1$-dimensional flag simplicial complex $\Delta$ on $[n]$, and we describe the corresponding Stanley--Reisner complex $\Delta^{[p]}$. The first problem we address concerns the dimension of $K[\Delta^{[p]}]$. We prove that it can attain only the two values $2p-1$ and $2p$ (see Proposition~\ref{Thm: dimension} (1)). Moreover, in Proposition~\ref{prop: veronese} we show that $\dim K[\Delta^{[p]}]=2p-1$ if and only if $I_{\Delta}^{[p]}$ is the squarefree Veronese ideal generated in degree $2p$ in $n$ variables. Combining Theorem~\ref{Thm: Facet, pure, CM}~(1), where the facets of dimension $2p-1$ are explicitly described, and Proposition Proposition~\ref{Thm: dimension} (2), we get Corollary~\ref{coro: facet}: $\dim K[\Delta^{[p]}]=2p$ if and only if $\Delta$ contains an induced subgraph isomorphic to $K_{i,\,2p-i}$ for some odd integer $1\leq i \leq p$. On the other hand, in Theorem~\ref{Thm: Facet, pure, CM}~(1), we also characterize when $\Delta^{[p]}$ is pure and when it is Cohen--Macaulay.


% %%%%%%%%%%%%%%%%%%%%%%%%%%%%%%%%%%%%%%%%%%%%%%%%%%%%%%%%%%%%%%%%%%%%%%%%%%%%%%%%


% In Section \ref{Sec: linearly related}, we investigate when the first syzygy module of $I_{\Delta}^{[p]}$ is generated by linear forms, namely, when $I_{\Delta}^{[p]}$ is linearly related, and we give a complete characterization in Corollary \ref{coro: lin related}.

% % \begin{Corollary}[Corollary \ref{coro: lin related}]
% % Let $\Delta$ be a $1$-dimensional flag simplicial complex on $[n]$, and let
% % $2 \leq p \leq \nu(\overline{\Delta})$. The following conditions are equivalent:
% % \begin{enumerate}
% %     \item $I_{\Delta}^{[p]}$ is linearly related;
% %     \item $\Delta$ does not contain an induced subgraph isomorphic to $K_{2+i,\,2p-i}$ for some even integer $i$ with $0 \leq i \leq p$.
% % \end{enumerate}
% % \end{Corollary}

% Actually, this result is a consequence of a stronger theorem, Theorem \ref{Thm:linrelated12}, that is one of the first main result of our paper and provides a complete combinatorial characterization of the linearity property for squarefree powers of edge ideals of graphs.



% %and it completely settles the problem arising from \cite[Conjecture 5.3]{BHZ}.% and the counterexample exhibited in \cite[page 12]{EHHM}.
% The proof relies on two main ingredients. The first is the study of the connectivity of the graph $G_I^{(u,v)}$ associated to a monomial ideal, introduced in \cite[Corollary 2.2]{BHZ}. The second is a result of Gasharov, Peeva, and Welker \cite{Gasharov1999} on the lcm-lattice of a monomial ideal, which is used in Proposition \ref{thm:betti} to prove that the graded Betti numbers $\beta_{1,j}(I(G)^{[p]})$ can be non-zero only in degrees $j=2p+1$ and $j=2p+2$.

% The section concludes with a combinatorial interpretation of the Betti number $\beta_{1,2p+2}(I_{\Delta}^{[p]})$, which provides the failure of the linearity of the first syzygy module in the graded free resolution of $I_{\Delta}^{[p]}$, and whose proof is strongly based on the study of simplicial homology in top degree together with Hochster's formula.

% \begin{Theorem}[Theorem \ref{Thm: betti 2,2p+2}]
% Let $\Delta$ be a $1$-dimensional flag simplicial complex on $[n]$, and let
% $2 \leq p \leq \nu(\overline{\Delta})$. Then
% $$
% \beta_{1,2p+2}(I_{\Delta}^{[p]})
% =
% \sum_{\substack{0 \leq i \leq p \\ i \text{ even}}}
% \big|
% \{\text{induced subgraphs of } \Delta \text{ isomorphic to } K_{2+i,\,2p-i}\}
% \big|.
% $$
% \end{Theorem}

% For $1$-dimensional flag simplicial complexes, one can go beyond the study of the linearity of the first syzygy module of squarefree powers and characterize the linearity of the entire graded free resolution. This is achieved in Section \ref{sec: linearity simpl complex}, where we establish the second main result of the paper, Theorem \ref{Thm: characterization linear resolution}.

% % \begin{Theorem}\label{Thm: characterization linear resolution}
% % Let $\Delta$ be a $1$-dimensional flag simplicial complex on $[n]$, and let $2 \leq p \leq \nu(\overline{\Delta})$. The following conditions are equivalent:
% % \begin{enumerate}
% % \item $I_{\Delta}^{[p]}$ does not have a linear resolution;
% % \item one of the following conditions holds:
% % \begin{enumerate}
% %     \item $\Delta$ contains an induced subgraph isomorphic to $K_{2+i,\,2p-i}$ for some even integer $i$ with $0 \leq i \leq p$;
% %     \item $\Delta$ contains an induced subgraph isomorphic to the crown graph $Cr(2p+1)$.
% % \end{enumerate}4
% % \end{enumerate}
% % \end{Theorem}



% As in the proof of Theorem~\ref{Thm: betti 2,2p+2}, our approach is based on methods from Combinatorial Topology and on the study of simplicial homology in top degree; see Propositions~\ref{Prop: Crown implies no linear} and~\ref{Prop: other implciation}.

% We conclude the paper by giving a combinatorial interpretation of the Betti number $\beta_{2p+1,4p+2}(I_{\Delta}^{[p]})$ in terms of the number of induced crown graphs $Cr(2p+1)$ contained in $\Delta$; see Corollary~\ref{coro: betti for crown}.

%%%%%%%%%%%%%%%%%%%%%%%%%%%%%%%%%%%%%%%%%%%%%%%%%%%%%%%%%%%%%%%%%%%%%%%%%%%%%%%%%%%%%%%%%%%%%%
%%%%%%%%%%%%%%%%%%%%%%%%%%%%%%%%%%%%%%%%%%%%%%%%%%%%%%%%%%%%%%%%%%%%%%%%%%%%%%%%%%%%%%%%%%%%%%

\section{Preliminaries}\label{Sec1: Preliminaries}
This section is devoted to preliminaries, where we introduce the definitions, notation, and classical results from the literature that will be used throughout the paper.

\subsection{Combinatorics and Stanley--Reisner Rings}\label{Subsec1} A \textit{simplicial complex} $\Delta$ on the vertex set $[n] := \{1,\dots,n\}$ is a nonempty family of subsets of $[n]$ with the property that, whenever $F' \in \Delta$ and $F \subseteq F'$, one has $F \in \Delta$. The elements of $\Delta$ are called \textit{faces}. For a face $F \in \Delta$, its \textit{dimension}, denoted by $\dim(F)$, is defined as $|F| - 1$. A face of dimension $0$ is called a \textit{vertex}, while a face of dimension $1$ is called an \textit{edge}. The \textit{dimension} of $\Delta$ is $\max\{\dim(F) : F \in \Delta\}$.

The maximal faces of $\Delta$ with respect to inclusion are called \textit{facets}, and the set of all facets is denoted by $\mathcal{F}(\Delta)$. For a subset $W \subseteq [n]$, the \textit{induced subcomplex} of $\Delta$ on $W$, denoted by $\Delta_W$, is the simplicial complex consisting of all faces of $\Delta$ that are contained in $W$. We denote by $\overline{\Delta}$ the \textit{complement} of $\Delta$, namely $\overline{\Delta} = \{\, [n] \setminus F : F \in \Delta \,\}$.

Given a collection $\mathcal{F} = \{F_1, \dots, F_m\}$ of subsets of $[n]$, we write $\langle F_1, \dots, F_m \rangle$, or simply $\langle \mathcal{F} \rangle$, for the simplicial complex whose faces are all subsets of $[n]$ contained in at least one $F_i$, for $i = 1,\dots,m$. If $\Delta = \langle F\rangle$, that is, if $\Delta$ is generated by a single facet, then $\Delta$ is called a \textit{simplex}. A simplicial complex $\Delta$ is said to be \textit{pure} if all its facets have the same dimension; in this case, the dimension of $\Delta$ coincides with the common dimension of its facets.

For each $0 \le i \le d-1$, the $i$-th skeleton of $\Delta$ is the simplicial complex $\Delta^{(i)}$ on $[n]$ whose faces are those faces $F$ of $\Delta$ with $|F| \le i+1$. We say that a simplicial complex $\Delta$ is \emph{connected} if, for any two facets $F$ and $G$, there exists a sequence of facets $F = F_0, F_1, \dots, F_q = G$ such that $F_i \cap F_{i+1} \neq \emptyset$ for all $i=0,\dots,q-1$. Observe that $\Delta$ is connected if and only if its $1$-skeleton $\Delta^{(1)}$ is connected.

Recall that a simplicial complex of dimension $1$ is called a \textit{graph}. In this case, the definitions above can be naturally interpreted in the language of graph theory. For instance, if $\Delta$ is a graph on $[n]$ and $W \subseteq [n]$, then the induced subcomplex $\Delta_W$ is referred to as the \textit{induced subgraph} of $\Delta$ on $W$.

Let $\Delta$ be a graph on $[n]$. A \textit{path} of length $r-1$ in $\Delta$ is a sequence of $r$ pairwise distinct vertices $v_1, \ldots, v_r \in [n]$ such that $\{v_i, v_{i+1}\} \in \Delta$ for every $i = 1, \ldots, r-1$. A \textit{cycle} of length $r$ in $\Delta$ (also called an \textit{$r$-cycle}) is a sequence of $r$ distinct vertices $v_1, \ldots, v_r \in [n]$, with $r \ge 3$, such that $\{v_i, v_{i+1}\} \in \Delta$ for every $i = 1, \ldots, r-1$, and $\{v_r, v_1\} \in \Delta$. A cycle of length three is called a \textit{triangle}.

A vertex $v$ of $\Delta$ is called \textit{isolated} if $\{v,w\} \notin \Delta$ for every vertex $w \neq v$. An \textit{independent set} of $\Delta$ is a set of vertices $V \subseteq [n]$ such that $\{v,w\} \notin \Delta$ for all $v,w \in V$ with $v \neq w$.

We now recall two classes of graphs that will play a central role in Sections~\ref{sec: 1-dim flag simplicial complexes} and \ref{sec: linearity simpl complex}.

A \textit{bipartite} graph is a graph whose vertex set can be partitioned into two disjoint independent subsets. Equivalently, a graph is bipartite if and only if it contains no cycles of odd length (see \cite[Lemma 9.1.1]{HHmonomialideals}). A bipartite graph is called \textit{complete} if every vertex in one part is adjacent to every vertex in the other. A complete bipartite graph with parts of sizes $m$ and $n$ is denoted by $K_{m,n}$ (see Figure~\ref{fig:bipartite + crown} (A)).

A \textit{crown graph} on $2n$ vertices, denoted by $Cr(n)$ (see Figure~\ref{fig:bipartite + crown} (B)), is a bipartite graph with vertex set $V(Cr(n))=\{u_1, \dots, u_n\}\sqcup\{v_1, \dots, v_n\}$ and edge set $E(Cr(n))=\{\{u_i,v_j\} : i,j\in [n],\ i\neq j\}.$
 %Equivalently, the crown graph is obtained from a complete bipartite graph by removing the edges of a perfect matching.

\begin{figure}[h]
    \centering
    \subfloat[]{\includegraphics[scale=0.5]{bip_graph.eps}}\qquad
    \subfloat[]{\includegraphics[scale=0.5]{exa_crown_graph.eps}}
    \caption{A complete bipartite $K_{2,4}$ and a crown graph $Cr(5)$.}
    \label{fig:bipartite + crown}
\end{figure}


Let $\Delta$ be a simplicial complex on $[n]$, and let $R = K[x_1, \dots, x_n]$ be the polynomial ring in $n$ variables over a field $K$, where $\deg(x_i)=1$ for all $i \in [n]$. Throughout this paper, $K$ denotes an arbitrary field unless otherwise specified.

A subset $F \subseteq [n]$ is called a \textit{non-face} of $\Delta$ if $F \notin \Delta$. We denote by $\mathcal{N}(\Delta)$ the collection of minimal non-faces of $\Delta$. For any subset $F = \{i_1, \dots, i_r\} \subseteq [n]$, we associate the squarefree monomial $x_F = x_{i_1} \cdots x_{i_r} \in R$.

The \textit{facet ideal} of $\Delta$ is the ideal $I(\Delta) \subset R$ generated by all squarefree monomials $x_F$ such that $F$ is a facet of $\Delta$. If $\Delta$ is a graph, then its facet ideal is called the \textit{edge ideal} of $\Delta$.

The \textit{Stanley--Reisner ideal} of $\Delta$ is the ideal $I_\Delta \subset R$ generated by the squarefree monomials $x_F$ with $F \notin \Delta$, that is,
\[
I_\Delta = (x_F : F \in \mathcal{N}(\Delta)).
\]
The quotient ring $R/I_\Delta$ is called the \textit{Stanley--Reisner ring} of $\Delta$ and is denoted by $K[\Delta]$.

By \cite[Corollary 6.3.5]{Villarreal}, the Krull dimension of $K[\Delta]$, denoted by $\dim K[\Delta]$, is equal to $\dim \Delta + 1$. Equivalently,
\[
\dim K[\Delta] = \max\{s : x_{i_1}\cdots x_{i_s} \notin I_\Delta,\ i_1 < \cdots < i_s\}.
\]

We say that $\Delta$ is \textit{Cohen--Macaulay} over $K$ if $K[\Delta]$ is Cohen--Macaulay. A pure simplicial complex of dimension $d-1$ is said to be \textit{connected in codimension one} if, for any two facets $F$ and $G$ of $\Delta$, there exists a sequence of facets
\[
F = F_0, F_1, \ldots, F_q = G
\]
such that $|F_i \cap F_{i+1}| = d-1$ for all $i = 0, \ldots, q-1$. It is known (see \cite[Lemma 9.1.12]{HHmonomialideals}) that every Cohen--Macaulay simplicial complex is connected in codimension one.

A simplicial complex $\Delta$ is said to be \textit{flag} if its Stanley--Reisner ideal $I_{\Delta}$ is generated in degree two. In particular, if $\Delta$ is a $1$-dimensional flag simplicial complex, then $\Delta$ is a triangle-free graph, and $I_{\Delta} = I(\overline{\Delta})$, that is, the Stanley--Reisner ideal of $\Delta$ coincides with the edge ideal of the complement graph of $\Delta$.

\subsection{Basics in Commutative Algebra and squarefree powers.}\label{SubSec2}

%\subsection{Preliminaries on Commutative Algebra, Simplicial Homology, and Squarefree Powers}

%Throughout the paper, let $K$ be an arbitrary field, and let $R = K[x_1, \ldots, x_n]$ denote the polynomial ring over $K$, where $\deg(x_i)=1$ for all $i \in [n]$.
Let $I$ be a homogeneous ideal of $R=K[x_1, \ldots, x_n]$ generated in degree $t$. By the Hilbert Syzygy Theorem, $I$ admits a minimal graded free resolution $\mathbb{F}(I)$, which is unique up to isomorphism and has finite length at most $n$. Explicitly, $\mathbb{F}(I)$ has the form
\[
0 \to \bigoplus_{j \in \mathbb{Z}} R(-j)^{\beta_{\ell,j}}
\xrightarrow{d_\ell} \cdots
\to \bigoplus_{j \in \mathbb{Z}} R(-j)^{\beta_{i,j}}
\xrightarrow{d_i} \cdots
\to \bigoplus_{j \in \mathbb{Z}} R(-j)^{\beta_{0,j}}
\xrightarrow{d_0} I \to 0,
\]
where $\ell \le n$. The integers $\beta_{i,j}$ are the \emph{graded Betti numbers} of $I$, and $\ell$ is the \textit{projective dimension} of $I$, denoted by $\mathrm{pd}(I)$.

The \emph{Castelnuovo--Mumford regularity}, or simply the \emph{regularity}, of $I$ is defined by
\[
\reg(I) = \max\{\, j - i \mid \beta_{i,j}(I) \neq 0 \text{ for some } i \,\}.
\]
Moreover, one has $\reg(R/I)=\reg(I)-1$ and $\operatorname{pd}(R/I)=\operatorname{pd}(I)+1$.

If $\beta_{i,j}(I) = 0$ for all $i \ge 0$ and all $j \neq i + t$, then $I$ is said to have a \textit{linear resolution}. We say that $I$ is \textit{linearly related} if $\beta_{1,j}(I) = 0$ for all $j \neq t + 1$.

The graded Betti numbers of a squarefree monomial ideal can be computed via Hochster’s formula. We recall that there is a one-to-one correspondence between squarefree monomial ideals and Stanley--Reisner ideals of simplicial complexes. Let $\Delta$ be a simplicial complex with Stanley--Reisner ideal $I_\Delta$. Then
\[
\beta_{i,j}(I_{\Delta}) =
\sum_{\substack{W \subseteq [n] \\ |W| = j}}
\dim_K \widetilde{H}_{j-i-2}(\Delta_W; K),
\]
where $\widetilde{H}_i(\Delta;K)$ denotes the $i$-th reduced simplicial homology group of $\Delta$ with coefficients in $K$. As a consequence of Hochster’s formula, one has $\reg K[\Delta] \le \dim \Delta + 1 = \dim K[\Delta]$.

For completeness, we briefly recall the construction of simplicial homology. Let $\Delta$ be a $d$-dimensional simplicial complex with vertex set $V(\Delta)$, and fix a total order on $V(\Delta)$. For each $i \in \{0,\dots,d\}$, let $C_i(\Delta)$ be the $K$-vector space generated by the oriented $i$-faces $[v_1,\dots,v_{i+1}]$ with $v_1 < \cdots < v_{i+1}$. The augmented chain complex of $\Delta$ is
\[
0 \to C_d(\Delta) \xrightarrow{\delta_d} C_{d-1}(\Delta) \xrightarrow{\delta_{d-1}} \cdots \xrightarrow{\delta_1} C_0(\Delta) \xrightarrow{\delta_0} K \to 0,
\]
where $\delta_0(v)=1$ for all $v \in V(\Delta)$ and, for $1 \le i \le d$, the boundary map $\delta_i : C_i(\Delta) \to C_{i-1}(\Delta)$ is given by
\[
\delta_i \big([v_1,\ldots,v_{i+1}]\big)
=
\sum_{j=1}^{i+1} (-1)^{j+1}
[v_1, \ldots, \hat{v_j}, \ldots, v_{i+1}],
\]
where $\hat{v_j}$ indicates that $v_j$ is omitted. The $i$-th reduced simplicial homology group of $\Delta$ over $K$ is defined as $\widetilde{H}_i(\Delta;K)=\ker \delta_i / \mathrm{Im}\, \delta_{i+1}$. By convention, $\widetilde{H}_{-1}(\emptyset;K)=K$ and $\widetilde{H}_i(\emptyset;K)=0$ for all $i \ge 0$. Moreover, if $\Delta \neq \emptyset$, then $\dim_K \widetilde{H}_0(\Delta;K)$ is one less than the number of connected components of $\Delta$ (see \cite[Chapter 1, Section 7]{MJ}).

We now introduce the notion of \textit{squarefree powers} of a squarefree monomial ideal, following \cite{BHZ}. Let $I$ be a squarefree monomial ideal in $R = K[x_1, \ldots, x_n]$. The $k$-th squarefree power of $I$ is defined as the ideal generated by the squarefree monomials among the minimal generators of $I^k$, and it is denoted by $I^{[k]}$. For instance, if $I=(x,y,z)$ in $K[x,y,z]$, then $I^{[2]}=(xy,yz,xz)$, $I^{[3]}=(xyz)$ and $I^{[k]}=(0)$ for all $k\geq 3$. 

Viewing $I$ as the facet ideal of a simplicial complex $\Delta$, the minimal generators of $I^{[k]}$ are in bijection with the $k$-matchings of $\Delta$.
A \emph{matching} of a simplicial complex $\Delta$ is a collection of facets that are pairwise disjoint. A matching consisting of $k$ facets is called a \textit{$k$-matching}. Such a $k$-matching is said to be \emph{maximal} if $\Delta$ admits no $(k+1)$-matching. The \emph{matching number} of $\Delta$, denoted by $\nu(\Delta)$, is the cardinality of a maximal matching of $\Delta$. Therefore,
\[
I^{[k]} = (x_{i_1}\cdots x_{i_k} : \{i_1, \ldots, i_k\} \text{ is a $k$-matching of } \Delta).
\]
In particular, $I^{[k]} \neq 0$ if and only if $1 \le k \le \nu(\Delta)$, and $I^{[\nu(\Delta)]}$ is the largest nonzero squarefree power of $I$. Moreover, if $\Delta$ is a $1$-dimensional simplicial complex, then $I_{\Delta}^{[p]} \neq (0)$ if and only if $1 \le p \le \nu(\overline{\Delta})$.

%%%%%%%%%%%%%%%%%%%%%%%%%%%%%%%%%%%%%%%%%%%%%%%%%%%%%%%%%%%%%%%%%%%%%%%%%%%%%%%%%%%


\section{Stanley-Reisner ideal of $1$-dimensional flag simplicial complex and its squarefree powers}\label{sec: 1-dim flag simplicial complexes}

In this section, we focus on the squarefree powers of the Stanley--Reisner ideal of a $1$-dimensional flag simplicial complex. Recall that if $\Delta$ is a $1$-dimensional flag simplicial complex, then $\Delta$ is a triangle-free graph and the Stanley--Reisner ideal of $\Delta$ can be identified with the edge ideal of the complement graph of $\Delta$, denoted by $\overline{\Delta}$.

Throughout the paper, we denote by $\Delta^{[p]}$ the Stanley--Reisner complex associated with the squarefree power $I_{\Delta}^{[p]}$. The first question to be addressed is the determination of its Krull dimension. We start with the following example, that can help the reader to preliminary understand its behavior. 


\begin{Example}\rm

\noindent
(1) Let $\Delta_1=\langle \{1,2\},\{1,3\},\{1,4\},\{4,5\}\rangle$ be the $1$-dimensional flag simplicial complex shown in Figure~\ref{fig:exa simpl complex} (A).

\begin{figure}[h]
    \centering
    \subfloat[]{\includegraphics[scale=0.5]{one_dim_complex1.eps}}\qquad\qquad
    \subfloat[]{\includegraphics[scale=0.5]{one_dim_complex2.eps}}
    \caption{A $1$-dimensional simplicial complex and its complement graph.}
    \label{fig:exa simpl complex}
\end{figure}

The Stanley--Reisner ideal of $\Delta_1$ is $I_{\Delta_1} = (x_1x_5,\, x_2x_3,\, x_3x_4,\, x_2x_5,\, x_3x_5),$ which can be identified with the edge ideal of the complement graph $\overline{\Delta_1}$ depicted in Figure~\ref{fig:exa simpl complex} (B). The second squarefree power of $I_{\Delta_1}$ is
\[
I_{\Delta_1}^{[2]} = (x_{1}x_{2}x_{3}x_{5},\, x_{1}x_{2}x_{4}x_{5},\, x_{1}x_{3}x_{4}x_{5},\, x_{2}x_{3}x_{4}x_{5}).
\]

The facets of the simplicial complex $\Delta_1^{[2]}$ are
\[
\{
\{3,4,5\}, \{2,4,5\}, \{1,4,5\}, \{2,3,5\},
\{1,3,5\}, \{1,2,5\}, \{1,2,3,4\}
\}.
\]
Note that $\dim S/I_{\Delta_1}^{[2]} = 4 = 2p.$ Moreover, the unique $3$-dimensional facet is $F=\{1,2,3,4\}$, and the induced subcomplex $(\Delta_1)_F$ is isomorphic to the complete bipartite graph $K_{1,3}$.

\medskip

\noindent
(2) Consider
$\Delta_2=\langle \{1,2\},\{1,3\},\{1,4\},\{5\}\rangle$.
Then
$$
I_{\Delta_2}^{[2]}
=
\left(
x_{1}x_{2}x_{3}x_{5},\,
x_{1}x_{2}x_{4}x_{5},\,
x_{1}x_{3}x_{4}x_{5},\,
x_{2}x_{3}x_{4}x_{5}
\right),
$$
and the facets of $\Delta_2^{[2]}$ are
$$
\{
\{3,4,5\},\,
\{2,4,5\},\,
\{1,4,5\},\,
\{2,3,5\},\,
\{1,3,5\},\,
\{1,2,5\},\,
\{1,2,3,4\}
\}.
$$
As in the previous example, $\dim S/I_{\Delta_2}^{[2]} = 4 = 2p,$
and the unique $3$-dimensional facet is $F=\{1,2,3,4\}$, with $(\Delta_2)_F$ isomorphic to $K_{1,3}$.

It is worth noting that, if we instead consider
$\Delta'_2=\langle \{1,2\},\{1,3\},\{1,4\}\rangle$,
then $I_{\Delta'_2}^{[2]}=(0),$ since $\nu(\overline{\Delta'_2})=1$.

\medskip

\noindent
(3) Let
$\Delta_3=\langle\{1,2\},\{3,4\},\{1,4\},\{5\}\rangle$.
Then
$$
I_{\Delta_3}^{[2]}
=
\left(
x_{1}x_{2}x_{3}x_{4},\,
x_{1}x_{2}x_{3}x_{5},\,
x_{1}x_{2}x_{4}x_{5},\,
x_{1}x_{3}x_{4}x_{5},\,
x_{2}x_{3}x_{4}x_{5}
\right),
$$
and the facets of $\Delta_3^{[2]}$ are
$$
\left\{
\{3,4,5\},
\{2,4,5\},
\{1,4,5\},
\{2,3,5\},
\{1,3,5\},
\{1,2,5\},
\{2,3,4\},
\{1,3,4\},
\{1,2,4\},
\{1,2,3\}
\right\}.
$$
Note that $\dim S/I_{\Delta_3}^{[2]} = 3 = 2p-1.$ Moreover, in this case, $I_{\Delta_3}^{[2]}$ coincides with the ideal generated by all squarefree monomials of degree $4$ in $K[x_1,\ldots,x_5]$, known as squarefree Veronese ideal generated in degree $4$ in $5$ variables.
\end{Example}

The phenomenons observed in the previous examples are not accidental. Indeed, we shall prove that the dimension can attain only the two values $2p-1$ and $2p$. More precisely, the dimension is equal to $2p$ whenever $\Delta$ contains suitable complete bipartite graphs as induced subgraphs. On the other hand, the dimension is equal to $2p-1$ precisely when $I_{\Delta}^{[p]}$ can be identified with the squarefree Veronese ideal generated in degree $2p$ in $n$ variables, which we formally introduce later.

\begin{Proposition}\label{Thm: dimension}
Let $\Delta$ be a $1$-dimensional flag simplicial complex on $[n]$ and $2\leq p\leq \nu(\overline{\Delta})$. Then:
\begin{enumerate}
\item $\dim S/I_{\Delta}^{[p]}\in \{2p-1,2p\}$.
\item If $\Delta$ contains an induced subgraph isomorphic to $K_{i,\,2p-i}$ for some odd integer $i$ with $1 \le i \le p$, then $\dim S/I_{\Delta}^{[p]} = 2p$.
\end{enumerate}

\begin{table}[h]
\centering
\begin{tabular}{c|l}
$p$ & Induced subgraphs \\ \hline
2 & $K_{1,3}$ \\
3 & $K_{1,5},\; K_{3,3}$ \\
4 & $K_{1,7},\; K_{3,5}$ \\
5 & $K_{1,9},\; K_{3,7},\; K_{5,5}$ \\
6 & $K_{1,11},\; K_{3,9},\; K_{5,7}$ \\
7 & $K_{1,13},\; K_{3,11},\; K_{5,9},\; K_{7,7}$
\end{tabular}
\caption{$K_{i,2p-i}$ with $2\leq p\leq 7$, $i$ odd and $1 \le i \le p$.}
\end{table}
\end{Proposition}

\begin{proof}

(1) We begin by showing that $\dim S/I_{\Delta}^{[p]} \leq 2p$. If $n=2p$, there is nothing to prove, since
$\dim S/I_{\Delta}^{[p]} \leq n = 2p$. We may therefore assume that $n>2p$. Since
$\dim S/I_{\Delta}^{[p]}=\max\{s: x_{i_1}\cdots x_{i_s}\notin I_{\Delta}^{[p]},\ i_1<\dots<i_s\}$,
it suffices to prove that
$x_{i_1}\dots x_{i_{2p}}x_{i_{2p+1}}\in I_{\Delta}^{[p]}$
for every $i_1<\dots<i_{2p}<i_{2p+1}$.\\
Let $I_1=\{i_1,\dots,i_{2p},i_{2p+1}\in [n]: i_1<\dots<i_{2p}<i_{2p+1}\}$.
We show that the induced graph of $\overline{\Delta}$ on $I_1$ admits a $p$-matching.\\
Suppose first that $\{i_{1},i_k\}\in \Delta$ for every $k=2,\dots,2p+1$. This implies that
$\{i_h,i_k\}\notin \Delta$ for every $h,k\in \{2,\dots,2p+1\}$ with $h\neq k$, since otherwise $\Delta$
would contain a triangle. Then
$$\{\{i_j,i_{j+1}\}: j\in\{2,\dots,2p\}\ \text{even}\}$$
forms a $p$-matching of $\overline{\Delta}$, as desired. Otherwise, there exists a vertex in $\{i_2,\dots,i_{2p+1}\}$ that is not adjacent to $i_1$ in $\Delta$.
Without loss of generality, we may assume that $\{i_{1},i_{2}\}\notin \Delta$, and we set
$I_2=I_1\setminus\{i_1,i_2\}$. If $\{i_{3},i_k\}\in \Delta$ for every $k=4,\dots,2p+1$, then
$\{i_h,i_k\}\notin \Delta$ for every $h,k\in \{4,\dots,2p+1\}$ with $h\neq k$, since otherwise $\Delta$
would contain a triangle. Therefore,
$$M_1=\{\{i_j,i_{j+1}\}: j\in\{4,\dots,2p\}\ \text{even}\}$$
forms a $(p-1)$-matching of $\overline{\Delta}$ and, as a consequence,
$M_1\cup \{\{i_1,i_2\}\}$ is a $p$-matching of $\overline{\Delta}$. Otherwise, if there exists a vertex, say $i_4$, such that $\{i_{3},i_{4}\}\notin \Delta$, we proceed as above.
Iterating this process, we obtain the desired conclusion, namely that the induced graph of $\overline{\Delta}$
on $I_1$ admits a $p$-matching. Hence
$x_{i_1}\cdots x_{i_{2p}}x_{i_{2p+1}}\in I_{\Delta}^{[p]}$, which implies
$\dim S/I_{\Delta}^{[p]}\leq 2p$, as claimed.\\
It is straightforward to observe that, if we have a set of $2p-1$ vertices of $\Delta$, namely $i_1,\dots,i_{2p-1}$, then we cannot form a $p$-matching in $\overline{\Delta}$, so $x_{i_1}\cdots x_{i_{2p-1}}\notin I_{\Delta}^{[p]}$, which implies
$\dim S/I_{\Delta}^{[p]}\geq 2p-1$. In conclusion, we have $\dim S/I_{\Delta}^{[p]}\in \{2p-1,2p\}$. 

(2) Suppose that $\Delta$ contains an induced subgraph $G$ isomorphic to $K_{i, 2p-i}$ for some odd integer $i$ with $1 \le i \le p$. Let $V(G) = V \cup W$, where $V = \{v_1, \dots, v_i\}$ and $W = \{w_1, \dots, w_{2p-i}\}$, and the edge set is defined as $E(G) = \{ \{v_h, w_k\} : h \in [i], k \in [2p-i] \}$. Then the monomial $x_{v_1} \cdots x_{v_i} x_{w_1} \cdots x_{w_{2p-i}}$ does not belong to $I_{\Delta}^{[p]}$ because, by construction, the matching number of induced subgraph of $\overline{\Delta}$ on $V \cup W$ is less than $p$. Therefore, we conclude that $\dim S/I_{\Delta}^{[p]} = 2p$, since by (1) we have $\dim S/I_{\Delta}^{[p]} \le 2p$.

%     \cup \{\{v_{k-1},v_{k+1}\}\}
%     \cup \{\{v_j,v_{j+1}\}: j=k+2,\dots,2p+1 \ \text{even}\} $$

%Since $G$ is properly contained in $\Delta$, there exists a vertex $v$ of $\Delta$ such that either $v$ is isolated, or $\{v,v_k\}$ is not an edge of $\Delta$ for some $k\in[i]$, or $\{v,w_h\}$ is not an edge of $\Delta$ for some $h\in[2p-i]$. Indeed, $\{v,v_k\}$ and $\{v,w_h\}$ cannot both be edges of $\Delta$, since otherwise $\Delta$ would contain the triangle $\{v,v_k,w_h\}$.

% Without loss of generality, we may assume that $\{v,v_k\}$ is not an edge of $\Delta$
% for some $k\in[i]$. The argument is analogous in the case where
% $\{v,w_h\}$ is not an edge of $\Delta$ for some $h\in[2p-i]$. We distinguish two cases.

% \begin{itemize}
%     \item If $k$ is odd, consider
%     $$M_1=\{\{v_j,v_{j+1}\}: j=0,\dots,k-1 \ \text{odd}\}
%     \cup \{\{v_j,v_{j+1}\}: j=k+1,\dots,2p+1 \ \text{even}\},$$
%     which is a $(p-1)$-matching of $\overline{\Delta}$. Then
%     $M_1\cup \{\{v,v_k\}\}$ is a $p$-matching of $\overline{\Delta}$.
    
%     \item If $k$ is even, consider
%     $$M_2=\{\{v_j,v_{j+1}\}: j=1,\dots,k-2 \ \text{odd}\}
%     \cup \{\{v_{k-1},v_{k+1}\}\}
%     \cup \{\{v_j,v_{j+1}\}: j=k+2,\dots,2p+1 \ \text{even}\},$$
%     which is a $(p-1)$-matching of $\overline{\Delta}$. Then
%     $M_2\cup \{\{v,v_k\}\}$ is a $p$-matching of $\overline{\Delta}$.
% \end{itemize}

\end{proof}

As a consequence of above result, the Castelnuovo--Mumford regularity of $S/I_{\Delta}^{[p]}$ can assume
only two possible values, namely $2p-1$ and $2p$.


\begin{Corollary}\label{Coro:reg}
Let $\Delta$ be a $1$-dimensional flag simplicial complex on $[n]$ and $2\leq p\leq \nu(\overline{\Delta})$. Then
$\reg S/I_{\Delta}^{[p]} \in \{2p-1,2p\}$.
\end{Corollary}

\begin{proof}
By Proposition~\ref{Thm: dimension}, we have $\reg K[\Delta^{[p]}] \le 2p$. Moreover, since
$I_{\Delta}^{[p]}$ is generated in degree $2p$, it follows that
$2p-1 \le \reg K[\Delta^{[p]}]$. Therefore,
$\reg S/I_{\Delta}^{[p]} \in \{2p-1,2p\}$.
\end{proof}
In the following, we show that $I_{\Delta}^{[p]}$ is a squarefree Veronese ideal generated
in degree $2p$ in $n$ variables in the case where
$\dim S/I_{\Delta}^{[p]}$ attains the lower value $2p-1$. Recall that the squarefree
Veronese ideal of degree $d$ in $n$ variables, denoted by $I_{n,d}$, is generated by all squarefree monomials of degree $d$ in $S=K[x_1,\dots,x_n]$. As shown in
\cite{HHmonomialideals}, the ideal $I_{n,d}$ is the Stanley--Reisner ideal of the
$(d-2)$-skeleton of the $(n-1)$-simplex. Consequently, the quotient ring $S/I_{n,d}$ is Cohen--Macaulay with $\dim(S/I_{n,d})=\depth(S/I_{n,d})=d-1$. Moreover, it is proved that $I_{n,d}$ has linear quotients and, hence, a linear resolution.



\begin{Proposition}\label{prop: veronese}
Let $\Delta$ be a $1$-dimensional flag simplicial complex on $[n]$ and let
$2\leq p \leq \nu(\overline{\Delta})$. Then the following conditions are equivalent:
\begin{enumerate}
    \item $\dim S/I_{\Delta}^{[p]}=2p-1$;
    \item $I_{\Delta}^{[p]}$ is a squarefree Veronese ideal generated in degree $2p$
    in $n$ variables.
\end{enumerate}
\end{Proposition}

\begin{proof}
$(1)\Rightarrow(2)$. By definition, $I_{\Delta}^{[p]}$ is generated by squarefree monomials of degree $2p$. Moreover, 
\[
\dim S/I_{\Delta}^{[p]}
=\max\{s: x_{i_1}\cdots x_{i_s}\notin I_{\Delta}^{[p]},\ i_1<\dots<i_s\}.
\]
Assume by contradiction that the set of generators of $I_{\Delta}^{[p]}$ does not
contain all squarefree monomials of degree $2p$. Then there exists a squarefree
monomial of degree $2p$ not belonging to $I_{\Delta}^{[p]}$, which implies
$\dim S/I_{\Delta}^{[p]}=2p$, contradicting the assumption
$\dim S/I_{\Delta}^{[p]}=2p-1$.

\smallskip
$(2)\Rightarrow(1)$. This implication is immediate.
\end{proof}


We now prove the main result of this section. In the case where the Krull dimension of $S/I_{\Delta}^{[p]}$ is $2p$, we provide a complete description of the $(2p-1)$-dimensional facets of $\Delta^{[p]}$ and establish a characterization of its purity and Cohen--Macaulayness. 

For the notation, we set $\mathcal{F}(\Delta, K_{i,j})=\{ F \in \Delta \mid \Delta_F \cong K_{i,j}\}$. Moreover, we denote by $\mathcal{F}_i(\Delta)$ the set of all $i$-dimensional facets of $\Delta$.


\begin{Theorem}\label{Thm: Facet, pure, CM}
    Let $\Delta$ be a 1-dimensional flag simplicial complex on $[n]$ and $2\leq p\leq \nu(\overline{\Delta})$. Assume that $\dim S/I_{\Delta}^{[p]}=2p$. Then
    \begin{enumerate}
   \item The $(2p-1)$-dimensional facets of $\Delta^{[p]}$ are given by 
   \[
   \mathcal{F}_{2p-1}(\Delta^{[p]})= 
   \begin{cases} 
       \mathcal{F}( \Delta, K_{1,2p-1}) \cup \mathcal{F}( \Delta, K_{3,2p-3}) \cup \ldots \cup \mathcal{F}( \Delta, K_{p,p}) & \text{if } p \text{ is odd } \\
     \mathcal{F}( \Delta, K_{1,2p-1}) \cup \mathcal{F}( \Delta, K_{3,2p-3}) \cup \ldots \cup \mathcal{F}( \Delta, K_{p-1,p+1})   & \text{if } p \text{ is even}
   \end{cases}
\]
   
   \item Then $\Delta^{[p]}$ is pure if and only if one of the following holds
     \begin{enumerate}
       \item[(i)] $\Delta \iso K_{2k+1,2\ell+1}$ and $n=(2k+1)+(2\ell+1)\ge 2p$,
       \item[(ii)] $\Delta \iso K_{2k+1,2\ell}$ and $\ell \geq p$,
        \item[(iii)] $\Delta \iso K_{2k,2\ell}$ and $k,\ell \geq p$.
    \end{enumerate}

\item Then  $\Delta^{[p]}$ is Cohen-Macaulay if and only if $\Delta \iso K_{1,n-1}$ and $n>2p$.
    \end{enumerate}
\end{Theorem}

\begin{proof}

We divide the proof into three parts, corresponding to the statements above.

\textbf{(1)} First, we show by induction on $p$ that $\Delta_F$ is bipartite for any facet $F$ of $\Delta^{[p]}$ of dimension $2p-1$. \\
Let $F$ be a $(2p-1)-$dimensional facet of $\Delta^{[p]}$. %In the case $p=1$, there is nothing to prove and,
In the case $p=2$, $\Delta_F$ is trivially bipartite since $\Delta_F$ does not contain any triangle. Now assume $p \ge 3$. Suppose by contradiction that $\Delta_F$ is not bipartite. Then $\Delta_F$ contains an odd cycle $C_{2s+1}$ as an induced subgraph, for some $s\geq 1$.

\textit{Case 1.} Suppose that $2s+1 = 2p-1$.
Set $F \setminus V(C_{2s+1}) = \{v\}$.
Since $\Delta$ does not contain triangles, there exists $w \in V(C_{2s+1})$ such that $\{v,w\} \notin \Delta$.
One can see that
$$
\left( \prod_{u \in C_{2s+1} \setminus \{w\}} u \right)\in I_{\Delta}^{[p-1]}.
$$
Hence,
$$
vw \left( \prod_{u \in C_{2s+1} \setminus \{w\}} u \right)\in I_{\Delta}^{[p]}.
$$
Therefore, $F$ cannot be a facet of $\Delta^{[p]}$, which is a contradiction.

\textit{Case 2.} Suppose that $2s+1 \le 2p-3$.
Then there exist $v_1, v_2, v_3$ in $F \setminus V(C_{2s+1})$.
Since $\Delta$ does not contain triangles, we may assume that $\{v_1,v_2\} \notin \Delta$.
By induction on the number of vertices, we have
$$
\left(\prod_{u \in F \setminus \{v_1, v_2\}} u\right) \in I_{\Delta}^{[p-1]}.
$$
Hence
$$
v_1 v_2 \left(\prod_{u \in F \setminus \{v_1, v_2\}} u\right) \in I_{\Delta}^{[p]}.
$$
Therefore, $F$ cannot be a facet of $\Delta^{[p]}$, which is a contradiction.

Now we know that $\Delta_F$ is bipartite and is contained, as a spanning subgraph, in a complete bipartite graph $K_{\ell,m}$, where $\ell$ and $m$ are either both even or both odd.
If both $\ell$ and $m$ are even, the complement graph $\overline{K_{\ell,m}}$ has a $p$-matching. Hence so does $\overline{\Delta_F}$. Therefore, $F$ cannot be a facet of $\Delta^{[p]}$.
If both $\ell$ and $m$ are odd and $\Delta_F$ is not the complete bipartite graph $K_{\ell,m}$, then the complement graph $\overline{\Delta_F}$ has a $p$ a $p$-matching. Hence $F$ cannot be a facet of $\Delta^{[p]}$.
Thus, if $F$ is a facet of $\Delta^{[p]}$, then $\Delta_F$ must be the complete bipartite graph $K_{\ell,m}$ with $\ell$ and $m$ odd.

On the other hand, $\overline{K_{\ell,m}}$ with $\ell$ and $m$ odd does not have a $p$-matching. Hence, if $\Delta_F$ is $K_{\ell,m}$ with $\ell$ and $m$ odd, then $F$ is a facet of $\Delta^{[p]}$.

\textbf{(2)}
We assume that $\Delta^{[p]}$ is pure.
We show that $\Delta$ is bipartite.
Suppose that $\Delta$ contains an odd cycle $C_{2s+1}$ with $s \ge 2$.
Then we can take $W \subset V$ with $|W| = 3$ such that $\Delta_W$ is the disjoint union of
an edge and an isolated vertex. However, this cannot be contained as an induced subgraph in a complete bipartite graph $K_{\ell,m}$ with $\ell + m = 2p$ and $\ell, m$ odd. Hence we get a contradiction and, then, $\Delta$ has to be bipartite.

Next, we show that $\Delta$ is complete bipartite.
Suppose that $\Delta$ is contained in a complete bipartite graph $K_{s,t}$ with
$X \sqcup Y$ as vertex set partition, and that $\Delta$ is not complete bipartite.
We may assume that $|X| \ge 2$. Let $\{x, y\} \notin \Delta$, $\{x', y\} \in \Delta$, where $x, x' \in X$, $y \in Y$.
Take subsets $X' \subset X$, $Y' \subset Y$ such that $x, x' \in X'$, $y \in Y'$, and $|X'| + |Y'| = 2p-1$.
Then $\Delta_{X' \cup Y'}$ is neither a complete bipartite graph nor an edgeless graph (that is, a graph in which no vertices are connected by edges). Hence $\Delta_{X' \cup Y'}$ cannot be an induced subgraph of a complete bipartite graph.
Thus $X' \cup Y'$ is not contained in a facet of dimension $2p-1$.
Therefore, $\Delta$ is complete bipartite.

Since $p \leq \nu(\overline{\Delta})$, we have $n \ge 2p$.

If $\Delta \cong K_{2k+1,2\ell}$ with $\ell < p$ and $k > 0$, then 
$K_{2p-2\ell-1,2\ell}$ cannot be an induced subgraph of the form 
$K_{2s+1, 2t+1}$ with $(2s+1) + (2t+1) = 2p$.
Hence $V(K_{2p-2\ell-1,2\ell})$ is a facet of $\Delta^{[p]}$ of dimension $2p-2$. Therefore $\Delta^{[p]}$ is not pure of dimension $2p-1$.
Hence $\ell \ge p$ is necessary.

If $\Delta \cong K_{1,2\ell}$ with $\ell < p$, then $2\ell + 1 \le 2p-1$,
which contradicts the assumption that $p \leq \nu(\overline{\Delta})$.
Hence $\ell \ge p$ is necessary.

If $\Delta \cong K_{2k,2\ell}$ with $k \le \ell$ and $k < p$,
then $K_{2k,2p-2k-1}$ cannot be an induced subgraph of the form 
$K_{2s+1, 2t+1}$ with $(2s+1) + (2t+1) = 2p$.
Hence $V(K_{2k,2\ell})$ is a facet of $\Delta^{[p]}$ of dimension $2p-2$. Therefore $\Delta^{[p]}$ is not pure of dimension $2p-1$.
Hence $k, \ell \ge p$ is necessary.

Next, we show the converse.
For any facet $F$ of $\Delta^{[p]}$ with $|F| = 2p-1$,
$\Delta_F$ is a complete bipartite graph of the form $K_{2s,2t+1}$ with $2s + (2t+1) = 2p-1$, which is an induced subgraph of $K_{2s+1,2t+1}$,
whose vertex set is a facet of $\Delta^{[p]}$. Thus $\Delta^{[p]}$ is pure.

\textbf{(3)}
Assume that $\Delta^{[p]}$ is Cohen--Macaulay.
Then $\Delta^{[p]}$ is pure.
By (2), $\Delta$ is isomorphic to $K_{s,t}$ for some $s,t > 0$ with $s + t = n$.

Assume that $s,t > 1$ and that both $s$ and $t$ are odd.
Take $3 \le \ell \le s$ and $3 \le m \le t$ such that $\ell + m = 2p+2$.
Then there exist $F \in \Delta^{[p]}$ such that $\Delta_F \cong K_{\ell-2,m}$
and $G \in \Delta^{[p]}$ such that $\Delta_G \cong K_{\ell,m-2}$.
Then $F$ and $G$ cannot be connected in codimension one.

Assume that $s \ge 2p$ is even and $t > 1$.
Take $3 \le \ell \le s$ and $3 \le m \le t$ such that $\ell + m = 2p+2$.
Then there exist $F \in \Delta^{[p]}$ such that $\Delta_F \cong K_{2p-1,1}$
and $G \in \Delta^{[p]}$ such that $\Delta_G \cong K_{2p-3,3}$.
Then $F$ and $G$ cannot be connected in codimension one.

If $\Delta^{[p]}$ is Cohen--Macaulay, then $\Delta \cong K_{1,n-1}$.
If $n = 2p$, then $\nu(\overline{\Delta}) = p-1$, which contradicts the assumption that $\nu(\overline{\Delta}) \ge p$. Therefore $n > 2p$.
\end{proof}

\begin{Corollary}\label{coro: facet}
Let $\Delta$ be a $1$-dimensional flag simplicial complex on $[n]$ and $2\leq p\leq \nu(\overline{\Delta})$. Then, $\dim S/I_{\Delta}^{[p]}= 2p$ if and only if $\Delta$ contains an induced subgraph isomorphic to $K_{i,\,2p-i}$ for some odd integer $i$ with $1 \le i \le p$.
\end{Corollary}

\begin{proof}
    It follows from Proposition  \ref{Thm: dimension} and Theorem \ref{Thm: Facet, pure, CM}.
\end{proof}
We conclude the section by giving the values of the projective dimension.
\begin{Theorem}\label{Thm: pd}
Let $\Delta$ be a $1$-dimensional flag simplicial complex on $[n]$, and let $2 \leq p \leq \nu(\overline{\Delta})$. \\
If $\overline{\Delta}$ has a unique $p$-matching then $I_{\Delta}^{[p]}$ is principal and $\operatorname{pd}(S/I_{\Delta}^{[p]})=1$.\\
If $\overline{\Delta}$ has more than one $p$-matching, then
\[
\operatorname{pd}(S/I_{\Delta}^{[p]})=
\begin{cases}
n-2p, &
\begin{array}{l}
\text{if } \Delta \cong K_{1,n-1} \text{ or } \Delta \text{ does not contain an induced subgraph}\\
\text{isomorphic to } K_{i,\,2p-i} \text{ for each odd integer } i \text{ with } 1\le i\le p,
\end{array}
\\[10pt]
n-2p+1, & \text{otherwise}.
\end{cases}
\]
\end{Theorem}

\begin{proof}
    Assume that $\Delta \cong K_{1,n-1}$. Since $2 \leq p \leq \nu(\overline{\Delta})$, we necessarily have $n\geq 2p+1$. By Theorem~\ref{Thm: dimension} (3), the ring $S/I_{\Delta}^{[p]}$ is Cohen--Macaulay. Moreover, $\Delta$ contains an induced subgraph isomorphic to $K_{1,\,2p-1}$, so $\dim S/I_{\Delta}^{[p]}= 2p$. Therefore from the Auslander--Buchsbaum formula, we get $\operatorname{pd}(S/I_{\Delta}^{[p]}) = n-2p$.

If $\Delta$ does not contain an induced subgraph isomorphic to $K_{i,\,2p-i}$ for each odd integer $i$ with $1 \le i \le p$, then by Corollary \ref{coro: facet} and Proposition \ref{Thm: dimension} (1), we have $\operatorname{dim}(S/I_{\Delta}^{[p]}) = 2p-1$. Hence, from Proposition \ref{prop: veronese}  $I_{\Delta}^{[p]}$ is a squarefree Veronese ideal generated in degree $2p$ in $n$ variables, and it is well-known that $\operatorname{pd}(S/I_{\Delta}^{[p]}) = n-2p$.

Now assume that $\Delta$ is not isomorphic to $K_{1,n-1}$ and $\Delta$ contains an induced subgraph isomorphic to $K_{i,\,2p-i}$ for some odd integer $i$ with $1 \le i \le p$. Hence  $\operatorname{dim}(S/I_{\Delta}^{[p]}) = 2p$ and $S/I_{\Delta}^{[p]}$ is not Cohen-Macaulay, that is, $\operatorname{dim}(S/I_{\Delta}^{[p]})>\operatorname{depth}(S/I_{\Delta}^{[p]})$. Therefore, from the Auslander--Buchsbaum formula, it follows that $\operatorname{pd}(S/I_{\Delta}^{[p]}) > n-2p$. On the other hand, since $I_{\Delta}^{[p]}$ is a squarefree monomial ideal generated in degree $2p$ in $n$ variables, Hochster's formula yields $\operatorname{pd}(I_{\Delta}^{[p]}) \le n-2p$, that is $\operatorname{pd}(S/I_{\Delta}^{[p]}) \le n-2p+1$. Hence  $\operatorname{pd}(S/I_{\Delta}^{[p]}) = n-2p+1$.
\end{proof}

%%%%%%%%%%%%%%%%%%%%%%%%%%%%%%%%%%%%%%%%%%%%%%%%%%%%%%%%%%%%%%%%%%%%%%%%%%%%%%%%%%%

\section{Linearly related squarefree powers of edge ideals and of Stanley-Reisner ideal of 1-dimensional flag simplicial complexes}\label{Sec: linearly related}

The main goal of this section is to characterize the graphs $G$ for which the squarefree power $I(G)^{[p]}$ is linearly related, that is, the first syzygy module is generated by linear forms. We show that this property admits a simple graph-theoretic description in terms of induced subgraphs.

%The main goal of this section is to determine, for a given graph $G$, the precise conditions on $G$ under which $I(G)^{[p]}$ is linearly related, that is, the first syzygy module is generated by linear forms.  We prove the following result.


\begin{Theorem}\label{Thm:linrelated1}
Let $G$ be a graph and $2 \leq p \leq \nu(G)$. Then $I(G)^{[p]}$ is linearly related if and only if there does not exist an induced subgraph $H$ of $G$ on $2p+2$ vertices such that $H$ is disconnected and $\nu(H) = p+1$.
\end{Theorem}

%The following example illustrates Theorem \ref{Thm:linrelated1}.
Before turning to the proof of Theorem~\ref{Thm:linrelated1}, we present an example illustrating the theorem and record an immediate consequence for Stanley--Reisner ideals of $1$-dimensional flag simplicial complexes.

\begin{Example}\rm\label{exa1}
Consider the graph $G_1$ in Figure \ref{fig: graphs} (A). Its edge ideal is
\[
I(G_1)=(x_1x_2,x_1x_5,x_2x_5,x_3x_5,x_3x_4,x_4x_5,x_5x_6,x_5x_7,x_6x_7),
\] 
and its second squarefree power is generated by
\[
\begin{aligned}
I(G_1)^{[2]} = (&x_1x_2x_3x_4,\, x_1x_2x_3x_5,\, x_1x_2x_4x_5,\, x_1x_3x_4x_5,\, x_2x_3x_4x_5,\, x_1x_2x_5x_6,\, x_3x_4x_5x_6,\\
&x_1x_2x_5x_7,\,x_3x_4x_5x_7,\,x_1x_2x_6x_7,\,x_3x_4x_6x_7,\,x_1x_5x_6x_7,\,x_2x_5x_6x_7,\,x_3x_5x_6x_7,\,x_4x_5x_6x_7).
\end{aligned}
\]

In according to Theorem \ref{Thm:linrelated1}, this ideal is not linearly related, since the induced subgraph of $G_1$ on the vertex set $\{1,2,3,4,6,7\}$ is disconnected and has matching number $3$. By using \textit{Macaulay2} \cite{M2}, its Betti table over a field of characteristic zero is the following:
\[
\begin{array}{r|cccc}
      & 0 & 1 & 2 & 3 \\
\hline
\mathrm{total:}  & 15 & 26 & 15 & 3\\
0:  & . & . & . & .\\
1:  & . & . & . & .\\
2:  & . & . & . & .\\
3:  & . & . & . & .\\
4:  & 15 & 24 & 12 & 3\\
5:  & . & 2 & 3 & .
\end{array}
\]

On the other hand, the second squarefree power of the edge ideal of $G_2$ is linearly related, since $G_2$ contains no induced subgraph on $6$ vertices that is disconnected and has matching number $3$.

\begin{figure}[h]
    \centering
    \subfloat[$G_1$]{\includegraphics[scale=0.6]{graph2.eps}}\qquad
    \subfloat[$G_2$]{\includegraphics[scale=0.6]{graph1.eps}}
    \caption{Graphs in the Example \ref{exa1}.}
    \label{fig: graphs}
\end{figure}

\end{Example}

If $\Delta$ is a $1$-dimensional flag simplicial complex, since $I_{\Delta} = I(\overline{\Delta})$, we may equivalently regard $I_{\Delta}$ as the edge ideal of a graph. Thus, the property of being linearly related $I_{\Delta}^{[p]}$, stated below, follows as a special case of Theorem \ref{Thm:linrelated1}.


\begin{Corollary}\label{coro: lin related}
Let $\Delta$ be a $1$-dimensional flag simplicial complex on $[n]$, and let
$2 \leq p \leq \nu(\overline{\Delta})$. The following conditions are equivalent:
\begin{enumerate}
    \item $I_{\Delta}^{[p]}$ is linearly related.
    \item $\Delta$ does not contain an induced subgraph isomorphic to $K_{2+i,\,2p-i}$ for some even integer $i$ with $0 \leq i \leq p$.
\end{enumerate}

\begin{table}[h]
\centering
\begin{tabular}{c|l}
$p$ & Induced subgraphs \\ \hline
2 & $K_{2,4}$ \\
3 & $K_{2,6},\; K_{4,4}$ \\
4 & $K_{2,8},\; K_{4,6}$ \\
5 & $K_{2,10},\; K_{4,8},\; K_{6,6}$\\
6 & $K_{2,12}$, \; $K_{4,10}$ \; $K_{6,8}$\\
7 & $K_{2,14}$, \; $K_{4,12}$ \; $K_{6,10}$ \; $K_{8,8}$\\
\end{tabular}
\caption{Complete bipartite graphs $K_{2+i,\,2p-i}$ with $i$ even and $0 \leq i \leq p$.}
\end{table}
\end{Corollary}



%Let $G$ be a graph and $I(G)$ its edge ideal. In \cite[Conjecture 5.3]{BHZ}, it was conjectured that $I(G)^{[p]}$ is linearly related if and only if $p \geq \nu_0(G)$, where $\nu_0(G)$ denotes the restricted matching number of $G$. Recall that a matching is called \emph{restricted} if one of its edges forms a gap with every other edge in the matching. However, a counterexample to this conjecture was later given in \cite[page 12]{EHHM}. 



To prove Theorem \ref{Thm:linrelated1}, we first establish several preliminary results. We begin by introducing some terminology and recalling \cite[Corollary 2.2]{BHZ}.

 Let $I$ be a monomial ideal generated in degree $d$. In \cite{BHZ}, authors associated a graph $G_I$ to $I$ as follows: $V(G_I)=\cG(I)$, and  $\{u,v\} \in E(G_I)$ if and only if $\deg(\mathrm{lcm}(u,v))=d+1$. Moreover, for all $u,v\in \cG(I)$, the induced subgraph of $G_I$ on the vertex set $\{w\in V(G_I)\mid w\text{ divides } \lcm(u,v)\}$ is denoted by $G^{(u,v)}_{I}$. The connectivity of $G^{(u,v)}_{I}$ characterizes when $I$ is linearly related, as shown in the following. 
  
		\begin{Theorem}\cite[Corollary 2.2]{BHZ}\label{ref:BHN}
		   Let $I$ be a monomial ideal generated in degree d. Then $I$ is linearly related if and only if for all $u,v\in \cG(I)$ there is a path in $G^{(u,v)}_{I}$ connecting $u$ and $v$. 
		\end{Theorem}
        
 In the following remark, we introduce the operation of \textit{switching} and some related properties, which will be used repeatedly in the proof of Proposition \ref{prop:connected}. Before, recall that given a monomial $u \in K[x_1, \ldots, x_n]$, we set $\supp(u) = \{\, i \mid x_i \text{ divides } u \,\}.$ Moreover, if $G$ is a graph, for any edge $e = \{a,b\} \in E(G)$, we define ${\bf x}_e = x_a x_b \in I(G)$. 
    
\begin{Remark}\label{rem:switch}
 {\em  
 Let $G$ be a graph and $u = {\bf x}_{e_1} \cdots {\bf x}_{e_p} \in I(G)^{[p]}$ for some $2 \leq p \leq \nu(G)$. 
Let $f \in E(G)$ be such that $f \notin \supp(u)$ and $f$ does not form a gap with $e_1$. 
Write $e_1 = \{a,b\}$ and $f = \{c,d\}$. 
Since $e_1$ and $f$ do not form a gap, we may assume that $\{a,c\} \in E(G)$. 
Consider the monomials $v' = (u/x_b)x_c$ and $v = (v'/x_a)x_d = (u/{\bf x}_{e_1}){\bf x}_{f}$. 
Then $\deg(\lcm(u,v')) = 2p + 1$ and $\deg(\lcm(v',v)) = 2p + 1$. 
We say that $v$ is obtained from $u$ by {\em switching $e_1$ with $f$}. 
Moreover, $u$ and $v$ are connected in $G_I^{(u,v)}$ by a path of length two given by $u, v', v$.
 }
\end{Remark}



\begin{Proposition}\label{prop:connected}
    Let $G$ be a connected graph with $2p+2$ vertices. Then $I(G)^{[p]}$ is linearly related. 
\end{Proposition}

\begin{proof}
Let $I = I(G)$ and let $u, v \in \mathcal{G}(I^{[p]})$ be two distinct monomials. 
By Theorem~\ref{ref:BHN}, it is enough to show that there exists a path in $G^{(u,v)}_{I^{[p]}}$ connecting $u$ and $v$. 
Let $m = \lcm(u,v)$. 
Then $\deg(m)$ is either $2p + 1$ or $2p + 2$. 
If $\deg(m) = 2p + 1$, then $u$ and $v$ are adjacent in $G^{(u,v)}_{I^{[p]}}$, and we are done. 

Assume that $\deg(m) = 2p + 2$. 
Then $\supp(u)$ and $\supp(v)$ differ in exactly two elements, and we have $\supp(u) \cup \supp(v) = V(G)$. 
Write $u = {\bf x}_{e_1} \cdots {\bf x}_{e_p}$ and $v = {\bf x}_{f_1} \cdots {\bf x}_{f_p}$, where $e_i, f_i \in E(G)$ for all $i = 1, \ldots, p$. 
Let $e_1 = \{i_1, j_1\}$ and $f_1 = \{k_1, \ell_1\}$. 
We may assume that $i_1 \notin \supp(v)$ and $k_1 \notin \supp(u)$.

\medskip
\noindent
{\bf Case 1:} Suppose that $j_1 \notin \supp(v)$, that is $e_1 \notin  \supp(v)$. We distinguish two subcases below.
\vspace{0.2cm}


\noindent
\textit{\textbf{ Sub-case 1.1:}} Suppose that $\ell_1 \notin \supp(u)$, that is $f_1 \notin  \supp(u)$
Then $u = {\bf x}_{e_1} w$ and $v = {\bf x}_{f_1} w$, where $w = {\bf x}_{e_2} \cdots {\bf x}_{e_p}$. 
Since $G$ is connected, there exists a path connecting $i_1$ and $\ell_1$. 
By interchanging the roles of $i_1, j_1$ and $k_1, \ell_1$, we may assume that there exist vertices $a_1, \ldots, a_q$ such that 
\[
P := i_1, j_1, a_1, \ldots, a_q, k_1, \ell_1
\]
is a path of minimal length connecting $i_1$ and $\ell_1$ and passing through $j_1$ and $k_1$. 
In other words, for all $1 \le r < s \le q$, we have $\{a_r, a_s\} \in E(G)$ if and only if $s = r + 1$. 
Moreover, for each $1 \le i \le q$, the vertex $a_i \in \supp(w)$ because $\supp(u) \cup \supp(v) = V(G)$.

If $q = 0$, then $\{j_1, k_1\} \in E(G)$. 
This gives $u' = (u /x_{i_1} )x_{k_1} = (v /x_{\ell_1} )x_{j_1}=x_{k_1} x_{j_1} w \in I^{[p]}$, and hence $u, u', v$ form a path in $G_{I^{[p]}}^{(u,v)}$, as required. 

Now assume that $q > 1$. 
We construct a sequence of edges $e_{t_1}, \ldots, e_{t_s} \subseteq \{e_2, \ldots, e_p\}$ as follows. 
For each $1 \le i \le q$, since $a_i \in \supp(w)$, there exists some $b_i \in \supp(w)$ such that $\{a_i, b_i\} \in \{e_2, \ldots, e_p\}$. 
We set $e_{t_i} = \{a_i, b_i\}$. 
Note that if $b_i = a_{i+1}$ for some $i$, then $e_{t_i} = e_{t_{i+1}}$. 
In this case, we omit $e_{t_{i+1}}$ from our sequence and obtain distinct edges $e_{t_1}, \ldots, e_{t_s}$ such that each $a_i$ appears in exactly one of these edges. 

Now we apply Remark~\ref{rem:switch} to obtain the desired path in $G_{I^{[p]}}^{(u,v)}$ connecting $u$ and $v$. 
In the sequence of edges 
\[
e_1 = e_{t_0}, e_{t_1}, \ldots, e_{t_s}, e_{t_{s+1}} = f_1,
\]
every pair of consecutive edges does not form a gap in $G$. 
By applying Remark~\ref{rem:switch} repeatedly, we obtain a sequence of monomials 
\[
v = v_0, v_1, \ldots, v_s, v_{s+1} = u
\]
such that $v_i$ is obtained from $v_{i-1}$ by switching $e_{t_i}$ with $e_{t_{i-1}}$ for all $i = 1, \ldots, s + 1$. 
This shows that each $v_i$ is connected to $v_{i-1}$ by a path of length two in $G_{I^{[p]}}^{(u,v)}$, and hence $u$ and $v$ are connected in $G_{I^{[p]}}^{(u,v)}$.

\vspace{0.2cm}

\textit{\textbf{ Sub-case 1.2:}}  Suppose that $\ell_1 \in \supp(u)$. 
Set $e_2 = \{\ell_1, j_2\}$. 
Since $\deg(m) = 2p + 2$ and $i_1, j_1 \notin \supp(v)$, we obtain $j_2 \in \supp(v)$ and set $f_2 = \{j_2, \ell_2\}$. 
Continuing in this way, we define 
\[
e_t = \{\ell_{t-1}, j_t\} \quad \text{and} \quad f_t = \{j_t, \ell_t\}
\]
for all $t = 2, \ldots, s$, where $s$ is chosen such that $\ell_s \notin \supp(u)$. 
Then $e_t = f_t$ for all $t$ with $s + 1 \le t \le p$, and $u={\bf x}_{e_1} \cdots {\bf x}_{e_s} u'$, $v={\bf x}_{f_1} \cdots {\bf x}_{f_s} u'$ and $u'={\bf x}_{e_{s+1}} \cdots {\bf x}_{e_p}$. Set 
\[
w_s = (u/x_{\ell_{s-1}})x_{\ell_s}
\quad \text{and} \quad
w_t = (w_{t+1}/x_{\ell_{t-1}})x_{\ell_t}
\quad \text{for all } t = 2, \ldots, s-1.
\]
It follows that $w_2, \ldots, w_s \in \mathcal{G}(I^{[p]})$ and $u, w_s, \ldots, w_2$ form a path in $G_{I^{[p]}}^{(u,v)}$. Moreover, $w_2$ and $v$ satisfy $w_2 = {\bf x}_{e_1} w$ and $v = {\bf x}_{f_1} w$, 
where $w = {\bf x}_{f_2} \cdots {\bf x}_{f_p}$. 
By repeating the argument of Case 1.1, we see that $w_2$ and $v$ are connected in $G_{I^{[p]}}^{(w_2,v)}$. 
Since $\lcm(w_2, v) = \lcm(u, v)$, it follows that $G_{I^{[p]}}^{(w_2,v)}$ is the same as $G_{I^{[p]}}^{(u,v)}$. 
Hence, $u$ and $v$ are connected in $G_{I^{[p]}}^{(u,v)}$.

\medskip
\noindent
{\bf Case 2:} Suppose that $j_1 \in \supp(v)$. 
We distinguish two subcases below.

\medskip
\noindent
\textit{\textbf{ Sub-case 2.1:}}  Suppose that $\ell_1 \notin \supp(u)$. 
This case is identical to Case~1.1 after interchanging the roles of $u$ and $v$.

\medskip
\noindent
\textit{\textbf{ Sub-case 2.2:}} Suppose that $\ell_1 \in \supp(u)$. 
Assume that there exists some $2 \le s \le p$ such that, after rearranging the indices, we have 
\[
e_i = \{\ell_{i-1}, j_i\} \quad \text{and} \quad f_i = \{j_{i-1}, \ell_i\}
\]
for all $2 \le i \le s$, and $j_s = \ell_s$. Then 
\[
u = (x_{i_1}x_{j_1})(x_{\ell_1}x_{j_2}) \cdots (x_{\ell_{s-1}}x_{j_s}){\bf x}_{e_{s+1}} \cdots {\bf x}_{e_p}
\quad \text{and} \quad
v = (x_{k_1}x_{\ell_1})(x_{j_1}x_{\ell_2}) \cdots (x_{j_{s-1}}x_{\ell_s}){\bf x}_{f_{s+1}} \cdots {\bf x}_{f_p}.
\]
with $j_s = \ell_s$. Set $u' = (u/x_{i_1})x_{k_1}$. 
Then $u' \in I^{[p]}$, and we obtain a path $u, u', v$ connecting $u$ and $v$ in $G_{I^{[p]}}^{(u,v)}$.

Now suppose that no such $s$ exists. 
In particular, $j_1 \ne \ell_1$. 
Since $\ell_1 \in \supp(u)$ and $j_1 \in \supp(v)$, there exist $j_2$ and $\ell_2$ such that we may set 
\[
e_2 = \{\ell_1, j_2\} \quad \text{and} \quad f_2 = \{j_1, \ell_2\}.
\]
By our assumption, $j_2 \ne \ell_2$. 
Continuing in this way, let $s$ be the minimal integer for which $j_s \notin \supp(v)$, and  
\[
e_i = \{\ell_{i-1}, j_i\}, \qquad f_i = \{j_{i-1}, \ell_i\} \quad \text{for all } 2 \le i \le s.
\]
Here we assume that at each step $\ell_i \in \supp(u)$ for all $i = 1, \ldots, s-1$; otherwise, we may interchange the roles of $u$ and $v$. 

If $s$ is even, set $u' = (u/x_{j_s})x_{k_1}$. 
We can arrange the variables in $u'$ so that 
\[
u' = {\bf x}_{e_1}{\bf x}_{f_1}{\bf x}_{e_3}{\bf x}_{f_3} \cdots {\bf x}_{e_{s-1}}{\bf x}_{f_{s-1}}
{\bf x}_{e_{s+1}} \cdots {\bf x}_{e_p} \in \mathcal{G}(I^{[p]}).
\]
Then $u, u', v$ form a path in $G_{I^{[p]}}^{(u,v)}$, as required. 

If $s$ is odd, construct $v' = (v/x_{\ell_{s-1}})x_{i_1}$. 
We can arrange the variables in $v'$ so that 
\[
v' = {\bf x}_{f_1}{\bf x}_{e_1}{\bf x}_{f_3}{\bf x}_{e_3} \cdots {\bf x}_{f_{s-2}}{\bf x}_{e_{s-2}}
{\bf x}_{f_s} \cdots {\bf x}_{f_p} \in \mathcal{G}(I^{[p]}).
\]
Then $v'$ and $v$ are adjacent in $G_{I^{[p]}}^{(u,v)}$. 
Also note that $e_{j_s} \notin \supp(v')$. By replacing the role of $v$ with $v'$ in Case~1, we obtain the desired path connecting $v'$ and $u$. 
This completes the proof.
\end{proof}


%%%%%%%%%%%%%%%%%%%%%%%%%%%%%%%%%%%%%%%%%%%%%%%%%%%%%%%%%%%%%%%%%%%%%%%%%%%%%%%%%%%%%%%%%%%%%%%%%%%%%%%%%%%%%%%%%%%%%%%%%%%%%%%%%%%%%%%%%%%%%%%%%%%%%%%%%%%%%%%%%%%%%%%%%%%%%%%%%%%%%%%%%%%%%%%%%%%%%%%%%%%%%%%%%%%%%%%%%%%%%%%%%%%%%%%%%%%%%%%%%%%%%%%%%%%%%%%%%%%%%%%%%%%%%%%%%%%%%%%%%%%%%%%%%%%%%%%%%%%%%%%%%%%%%%%%%%%%%%%%%%%%%%%%%%%%%%%%%%%%%%%%%%%%%%%%%%%%%%%%%%%%%%%%%%%%%%%%%%%%%%%%%%%%%%%%%

 We now investigate $\beta_{1,j}(I(G)^{[p]})$, with particular emphasis on determining for which degrees $j$ they vanish. To this end, we first introduce the following notation. Let $u,v \in \mathcal{G}(I(G))^{[p]})$ be two distinct monomials. We fix a presentation of $u$ and $v$ given by $
u = {\bf x}_{e_1} \cdots {\bf x}_{e_p}$ and $v = {\bf x}_{f_1} \cdots {\bf x}_{f_p},$
where each $e_i$ and $f_i$ is an edge of $G$. We denote by $d(u,v)$ the maximum number of common edges in $\{e_1, \ldots, e_p\}$ and $\{f_1, \ldots, f_p\}$.
  
	\begin{Proposition}\label{thm:betti}
		Let $G$ be a graph and $2\leq p \leq \nu(G)$. Then $\beta_{1,j}(I(G)^{[p]}) = 0$ if $j \notin \{ 2p+1,2p+2\}$.
	\end{Proposition}
		
		\begin{proof}
Let $I = I(G)$. 
Consider the Taylor complex associated with $I^{[p]}$ (see \cite[Section~7.1]{HHmonomialideals}). 
From its construction, we observe that $\beta_{1,m}(I^{[p]}) = 0$ unless $m = \lcm(u_1, u_2)$ for some $u_1, u_2 \in G(I^{[p]})$. 
Furthermore, by a result of Gasharov, Peeva, and Welker~\cite{Gasharov1999}, for all $m \in L(I^{[p]})$, we have
\[
\beta_{1,m}(I^{[p]}) = \dim_K \tilde{H}_0(\Delta((1,m)); K).
\]
Recall that $\dim_K \tilde{H}_0(\Delta((1,m)); K) = c - 1$, where $c$ denotes the number of connected components of $\Delta((1,m))$ (see \cite[Chapter~1, Section~7]{MJ}). 
Since the largest possible degree of $m$ is $4p$, it follows that $\beta_{1,j}(I^{[p]}) = 0$ for all $j > 4p$. 

Hence, to establish the claim, it suffices to prove the following:  
if $m$ is the least common multiple of two minimal generators of $I^{[p]}$, with $2p + 1 < \deg(m) \le 4p$ and $\deg(m) \ne 2(p + 1)$, then the open interval $(1, m)$ in the lcm-lattice $L(I^{[p]})$ is connected. 
To show that the interval $P = (1, m)$ is connected, it is enough to verify that its comparability graph $G_P$ is connected.

By definition, the vertices of $G_P$ correspond to the monomials in $L(I^{[p]})$ that are different from $1$ and strictly divide $m$. 
Every monomial in $V(G_P)$ of degree greater than $2p$ is adjacent to at least one vertex of degree $2p$, and these degree $2p$ vertices are exactly the minimal generators of $I^{[p]}$. 
Therefore, it suffices to prove that for any two distinct $u, v \in V(G_P)$ with $\deg(u) = \deg(v) = 2p$, there exists a path in $G_P$ connecting $u$ and $v$. 
If $\lcm(u,v) \in V(G_P)$, then $\lcm(u,v)$ serves as a common neighbor of $u$ and $v$, and the claim follows immediately. 
Otherwise, since both $u$ and $v$ strictly divide $m$, we have $\lcm(u,v) \notin V(G_P)$ precisely when $\lcm(u,v) = m$.

Fix presentations of $u$ and $v$ given by $u = {\bf x}_{e_1} \cdots {\bf x}_{e_p}$ and $v = {\bf x}_{f_1} \cdots {\bf x}_{f_p}$, where each $e_i$ and $f_i$ is an edge of $G$. After a suitable relabeling of indices, we may assume that $\{e_1, \ldots, e_k\} \cap \{f_1, \ldots, f_k\} = \emptyset$ and $\{e_{k+1}, \ldots, e_p\} = \{f_{k+1}, \ldots, f_p\}.$ Let $m = \lcm(u,v)$, and assume that $2p + 2 < \deg(m) < 4p$. 
Then $k > 2$. 

Let $A = \{e_1, \ldots, e_k\}$ and $B= \{f_1, \ldots, f_k\}$. 
Consider the bipartite graph $H$ on the vertex set $A \cup B$ with
\[
E(H) = \bigl\{ \{e_i, f_j\} : e_i \cap f_j \ne \emptyset \bigr\}.
\]
Note that $H$ does not contain multiple edges, since $A$ and $B$ have no common elements. 
Moreover, each vertex in $H$ has degree at most two. 
Observe that there must exist at least one vertex in $H$ of degree at most one; otherwise, $\supp(u) = \supp(v)$ and hence $u = v$, a contradiction. Without loss of generality, we may assume that $\deg(e_k) \le 1$. Moreover, if $\deg(e_k) = 1$, then we may assume that $f_k$ is the unique neighbor of $e_k$. Define $v_1 = (v / {\bf x}_{f_k}) {\bf x}_{e_k}.$ By the choice of $e_k$, we have $e_k \cap \bigl(\supp(v) \setminus \{f_k\}\bigr) = \emptyset$, and therefore $v_1 \in \mathcal{G}(I^{[p]})$. 
Furthermore:
\begin{itemize}
    \item[(i)] If $\deg(e_k) = 0$, then $\deg(\lcm(v, v_1)) = 2p + 2 < \deg(m)$.
    \item[(ii)] If $\deg(e_k) = 1$, then $\deg(\lcm(v, v_1)) = 2p + 1 < \deg(m)$.
\end{itemize}

In both cases \textup{(i)} and \textup{(ii)} above, the monomial $\lcm(v, v_1)$ strictly divides $m$. 
Hence, $\lcm(v, v_1) \in V(G_P)$, and we obtain a path 
\[
v, \ \lcm(v, v_1), \ v_1
\]
in $G_P$. 
Moreover, we have $d(u, v)=p-k < d(u, v_1)=p-k+1$. If $k=2$, then 
\[
v, \ \lcm(v, v_1), \ v_1, \lcm(v_1, u), \ u
\]
forms a path in $G_P$ connecting $u$ and $v$. If $k>2$, then we set $v_0=v$ and repeat the above argument by replacing the role of $v_i$ with $v_{i+1}$, for $i\ge 0$. After  $j$ steps where $2 \leq j\leq k-1$, we obtain monomials $v_2, \ldots, v_{j}$ such that 
\[
v, \ \lcm(v, v_1), \ v_1, \lcm(v_1, v_2), \ v_2\ \ldots, \ v_{j}, \ \lcm(v_{j}, u), \ u
\]
forms a path in $G_P$ connecting $u$ and $v$.
\end{proof}



%%%%%%%%%%%%%%%%%%%%%%%%%%%%%%%%%%%%%%%%%%%%%%%%%%%%%%%%%%%%%%%%%%%%%%%%%%%%%%%%%%%%%%%%%%%%%%%%%%%%%%%%%%%%%%%%%%%%%%%%%%%%%%%%%%%%%%%%%%%%%%%%%%%%%%%%%%%%%%%%%%%%%%%%%%%%%%%%%%%%%%%%%%%%%%%%%%%%%
   Now we present the proof of Theorem~\ref{Thm:linrelated1}. Equivalently, we prove the following. 
   
    \begin{Theorem}\label{Thm:linrelated}
Let $G$ be a graph and $2\leq p \leq \nu(G)$. Then the following statements are equivalent.
\begin{enumerate}
    \item $I(G)^{[p]} $ is not linear related.
    \item There exists an induced subgraph $H$ of $G$ on $2p+2$ vertices such that $H$ is disconnected with $\nu(H)=p+1$. 
\end{enumerate} 
\end{Theorem}
\begin{proof}
(2) $\implies$ (1) Assume that there exist an induced subgraph $H$ of $G$ on $2p+2$ vertices such that $H$ is disconnected and $\nu(H)=p+1$. Let $M=\{e_1, \ldots, e_{p+1}\}$ be a matching of $H$ of size $p+1$ and $H_1$ and $H_2$ be subgraphs of $H$ such that $V(H)$ is disjoint union of $V(H_1)$ and $V(H_2)$, and $E(H)$ is disjoint union of $E(H_1)$ and $E(H_2)$. Since $|V(H)|=2p+2$ and $\nu(H)=p+1$, it follows that $M\cap E(H_i) \neq \emptyset$ for $i=1,2$. We may assume that $e_1 \in E(H_1)$ and $e_{p+1}\in E(H_2)$. Our assumption on $|V(H)|$ and $\nu(H)$ yields that $H_1$ and $H_2$ have an even number of vertices. Then for some positive integers $a$ and $b$, we have  $|V(H_1)|=2a$ and $|V(H_2)|=2b$ where $2a+2b=2p+2$. 

Let $w = {\bf x}_{e_2} \cdots {\bf x}_{e_{p}}$. Then $u = {\bf x}_{e_1} w$ and $v = {\bf x}_{e_{p+1}} w$ belong to $\mathcal{G} (I^{[p]})$. To show that $I$ is not linearly related, we apply Theorem~\ref{ref:BHN} and show that $G^{(u,v)}_{I^{[p]}}$ is disconnected. Since $H$ is an induced subgraph of $G$ on vertices appearing in $M$, it is enough to show that $G^{(u,v)}_{J^{[p]}}$ is disconnected, where $J=I(H)$. Since $\lcm(u,v)= {\bf x}_{e_1}  \cdots {\bf x}_{e_{p+1}}$. We first observe that $G^{(u,v)}_{J^{[p]}}= G_{J^{[p]}}$. Moreover, for any monomial $m \in \mathcal{G}(J^{[p]})$, one of the following conditions holds. 
\begin{enumerate}
\item[{(i)}] $|\supp(m) \cap V(H_1)| = 2(a-1)$ and $|\supp(m) \cap V(H_1) |= 2b$ 
\item[{(ii)}]$|\supp(m) \cap V(H_1) |= 2a$ and $|\supp(m) \cap V(H_1) |= 2(b-1)$ 
\end{enumerate}

Let $A$ be the set of all monomials in $\mathcal{G}(J^{[p]})$ satisfying (i) and $B$ be the set of all monomials in $\mathcal{G}(J^{[p]})$ satisfying (ii). It follows immediately that none of the monomials in $A$ is adjacent to any monomial in $B$, and therefore the monomials in $A$ and $B$ belong to different connected components of $G_J$. Since $u\in A$ and $v\in B$,  we obtain the desired conclusion.


(1) $\implies$ (2) Assume that $I^{[p]}$ is not linearly related. It follows from Proposition~\ref{thm:betti} that $\beta_{1,2p+2}(I^{[p]}) \neq 0$, and following the proof of Proposition~\ref{thm:betti} and the notations introduced therein, we obtain two monomials $u$ and $v$  in $\mathcal{G}(I^{[p]})$ such that the degree of $m=\lcm(u,v)$ is $2p+2$ and there is no path connecting $u$ and $v$ in the comparability graph $G_P$ of $P=(1,m)$. Let $H$ be the induced subgraph of $G$ on the vertex set $\supp(u)$ and $\supp(v)$, and set $J=I(H)$. Then $H$ has $2p+2$ vertices. Also, $u,v \in \mathcal{G}(J^{[p]})$, and $\beta_{1,2p+2}(J^{[p]}) \neq 0$. This shows that $J^{[p]}$ is not linearly related. By Proposition~\ref{prop:connected}, it follows that $H$ is disconnected. It only remains to show that $\nu(H)=p+1$. Since $u,v \in J^{[p]}$ and $H$ has $2p+2$ vertices, it follows that $p+1\geq \nu(H)\geq p$. From \cite[Theorem 5.1]{BHZ}, it follows that $J^{[\nu(H)]}$ has linear resolution. This gives $
\nu(H)=p+1$, as required. 
\end{proof}


%%%%%%%%%%%%%%%%%%%%%%%%%%%%%%%%%%%%%%%%%%%%%%%%%%%%%%%%%%%%%%%%%%%%%%%%%%%%%%%%%%%%%%%%%%%%%%%%%%%%%%%%%%%%%%%%%%%%%%%%%%%%%%%%%%%%%%%%%%%%%%%%%%%%%%%%%%%%%%%%%%%%%%%%%%%%%%%%%%%%%%%%%%%%%%%%%%



% \subsection{When $I_{\Delta}^{[p]}$ has a linear first syzygy module} We begin by examining the property of being linearly related and specialize Theorem \ref{Thm:linrelated} to this setting in the following corollary. 

% \begin{Corollary}\label{coro: lin related}
% Let $\Delta$ be a $1$-dimensional flag simplicial complex on $[n]$, and let
% $2 \leq p \leq \nu(\overline{\Delta})$. The following conditions are equivalent:
% \begin{enumerate}
%     \item $I_{\Delta}^{[p]}$ is linearly related.
%     \item $\Delta$ does not contain an induced subgraph isomorphic to $K_{2+i,\,2p-i}$ for some even integer $i$ with $0 \leq i \leq p$.
% \end{enumerate}

% \begin{table}[h]
% \centering
% \begin{tabular}{c|l}
% $p$ & Induced subgraphs \\ \hline
% 2 & $K_{2,4}$ \\
% 3 & $K_{2,6},\; K_{4,4}$ \\
% 4 & $K_{2,8},\; K_{4,6}$ \\
% 5 & $K_{2,10},\; K_{4,8},\; K_{6,6}$\\
% 6 & $K_{2,12}$, \; $K_{4,10}$ \; $K_{6,8}$\\
% 7 & $K_{2,14}$, \; $K_{4,12}$ \; $K_{6,10}$ \; $K_{8,8}$\\
% \end{tabular}
% \caption{Complete bipartite graphs $K_{2+i,\,2p-i}$ with $i$ even and $0 \leq i \leq p$.}
% \end{table}
% \end{Corollary}

We can now give a proof of Corollary \ref{coro: lin related}.

\begin{proof}[Proof of Corollary \ref{coro: lin related}]
$(1) \Rightarrow (2)$.  
Suppose that $\Delta$ contains an induced subgraph $H$ isomorphic to $K_{2+i,\,2p-i}$ for some even integer $i$ with $0 \leq i \leq p$. Then the induced subgraph of $\overline{\Delta}$ on $V(H)$ is disconnected and satisfies $\nu(H)=p+1$. Since $I_{\Delta}^{[p]} = I(\overline{\Delta})^{[p]}$, it follows from Theorem~\ref{Thm:linrelated} that $I_{\Delta}^{[p]}$ is not linearly related, which is a contradiction.

$(2) \Rightarrow (1)$.  
Assume that $I_{\Delta}^{[p]}$ is not linearly related. By Theorem~\ref{Thm:linrelated}, there exists an induced subgraph $H$ of $\overline{\Delta}$ on $2p+2$ vertices such that $H$ is disconnected and $\nu(H)=p+1$. Let
\[
V(H)=\{v_1,\dots, v_{2(p+1)}\}
\quad \text{and} \quad
M=\bigl\{\{v_k,v_{k+1}\} : k=1,\dots,2p+1 \text{ odd}\bigr\}.
\]
be the matching of $H$ of size $p+1$. We know that $\Delta$ does not contain any triangle, equivalently, $\overline{\Delta}$ does not contain three isolated vertices. Since $H$ is disconnected, it yields that $H$ has exactly two connected components, say $H_1$ and $H_2$. Moreover, given any vertices $v_r, v_s  \in H_1$ and $v_t \in H_2$, since $v_r$ and $v_s$ are not adjacent with $v_t$, we must have $\{v_r,v_s\} \in H_1$. Hence $H_1$ is a complete graph, and similarly $H_2$ is also a complete graph. Let $|M\cap E(H_1)| = j$ and $|M\cap E(H_2)| = p+1-j$. Then it follows that $H_1$ is a complete graph on $2j$ vertices and $H_2$ is a complete graph on $2p+2-2j$ vertices. Equivalently, $\Delta$ contains an induced subgraph isomorphic to $K_{2j, 2p+2-2j}$, as required.
\end{proof}

\begin{Example}\rm
Consider the $1$-dimensional flag simplicial complex $\Delta$ in Figure \ref{fig: 1-dimensional flag simplicial complex}.

\begin{figure}[h]
    \centering
    \includegraphics[scale=0.6]{graph_lin_rel.eps}
    \caption{A $1$-dimensional flag simplicial complex.}
    \label{fig: 1-dimensional flag simplicial complex}
\end{figure}

By Corollary~\ref{coro: lin related}, $I_{\Delta}^{[2]}$ is not linearly related, since, for instance, the induced subgraph of $\Delta$ on the vertex set $\{1,2,4,5,6,7\}$ is isomorphic to $K_{2,4}$. Using \textit{Macaulay2} \cite{M2}, the Betti table of $I_{\Delta}^{[2]}$ over a field of characteristic zero is the following:
\[
\begin{array}{r|cccccc}
      & 0 & 1 & 2 & 3 & 4 & 5 \\
\hline
\mathrm{total:} & 91 & 336 & 516 & 406 & 162 & 26 \\
0: & . & . & . & . & . & . \\
1: & . & . & . & . & . & . \\
2: & . & . & . & . & . & . \\
3: & . & . & . & . & . & . \\
4: & 91 & 329 & 497 & 389 & 157 & 26 \\
5: & . & 7 & 19 & 17 & 5 & .
\end{array}
\]

It is worth noting that $\beta_{1,6}$, which detects the failure of the linearly related property of $I_{\Delta}^{[2]}$, coincides with the number of induced subgraphs of $\Delta$ that are isomorphic to $K_{2,4}$. In this case, there are exactly $7$ such subgraphs, induced by the vertex sets $\{1,2,3,4,5,6\}$, $\{1,2,4,5,6,8\}$, $\{1,2,4,5,7,8\}$, $\{1,2,4,6,7,8\}$, $\{1,2,5,6,7,8\}$, $\{2,3,5,6,7,8\}$, and $\{1,3,5,6,7,8\}$.

This phenomenon is not restricted to the present example, but holds in general, as we show in the next theorem. More precisely, we provide a combinatorial description of the Betti number $\beta_{1,2p+2}(I_{\Delta}^{[p]})$, which measures the failure of the linearly related property of $I_{\Delta}^{[p]}$, in terms of the number of induced subgraphs isomorphic to $K_{2+i,\,2p-i}$, with $0\leq i\leq p$ even.
\end{Example}


\begin{Theorem}\label{Thm: betti 2,2p+2}
Let $\Delta$ be a $1$-dimensional flag simplicial complex on $[n]$ and let
$2 \leq p \leq \nu(\overline{\Delta})$. %For each even integer $i$ with $0 \leq i \leq p$, let
% \[
% N_i(\Delta)
% =
% \{\text{induced subgraphs of } \Delta \text{ isomorphic to } K_{2+i,\,2p-i}\}.
% \]
Then
\[
\beta_{1,2p+2}(I_{\Delta}^{[p]})
=
\sum_{\substack{0 \leq i \leq p \\ i \text{ even}}} \big\vert\{\text{induced subgraphs of } \Delta \text{ isomorphic to } K_{2+i,\,2p-i}\}\big\vert.% \vert N_i(\Delta)\vert.
\]
\end{Theorem}

\begin{proof}
If $\Delta$ does not contain an induced subgraph isomorphic to $K_{2+i,\,2p-i}$ for some even integer $i$ with $0 \leq i \leq p$, there is nothing to prove, because $I_{\Delta}^{[p]}$ is linearly related by Corollary \ref{coro: lin related} and so $\beta_{1,2p+2}(I_{\Delta}^{[p]})=0$, as desired.

Assume that $\Delta$ contains an induced subgraph isomorphic to $K_{2+i,\,2p-i}$ for some even integer $i$ with $0 \leq i \leq p$. We divide the proof into two cases, depending on whether $\Delta$ is of the form $K_{2+i,\,2p-i}$ or not.\\


\textbf{Case 1)} We first prove that the claim holds when $\Delta = K_{2+i,\,2p-i}$ for some even
integer $i$ with $0 \leq i \leq p$. In this case, we show that
\[
\beta_{1,2p+2}(I_{\Delta}^{[p]}) = 1.
\]
By Hochster's formula, we have
\[
\beta_{1,2p+2}(I_{\Delta}^{[p]}) = \dim_K \tilde{H}_{2p-1}(\Delta^{[p]}),
\]
where $\tilde{H}_{2p-1}(\Delta^{[p]}) = \ker(\delta_{2p-1}) / \mathrm{Im}(\delta_{2p})$.
Observe that $C_{2p} = (0)$, since there are no $2p$-dimensional facets by Proposition \ref{Thm: dimension}; hence $\mathrm{Im}(\delta_{2p}) = (0)$. Therefore, it suffices to show that there is a unique nonzero $(2p-1)$-dimensional cycle in $\Delta^{[p]}$, up to factors in $K$.\\

\textbf{\textit{Existence.}} Let us examine the $(2p-1)$-dimensional facets of $\Delta^{[p]}$ by using Theorem \ref{Thm: Facet, pure, CM}. Indeed, a complete bipartite graph $K_{j,\,2p-j}$, with $j$ an odd integer and $1 \leq j \leq p$, is an induced subgraph of $K_{2+i,\,2p-i}$ if and only if $j \leq 2+i$ and $2p-j \leq 2p-i$, which is equivalent to $0 < j-i \leq 2$. Since $j$ is odd and $i$ is even, this condition is equivalent to $j = i+1$. 


Assume that $V(\Delta)=\{1,\dots, i+2\} \sqcup \{i+3,\dots,2p+2\}$ and $
E(\Delta)=\bigl\{\{h,k\} \mid h\in [i+2],\ k\in [2p+2]\setminus[i+2]\bigr\}.$ The $(2p-1)$-dimensional facets of $\Delta^{[p]}$ are then given by
\[
\left[1,\dots,\widehat{h},\dots,i+2,\,i+3,\dots,\widehat{k},\dots,2p+2\right],
\]
for all $h\in [i+2]$ and $k\in [2p+2]\setminus[i+2]$, where $\widehat{h}$ and $\widehat{k}$ denote the removal of the vertices $h$ and $k$, respectively. 
For example, consider $\Delta = K_{2,4}$, and assume that 
\[
V(\Delta)=\{1,2\}\cup \{3,4,5,6\}
\quad \text{and} \quad
E(\Delta)=\bigl\{\{i,j\} : i=1,2,\ j=3,4,5,6\bigr\}.
\]
We can represent the eight $3$-dimensional facets of $\Delta^{[2]}$ as
\begin{align*}
\begin{aligned}
&[\widehat{1},2,\widehat{3},4,5,6]=[2,4,5,6]\\
&[\widehat{1},2,3,\widehat{4},5,6]=[2,3,5,6]\\
&[\widehat{1},2,3,4,\widehat{5},6]=[2,3,4,6]\\
&[\widehat{1},2,3,4,5,\widehat{6}]=[2,3,4,5]
\end{aligned}
\qquad
\begin{aligned}
&[1,\widehat{2},\widehat{3},4,5,6]=[1,4,5,6]\\
&[1,\widehat{2},3,\widehat{4},5,6]=[1,3,5,6]\\
&[1,\widehat{2},3,4,\widehat{5},6]=[1,3,4,6]\\
&[1,\widehat{2},3,4,5,\widehat{6}]=[1,3,4,5]
\end{aligned}
\end{align*}

Define the following nonzero vector in $C_{2p-1}$:
\[
\Gamma = \sum_{h=1}^{i+2} \sum_{k=i+3}^{2p+2} (-1)^{h+k} 
\left[1,\dots,\widehat{h},\dots,i+2,\,i+3,\dots,\widehat{k},\dots,2p+2 \right],
\]
and we aim to show that $\Gamma$ is a $(2p-1)$-cycle, that is, $\delta_{2p-1}(\Gamma)=0$.

In our example, $\Gamma$ is explicitly given by
\begin{align*}
    &+[\widehat{1},2,\widehat{3},4,5,6]
    -[\widehat{1},2,3,\widehat{4},5,6]
    +[\widehat{1},2,3,4,\widehat{5},6]
    -[\widehat{1},2,3,4,5,\widehat{6}] \\
    &-[1,\widehat{2},\widehat{3},4,5,6]
    +[1,\widehat{2},3,\widehat{4},5,6]
    -[1,\widehat{2},3,4,\widehat{5},6]
    +[1,\widehat{2},3,4,5,\widehat{6}].
\end{align*}

\noindent The sign is determined by the parity of the sum of the indices being removed. Applying $\delta_3$ to $\Gamma$ produces $32$ two-dimensional facets. Rather than listing all of them, we explain why their sum cancels out. A $2$-dimensional facet appearing in $\delta_3(\Gamma)$ is obtained by removing three elements from $[1,2,3,4,5,6]$: two of them are of the form $h$ and $k$, with $h \in \{1,2\}$ and $k \in \{3,4,5,6\}$, while the third element can be any remaining index distinct from $h$ and $k$.

\noindent For instance, the facet
\[
[4,5,6] = [\widehat{1},\widehat{2},\widehat{3},4,5,6]
\]
appears twice in $\delta_3(\Gamma)$:
\begin{itemize}
\item from $[\widehat{1},2,\widehat{3},4,5,6]$, obtained by removing $1$ and $3$ from $[1,2,3,4,5,6]$, and then removing $2$;
\item from $[1,\widehat{2},\widehat{3},4,5,6]$, obtained by removing $2$ and $3$ from $[1,2,3,4,5,6]$, and then removing $1$.
\end{itemize}
These two contributions appear with opposite signs and therefore cancel each other.



In general, we have
\begin{equation*}
\begin{aligned}
\delta_{2p-1}(\Gamma)
 = {} & \delta_{2p-1}\left(
 \sum_{h=1}^{i+2} \sum_{k=i+3}^{2p+2} (-1)^{h+k}
 \left[1,\dots,\widehat{h},\dots,i+2,i+3,\dots,\widehat{k},\dots,2p+2 \right]
 \right)= \\
={} & \sum_{\ell=1}^{2p} (-1)^{\ell+1} \left(\sum_{h=1}^{i+2} \sum_{k=i+3}^{2p+2} (-1)^{h+k}
\left[1,\dots,\widehat{h},\dots,\widehat{a_\ell},\dots,i+2,i+3,\dots,\widehat{k},\dots,2p+2 \right]
\right),
\end{aligned}
\end{equation*}
where $a_\ell$ ranges over
$([i+2]\setminus\{h\}) \cup (([2p+2]\setminus[i+2])\{k'\})$. 

For simplicity, denote
$$
\gamma(h,k;\ell)=\left[1,\dots,\widehat{h},\dots,\widehat{a_\ell},\dots,i+2,i+3,\dots,\widehat{k},\dots,2p+2 \right]
$$
where $h\in [i+2]$, $k \in [2p+2]\setminus[i+2]$ and $\ell \in [2p]$. In conclusion, we have
\begin{equation}\label{eq1}
\delta_{2p-1}(\Gamma)= \sum_{\ell=1}^{2p} \sum_{h=1}^{i+2} \sum_{k=i+3}^{2p+2} (-1)^{h+k+\ell+1} \gamma(h,k;\ell).
\end{equation}

We now study the terms in Equation~\eqref{eq1}. Note that the occurrence $\gamma(h,k;l)$ can appear in Equation \eqref{eq1} in two different ways. 
\begin{itemize}
    \item If $a_{\ell}\in [i+2]\setminus\{h\}$, then one occurrence is obtained by deleting
$h$ and $a_{\ell}$ from $$\left[1,\dots,h,\dots, a_{\ell}, \dots, i+2,\,i+3,\dots,\widehat{k},\dots, 2p+2\right],$$ while the second occurrence is obtained by exchanging the roles of $h$ and $a_{\ell}$;
in both cases, $k$ is kept fixed.
    \item Similarly, if $a_{\ell}\in ([2p+2]\setminus[i+2])\setminus\{k\}$, one occurrence is obtained by deleting
$k$ and $a_{\ell}$ from $$\left[1,\dots,\widehat{h}, \dots, i+2,\,i+3,\dots,a_{\ell},\dots,k,\dots, 2p+1\right],$$ and the second one by exchanging the roles of $k$ and $a_{\ell}$; while $h$ is kept fixed.
\end{itemize}

Fix $h_1,h_2 \in [i+2]$, $k_1,k_2 \in [2p+2]\setminus[i+2]$, and $\ell_1,\ell_2 \in [2p]$ such that $\gamma(h_1,k_1;\ell_1)=\gamma(h_2,k_2;\ell_2)$ and we study 
$$
(-1)^{h_1+k_1+\ell_1+1}\gamma(h_1,k_1;\ell_1)
+
(-1)^{h_2+k_2+\ell_2+1}\gamma(h_2,k_2;\ell_2).
$$
We distinguish the following cases, according to whether $a_{\ell_1}$ ranges over $([i+2]\setminus\{h_1\})$ or $([2p+2]\setminus[i+2])\setminus\{k_1\}$.\\

(1) Let $a_{\ell_1} \in [i+2]\setminus \{h_1\}$. If $a_{\ell_1} > h_1$, then
\[
\begin{cases}
h_2 = \ell_1 + 1, \\
k_2 = k_1, \\
\ell_2 = h_1.
\end{cases}
\quad \Rightarrow \quad
(-1)^{h_1+k_1+\ell_1+1}\gamma(h_1,k_1;\ell_1)
+
(-1)^{h_2+k_2+\ell_2+1}\gamma(h_2,k_2;\ell_2)=0
\]

Similarly, if $a_{\ell_1} < h_1$, then
\[
\begin{cases}
h_2 = \ell_1, \\
k_2 = k_1, \\
\ell_2 = h_1 - 1.
\end{cases}
\quad \Rightarrow \quad
(-1)^{h_1+k_1+\ell_1+1}\gamma(h_1,k_1;\ell_1)
+
(-1)^{h_2+k_2+\ell_2+1}\gamma(h_2,k_2;\ell_2)=0
\]

(2) Let $a_{\ell_1} \in [2p-i]\setminus \{k_1\}$. If $a_{\ell_1} > k_1$, then
\[
\begin{cases}
h_2 = h_1, \\
k_2 = \ell_1 + 2, \\
\ell_2 = k_1-1.
\end{cases}
\quad \Rightarrow \quad
(-1)^{h_1+k_1+\ell_1+1}\gamma(h_1,k_1;\ell_1)
+
(-1)^{h_2+k_2+\ell_2+1}\gamma(h_2,k_2;\ell_2)=0
\]

Similarly, if $a_{\ell_1} < k_1$, then
\[
\begin{cases}
h_2 = h_1, \\
k_2 = \ell_1+1, \\
\ell_2 = k_1 - 2.
\end{cases}
\quad \Rightarrow \quad
(-1)^{h_1+k_1+\ell_1+1}\gamma(h_1,k_1;\ell_1)
+
(-1)^{h_2+k_2+\ell_2+1}\gamma(h_2,k_2;\ell_2)=0
\]

Therefore, in Equation~\eqref{eq1}, the terms cancel pairwise, and hence $\delta_{2p-1}(\Gamma)=0$, as desired.\\


\textbf{\textit{Uniqueness.}} We now need to prove that $\ker(\delta_{2p-1})\subseteq \mathrm{span}_K\{\Gamma\}$.  Let $\Lambda\in \ker(\delta_{2p-1})\setminus\{0\}$.  Consider 
$$ \Lambda=\sum_{h=1}^{i+2} \sum_{k=1}^{2p-i}\lambda_{h,k}\left[1,\dots,\widehat{h},\dots,i+2,i+3,\dots,\widehat{k},\dots,2p+2 \right],$$
where $\lambda_{h,k}\in K$ and $\lambda_{h,k}\neq 0$ for some $h\in [i+2]$ and $k\in[2p+2]\setminus[i+2]$. Then, we have
\begin{equation}\label{eq2}
\delta_{2p-1}(\Lambda)=\sum_{\ell=1}^{2p} \sum_{h=1}^{i+2} \sum_{k=1}^{2p-i} (-1)^{\ell+1} \lambda_{h,k}
\left[1,\dots,\widehat{h},\dots, \widehat{a_\ell},\dots,i+2,i+3,\dots,\widehat{k},\dots,2p+2 \right]=0.
\end{equation}

We divide the proof into two claims.\\

\textbf{Claim 1.} We first show that $\lambda_{h,k} \neq 0$ for all $h \in [i+2]$ and $k \in [2p+2]\setminus[i+2]$.

\textit{Proof of Claim 1.} Indeed, suppose by contradiction that there exist $h_1 \in [i+2]$ and $k_1 \in [2p+2]\setminus[i+2]$ such that $\lambda_{h_1,k_1}=0$. This implies that an occurrence of
\[
\left[1,\dots,\widehat{h_1},\dots,\widehat{a_\ell},\dots,i+2,i+3,\dots,\widehat{k_1},\dots,2p+2 \right],
\]
for every $\ell \in [2p]$, does not appear in Equation~\eqref{eq2}. Consequently, the corresponding term obtained by switching the role of either $h_1$ or $k_1$ with $a_\ell$ cannot appear in Equation~\eqref{eq2} either; otherwise, it would not cancel out, and therefore $\delta_{2p-1}(\Lambda)$ would be nonzero, a contradiction. Hence, we obtain $\lambda_{h,k_1}=0$ for all $h\in [i+2]\setminus \{h_1\}$ and $\lambda_{h_1,k}=0$ for all $k\in [2p+2]\setminus([i+2]\cup \{k_1\})$. Now fix an arbitrary $h\in [i+2]$. Then $\lambda_{h,k_1}=0$ by the previous argument; moreover, applying the same argument with $h$ and $k_1$ in place of $h_1$ and $k_1$, respectively, we obtain $\lambda_{h,k}=0$ for all $k\in [2p+2]\setminus([i+2]\cup \{k_1\})$. In conclusion, $\lambda_{h,k}=0$ for all $h\in [i+2]$ and $k\in [2p+2]\setminus[i+2]$, that is, $\Lambda=0$, a contradiction. Hence, Claim~1 is proved. $\square$\\


\textbf{Claim 2.} There exists $c\in K\setminus\{0\}$ such that $\lambda_{h,k} = (-1)^{h+k} c$, for all $h\in [i+2]$ and $k\in [2p+2]\setminus[i+2]$. 

\textit{Proof of Claim 2.} From the fact that $\lambda_{h,k} \neq 0$ for all $h \in [i+2]$ and $k \in [2p+2]\setminus[i+2]$, we have that the $(2p-2)$-dimensional facet
\[
\left[1,\dots,\widehat{h},\dots, \widehat{a_\ell},\dots,i+2,\dots,\widehat{k},\dots,2p+2 \right]
\]
appears in Equation~\ref{eq2}, for all $h \in [i+2]$, $k \in [2p+2]\setminus[i+2]$ and $\ell \in [2p]$. Moreover, note that a $(2p-2)$-dimensional facet appears in Equation~\ref{eq2} exactly two times.
Since $C_{2p-2}(\Delta)$ is a $K$-vector space with basis given by the
$(2p-2)$-dimensional facets, the linear independence of this basis implies that
the sum of the coefficients corresponding to each such facet must vanish.\\
Fix $h_1,h_2\in [i+2]$ with $h_1<h_2$. For all $k\in [2p+2]\setminus[i+2]$, consider the coefficients of $[1,\dots,\widehat{h_1}, \dots,\widehat{h_2},\dots, i+2,\,i+3,\dots,\widehat{k},\dots,2p+2]$ in Equation \eqref{eq2}, whose sum is zero from what we said above, that is,
$$
(-1)^{h_1+1}\lambda_{h_2,k}+(-1)^{h_2}\lambda_{h_1,k}=0\quad \Rightarrow \quad\frac{\lambda_{h_1,k}}{(-1)^{h_1}}=\frac{\lambda_{h_2,k}}{(-1)^{h_2}}
$$
This means that, for all $h\in [i+2]$ and $k\in [2p+2]\setminus[i+2]$ there exists $C_k\in K\setminus\{0\}$ such that 
\begin{equation}\label{eq lambda1}
    \lambda_{h,k}=(-1)^hC_k
\end{equation}
By a similar argument, there exists $D_h\in K\setminus\{0\}$ such that \begin{equation}\label{eq lambda2}
    \lambda_{h,k}=(-1)^kD_h.
\end{equation}

By the Equations \ref{eq lambda1} and \ref{eq lambda2}, we get the there exists $c\in K\setminus\{0\}$ such that 
$$
\frac{C_k}{(-1)^k}=\frac{D_h}{(-1)^h}=c\quad \Rightarrow\quad \begin{cases}
    C_k= c (-1)^k\\
    D_h = c (-1)^h
\end{cases}
$$

Consequently, we obtain the relations $\lambda_{h,k} = (-1)^{h+k} c$ for all $h\in [i+2]$ and $k\in [2p+2]\setminus[i+2]$, that is, Claim 2 is proved. $\square$\\



From Claim 2 it follows that $\Lambda= c \Gamma$, so $\Lambda\in \mathrm{span}_K\{\Gamma\}$. Therefore, $\tilde{H}_{2p-1}(\Delta^{[p]})=\mathrm{span}_K\{\Gamma\}$, so $\beta_{1,2p+2}(I_{\Delta}^{[p]}) = \dim_K \tilde{H}_{2p-1}(\Delta^{[p]})=1.$\\


\textbf{Case 2)} Suppose that $\Delta$ properly contains an induced subgraph isomorphic to
$K_{2+i,\,2p-i}$ for some even integer $i$ with $0 \leq i \leq p$. By Hochster’s formula, we have
\[
\beta_{1,2p+2}(I_{\Delta}^{[p]}) =
\sum_{\substack{W \subseteq [n] \\ |W| = 2p+2}}
\dim_K \widetilde{H}_{2p-1}(\Delta_W^{[p]}; K).
\]
Let $W$ be a subset of $2p+2$ vertices of $\Delta$.
If $\Delta_W= K_{2+i,\,2p-i}$ for some even integer $i$ with $0 \leq i \leq p$,
then $\dim_K \widetilde{H}_{2p-1}(\Delta_W^{[p]}; K)=1$ by Case~(1).
If $\Delta_W$ is not of the form $K_{2+i,\,2p-i}$, then
$\dim_K \widetilde{H}_{2p-1}(\Delta_W^{[p]}; K)=0$; otherwise,
$I_{\Delta_W}^{[p]}$ would fail to be linearly related, contradicting
Corollary~\ref{coro: lin related}.
This concludes the proof of the theorem.
\end{proof}

%\section{Linearity of squarefree powers of Stanley-Reisner ideal of $1$-dimensional flag simplicial complex}\label{sec: linearity simpl complex}

\section{Beyond the linearity of the first syzygy module of squarefree powers of Stanley-Reisner ideal of $1$-dimensional flag simplicial complexes}\label{sec: linearity simpl complex}

In this section, we investigate and characterize the linearity of the graded free resolutions of the squarefree powers of the Stanley--Reisner ideal associated with a $1$-dimensional flag simplicial complex. The precise statement is given in the following theorem.

 

\begin{Theorem}\label{Thm: characterization linear resolution} Let $\Delta$ be a $1$-dimensional flag simplicial complex on $[n]$, and let $2 \leq p \leq \nu(\overline{\Delta})$. The following conditions are equivalent: 
\begin{enumerate} 
\item $I_{\Delta}^{[p]}$ has linear resolution. 
\item $\Delta$ does not contain an induced subgraph isomorphic to the crown graph $Cr(2p+1)$ or to $K_{2+i,\,2p-i}$ for some even integer $i$ with $0 \leq i \leq p$. 
\end{enumerate} 
\end{Theorem}

%As condition~(2) in the above theorem shows, there are two distinct types of obstructions: one affects the property of being linearly related, while we will see that the other affects linearity at higher syzygy levels. 

%We analyze the case in which squarefree powers of the Stanley--Reisner ideal associated with a $1$-dimensional flag simplicial complex can be linearly related but fail to have a linear resolution. 
We now develop the intermediate results needed to establish the proof of Theorem~\ref{Thm: characterization linear resolution}.

\begin{Proposition}\label{Prop: Crown implies no linear}
Let $\Delta$ be a $1$-dimensional flag simplicial complex on $[n]$, and let
$2 \leq p \leq \nu(\overline{\Delta})$.
If $\Delta$ contains an induced subgraph isomorphic to the crown graph $Cr(2p+1)$, then the ideal $I_{\Delta}^{[p]}$ does not have a linear resolution.
\end{Proposition}

\begin{proof}
For convenience, set $V(Cr(2p+1))=\{1, \dots, 2p+1\}\sqcup\{1', \dots, (2p+1)'\}$ and 
$E(Cr(2p+1))=\{\{i,j'\}: i,j\in [2p+1],\ i\neq j\}$. By Hochster’s formula, we have
\[
\beta_{2p+1,\,4p+2}\!\left(I_{\Delta}^{[p]}\right) =
\sum_{\substack{W \subseteq [n] \\ |W| = 4p+2}}
\dim_K \widetilde{H}_{2p-1}(\Delta_W^{[p]}; K).
\]
Let $W$ be the set of vertices of $Cr(2p+1)$. We show that
\[
\widetilde{H}_{2p-1}(\Delta_W^{[p]})
= \ker(\delta_{2p-1}) / \mathrm{Im}(\delta_{2p}) \neq (0).
\]
The proof follows a similar strategy to that used in the proof of Theorem~\ref{Thm: betti 2,2p+2} -- Case~1) (Existence). Indeed, it suffices to show that there exists a non-zero $(2p-1)$-dimensional cycle $\Gamma$ in $\Delta_W^{[p]}$, up to multiplication by scalars in $K$.
%To describe the construction of this cycle, we first fix an odd integer $i \in \{1, \dots, 2p\}$. Intuitively, the cycle will be defined as the alternating sum of all facets in $\mathcal{F}(\Delta, K_{i, 2p-i})$ with appropriately chosen signs.

We now describe the construction of the cycle $\Gamma$. We first fix an odd integer $i \in \{1, \dots, 2p\}$. %To explain the procedure of constructing $\Gamma$, we provide a running example for clarity. As a running example, we consider the case $p=2$ and the crown graph $Cr(5)$ with $V(Cr(5))=\{1,2,3,4,5\}\sqcup\{1',2',3',4',5'\}.$ In this case, $i=1$, and the $3$-dimensional facets of $\Delta_W^{[2]}$ are precisely the elements of $\mathcal{F}(\Delta,K_{1,3})$.



\begin{enumerate}
\item Let $H=\{h_1,\ldots,h_{2p+1-i}\}\subseteq [2p+1],$ with $h_1<\cdots<h_{2p+1-i}.$ Denote by $H^c=[2p+1]\setminus H=\{a_1<\cdots<a_i\}$ the complement of $H$ in $[2p+1]$. We define
\[
F(H):=[a_1,\ldots,a_i,h_1',\ldots,h_{2p+1-i}'].
\]
%Thus, $F(H)$ consists of the non-primed vertices not contained in $H$, together with the primed copies of the vertices in $H$.

% In our running example, there are $5$ ways to choose $H\subseteq [5]$ with $|H|=4$. The possible choices are $\{2,3,4,5\},\;
% \{1,3,4,5\},
% \{1,2,4,5\},
% \{1,2,3,5\},
% \{1,2,3,4\}.$
% For instance, if $H=\{2,3,4,5\},$ then $H^c=\{1\}$ and hence $F(H)=[1,2',3',4',5'].$ Similarly,
% \[
% \begin{aligned}
% F(\{1,3,4,5\})&=[2,1',3',4',5'], \qquad
% F(\{1,2,4,5\})=[3,1',2',4',5'],\\
% F(\{1,2,3,5\})&=[4,1',2',3',5'], \qquad
% F(\{1,2,3,4\})=[5,1',2',3',4'].
% \end{aligned}
% \]

\item We define $\partial'(F(H))$ by applying a differential only to the primed part, that is, we set
\[
\partial'(F(H))
=
\sum_{r=1}^{2p+1-i}
(-1)^{r+1}
[a_1,\ldots,a_i,h_1',\ldots,\widehat{h_r'},\ldots,h_{2p+1-i}'].
\]
% In our running example, we obtain
% $$
% \begin{aligned}
% \partial'(F(\{2,3,4,5\}))
% &=
% [1,3',4',5']
% -[1,2',4',5']
% +[1,2',3',5']
% -[1,2',3',4'],
% \\
% \partial'(F(\{1,3,4,5\}))
% &=
% [2,3',4',5']
% -[2,1',4',5']
% +[2,1',3',5']
% -[2,1',3',4'],
% \\
% \partial'(F(\{1,2,4,5\}))
% &=
% [3,2',4',5']
% -[3,1',4',5'] 
% +[3,1',2',5']
% -[3,1',2',4'],
% \\
% \partial'(F(\{1,2,3,5\}))
% &=
% [4,2',3',5']
% -[4,1',3',5']
% +[4,1',2',5']
% -[4,1',2',3'],
% \\
% \partial'(F(\{1,2,3,4\}))
% &=
% [5,2',3',4']
% -[5,1',3',4']
% +[5,1',2',4']
% -[5,1',2',3'].
% \end{aligned}
% $$

\item Finally, we define
$$
\Gamma
=
\sum_{\substack{H\subseteq [2p+1]\\ |H|=2p+1-i}}
(-1)^{\sum_{h\in H} h}
\,\partial'(F(H)).
$$
\end{enumerate}

We now show that $\delta_{2p-1}(\Gamma)=0$. In order to do that, it can be useful to regard the the terms in $\Gamma$ and $\delta_{2p-1}(\Gamma)$ as follows.

Every vector $F(H)$ can be combinatorially represented by a $3 \times (2p+1)$ table. The entries in the first row are the integers $\{1,\dots,2p+1\}$ arranged in increasing order from left to right. The second row corresponds to the set $\{a_1,\dots,a_i\}$, namely, the $j$-th column is marked with a bullet whenever $j\in \{a_1,\dots,a_i\}$. Similarly, the third row corresponds to the set $\{h_1,\dots,h_{2p+1-i}\}$. Hence, every facet of $\Gamma$ can be represented by such a table after removing one point from the bottom row. Consequently, every term in $\delta_{2p-1}(\Gamma)$ can be represented by such a table after removing either one point from the second row and one point from the bottom row, or two points from the bottom row. For example, for $p=5$ and $i=7$, Figure~\ref{fig: for tables1}(A) represents
$F(H)=[2,5,9,1',3',4',6',7',8',10',11']$.
In Figure~\ref{fig: for tables1}(B), where $h_r'=7'$ with $r=5$, we obtain the corresponding facet of $\Gamma$, that is,
$[2,5,9,1',3',4',6',\widehat{7'},8',10',11']$. In Figure \ref{fig: for tables2} we have the representations of two terms in $\delta_{2p-1}(\Gamma)$, namely $[2,\widehat{5},9,1',3',4',\widehat{6'},7',8',10',11']$ and $[2,5,9,1',3',4',6',\widehat{7'},\widehat{8'},10',11']$.

\begin{figure}[h]
\centering
\subfloat[]{\includegraphics[scale=0.075]{table1.eps}}\
\subfloat[]{\includegraphics[scale=0.075]{table4.eps}}
\caption{Example for $p=5$ and $i=7$.}
\label{fig: for tables1}
\end{figure}


\begin{figure}[h]
\centering
\subfloat[]{\includegraphics[scale=0.075]{table2.eps}}\
\subfloat[]{\includegraphics[scale=0.075]{table3.eps}}
\caption{Example for $p=5$ and $i=7$: terms in $\delta_{2p-1}(\Gamma)$.}
\label{fig: for tables2}
\end{figure}

% In our running example, we obtain
% $$
% \begin{aligned}
% \Gamma
% ={}&
% \partial'(F(\{2,3,4,5\}))
% -\partial'(F(\{1,3,4,5\})) 
% +\partial'(F(\{1,2,4,5\}))
% -\partial'(F(\{1,2,3,5\}))
% +\partial'(F(\{1,2,3,4\})).
% \end{aligned}
% $$

We can now compute  $\delta_{2p-1}(\Gamma)$. Observe first that, in the expression of $\delta_{2p-1}(\Gamma)$, the sum of the terms obtained by deleting two primed vertices is zero because it is the result of the application of
two simplicial differentials. Therefore, it suffices to consider the terms obtained by deleting exactly one non-primed vertex and one primed vertex, that is, in other words, the $3 \times (2p+1)$ tables with one missing point in the second row and another in the bottom row.

Fix $H=\{h_1,\ldots,h_{2p+1-i}\}\subseteq [2p+1]$, with $h_1<\cdots<h_{2p+1-i}$, and write $H^c=\{a_1<\cdots<a_i\}$. Choose $h_r\in H$ and $a_s\in H^c$. Consider 
\[
\gamma_1:=
[a_1,\ldots,\widehat{a_s},\ldots,a_i,
h_1',\ldots,\widehat{h_r'},\ldots,h_{2p+1-i}'],
\]
The occurrence given by $\gamma_1$ can occur only in an another alternative way in $\delta_{2p-1}(\Gamma)$, that we now construct.

Set $\widetilde H=(H\setminus\{h_r\})\cup\{a_s\},$ so
$\widetilde H^c=(H^c\setminus\{a_s\})\cup\{h_r\}$ (see Figure \ref{fig:table last}). Writing
$\widetilde H=\{w_1<\cdots<w_{2p+1-i}\}$ and
$\widetilde H^c=\{c_1<\cdots<c_i\},$ consider $F(\widetilde H)=[c_1,\ldots,h_r,\ldots,c_i,
w_1',\ldots,a_s',\ldots,w_{2p+1-i}']$ and delete $a_s'$ and $h_r$ from $F(\widetilde H)$, getting
\[
\gamma_2:=
[c_1,\ldots,\widehat{h_r},\ldots,c_i,
w_1',\ldots,\widehat{a_s'},\ldots,w_{2p+1-i}'].
\]
By construction, $\gamma_1=\gamma_2$. Moreover, it is easy to observe that this is the only way to obtain the same vector in $\delta_{2p-1}(\Gamma)$. 

\begin{figure}[h]
    \centering
    \includegraphics[scale=0.1]{table5.eps}
    \caption{Example of $F(\widetilde H)$ and $\gamma_2$: $a_s$ moves to the bottom row and is removed by (2) in the construction of $\Gamma$, while the entry $h_r$ is lifted to the second row and is removed by applying $\delta_{2p-1}$}
    \label{fig:table last}
\end{figure}


We now study the signs of $\gamma_1$ and $\gamma_2$ in the expression of $\delta_{2p-1}(\Gamma)$. Assume first that $a_s<h_r$ (see Figure~\ref{fig:table last}). The sign of $\gamma_1$, denoted by $\operatorname{sign}(\gamma_1)$, is given by $(-1)^{e(\gamma_1)}$, where $e(\gamma_1)$ is the sum of three contributions: two arising from (2) and (3) in the construction of $\Gamma$, and the third coming from the application of $\delta_{2p-1}$ to $\Gamma$. More precisely,
\[
e(\gamma_1)=\underbrace{\sum_{j=1}^{2p+1-i} h_j}_{\text{From (3) of definition of $\Gamma$}}
+
\underbrace{(r+1)}_{\text{From (2) of definition of $\Gamma$}}
+
\underbrace{(s+1)}_{\text{From $\delta_{2p-1}(\gamma_1)$}}
\]


The sign of $\gamma_2$ is given by $(-1)^{e(\gamma_2)}$, where $e(\gamma_2)$ also consists of three contributions, that we express in terms of $a_s$, $h_r$, $s$, and $r$ in order to compare them later with $e(\gamma_1)$. \\
Since $a_s<h_r$, using the combinatorial representation of $\gamma_2$, as shown in Figure~\ref{fig:table last}, where $a_s'$ and $h_r$ in $\gamma_2$ play the role of $h_r'$ and $a_s$ in $\gamma_1$, respectively, we obtain:
\[
e(\gamma_2)= 
\underbrace{\left(\sum_{\substack{j=1\\ j\neq t}}^{2p+1-i} h_j+a_s\right)}_{\text{From (3) of definition of $\Gamma$}}+\underbrace{\Big(\big(\vert\{h\in H\mid h<a_s\}\vert+1\big)+1\Big)}_{\text{From (2) of definition of $\Gamma$}}+\underbrace{\Big(\vert\{a\in H^c\mid a<h_r\}\vert+1\Big)}_{\text{From $\delta_{2p-1}(\gamma_2)$}}
\]
Since $\vert\{h\in H\mid h<a_s\}\vert=a_s-s$ and $\vert\{a\in H^c\mid a<h_r\}\vert=h_r-r$, we have:
$$e(\gamma_2)=\left(\sum_{\substack{j=1\\ j\neq r}}^{2p+1-i} h_j+a_s\right)
+(a_s-s+2)+(h_r-r+1).
$$
Therefore,
\[
\frac{\operatorname{sign}(\gamma_1)}
{\operatorname{sign}(\gamma_2)}
=
(-1)^{2a_s-2r-2s+1}
=
-1\quad\Rightarrow\quad \operatorname{sign}(\gamma_1)=-\operatorname{sign}(\gamma_2)
\]

Hence $\gamma_1$ and $\gamma_2$ cancel pairwise in $\delta_{2p-1}(\Gamma)$. The case $a_s>h_r$ can be treated analogously. Thus $\Gamma\in \ker(\delta_{2p-1})$, and so $\ker(\delta_{2p-1})\neq (0)$. Consequently,
\[
\widetilde H_{2p-1}(\Delta_W^{[p]};K)
=
\ker(\delta_{2p-1})/\operatorname{Im}(\delta_{2p})
\neq (0),
\]
which implies that $\beta_{2p+1,4p+2}(I_\Delta^{[p]})\neq 0$.
\end{proof}





%%%%%%%%%%%%%%%%%%%%%%%%%%%%%%%%%%%%%%%%
%\begin{enumerate}   
% \item First, we select $2p+1-i$ non-primed integers to be removed from 
% $[1,\dots,2p+1,1',\dots,(2p+1)']$. Let
% \[
% \{h_1<\dots<h_{2p+1-i}\} \subset \{1,\dots,2p+1\}.
% \]

% In our example, there are $5$ ways to choose $4$ non-primed integers to be removed from $[1,2,3,4,5,1',2',3',4',5']$, namely 
% \[
% \{2,3,4,5\},\ 
% \{1,3,4,5\},\ 
% \{1,2,4,5\},\ 
% \{1,2,3,5\},\ 
% \{1,2,3,4\}.
% \]

% \item Next, we choose $i$ primed integers as the complement of the removed non-primed integers, that is,
% \[
% \{k_1',\dots,k_i'\} = \{1',\dots,(2p+1)'\} \setminus \{h_1',\dots,h_{2p+1-i}'\}.
% \]

% In our example, this yields the correspondence
% \[
% \{2,3,4,5\}\rightarrow \{1'\},\quad
% \{1,3,4,5\}\rightarrow \{2'\},\quad
% \{1,2,4,5\}\rightarrow \{3'\},\quad
% \{1,2,3,5\}\rightarrow \{4'\},\quad
% \{1,2,3,4\}\rightarrow \{5'\}.
% \]

% \item We then define
% \[
% F(h_1,\dots,h_{2p+1-i}) =
% [1,\dots,\widehat{h_1},\dots,\widehat{h_{2p+1-i}},\dots,2p+1,
% 1',\dots,\widehat{k_1'},\dots,\widehat{k_i'},\dots,(2p+1)'] .
% \]

% In our case we have:
% \[
% \begin{aligned}
% F(\{2,3,4,5\})&=[1,\widehat{2},\widehat{3},\widehat{4},\widehat{5},\widehat{1'},2',3',4',5']
% =[1,2',3',4',5'],\\
% F(\{1,3,4,5\})&=[\widehat{1},2,\widehat{3},\widehat{4},\widehat{5},1',\widehat{2'},3',4',5']
% =[2,1',3',4',5'],\\
% F(\{1,2,4,5\})&=[\widehat{1},\widehat{2},3,\widehat{4},\widehat{5},1',2',\widehat{3'},4',5']
% =[3,1',2',4',5'],\\
% F(\{1,2,3,5\})&=[\widehat{1},\widehat{2},\widehat{3},4,\widehat{5},1',2',3',\widehat{4'},5']
% =[4,1',2',3',5'],\\
% F(\{1,2,3,4\})&=[\widehat{1},\widehat{2},\widehat{3},\widehat{4},5,1',2',3',4',\widehat{5'}]
% =[5,1',2',3',4'].
% \end{aligned}
% \]

% \item Finally, we apply a simplicial differential only to the primed part and set
% \[
% \partial'(F(h_1,\dots,h_{2p+1-i})) =
% \sum_{\ell=1}^{2p-i} (-1)^{\ell+1}
% [1,\dots,\widehat{h_1},\dots,\widehat{h_{2p+1-i}},\dots,2p+1,
% 1',\dots,\widehat{a_\ell'},\dots,\widehat{k_1'},\dots,\widehat{k_i'},\dots,(2p+1)'] .
% \]

% In our example we have:
% \begin{align*}
% &\partial'(F(\{2,3,4,5\}))=\partial'([1,2',3',4',5']) = [1,3',4',5']
% -[1,2',4',5']
% +[1,2',3',5']
% -[1,2',3',4'], \\
% &\partial'(F(\{1,3,4,5\}))=  \partial'([2,1',3',4',5'])= [2,3',4',5']
% -[2,1',4',5']
% +[2,1',3',5']
% -[2,1',3',4'], \\ 
% &\partial'(F(\{1,2,4,5\}))= \partial'([3,1',2',4',5']) = [3,2',4',5']
% -[3,1',4',5']
% +[3,1',2',5']
% -[3,1',2',4'], \\
% &\partial'(F(\{1,2,3,5\}))= \partial'([4,1',2',3',5']) = [4,2',3',5']
% -[4,1',3',5']
% +[4,1',2',5']
% -[4,1',2',3'], \\
% &\partial'(F(\{1,2,3,4\}))= \partial'([5,1',2',3',4']) = [5,2',3',4']
% -[5,1',3',4']
% +[5,1',2',4']
% -[5,1',2',3'] .
% \end{align*}

% \item The cycle $\Gamma$ is obtained by summing over all choices of non–primed integers:
% \[
% \Gamma =
% \sum_{h_1<\dots<h_{2p+1-i}}
% (-1)^{\sum_{j=1}^{2p+1-i} h_j}
% \,\partial'\!\left(F(h_1,\dots,h_{2p+1-i})\right).
% \]

% Concluding our example, $\Gamma$ will have the following form:
% \[
% \Gamma
% =\;\partial'(F(\{2,3,4,5\}))
% -\partial'(F(\{1,3,4,5\}))
% +\partial'(F(\{1,2,4,5\}))
% -\partial'(F(\{1,2,3,5\}))
% +\partial'(F(\{1,2,3,4\})).
% \]
% \end{enumerate}

% We need to show now that $\delta_{2p-1}(\Gamma)=0$. Let us start by observing that, in the expression of $\delta_{2p-1}(\Gamma)$, the sum of the terms of the type 
% $$
% [1,\dots,\widehat{h_1},\dots,\widehat{h_{2p+1-i}},\dots,2p+1,
% 1',\dots,\widehat{a_\ell'},\dots,\widehat{b_\epsilon},\dots,\widehat{k_1'},\dots,\widehat{k_i'},\dots,(2p+1)']
% $$
% where $b_\epsilon$ runs over $[(2p+1)']\setminus \{a_\ell', k_1',\dots,k_i'\}$, is zero because it is the result of the application of two simplicial differentials.
% Therefore, in the expression of $\delta_{2p-1}(\Gamma)$, it suffices to study only the terms of the type:
% $$
% [1,\dots,\widehat{b_\epsilon},\dots,\widehat{h_1},\dots,\widehat{h_{2p+1-i}},\dots,2p+1,
% 1',\dots,\widehat{a_\ell'},\dots,\widehat{k_1'},\dots,\widehat{k_i'},\dots,(2p+1)']
% $$
% where $b_\epsilon$ runs over $[2p+1]\setminus \{h_1,\dots,h_{2p+1-i}\}$. %In such a case the philosophy of the proof is same as in the proof of Theorem~\ref{Thm: betti 2,2p+2} -- Case~1) (Existence).
% We now show how to proceed.

% Fix $h_1,\dots, h_{2p+1-i}$, then $k_1',\dots,k_i'$, and consider $\ell \in [2p-i]$. We now fix $b_\epsilon$ in $[2p+1]\setminus \{h_1,\dots,h_{2p+1-i}\}$. Consider the vector 
% $$
% \gamma_1:=[1,\dots,\widehat{b_\epsilon},\dots,\widehat{h_1},\dots,\widehat{h_{2p+1-i}},\dots,2p+1,
% 1',\dots,\widehat{a_\ell'},\dots,\widehat{k_1'},\dots,\widehat{k_i'},\dots,(2p+1)'].
% $$

% The occurrence given by $\gamma_1$ can occur only in an another alternative way, that now construct. Since $a_\ell = h_t$ for some $t$ and $b_\epsilon = k_r$ for some $r$, the only new occurrence that equals $\gamma_2$ is given by the following: %we then construct a two new sequences by taking $\{h_1,\dots,h_{2p+1-i}\}$ removing $h_t$ and inserting $k_r$, and analogously, another one by removing $k_r'$ from $\{k_1,\dots,k_i\}$ and inserting $h_t'$. Accordingly, we replace $b_\epsilon$ and $a_\ell$ with $h_t$ and $k_r$, respectively. Formally:
% \[
% \gamma_2:=[1,\dots,\widehat{w_1},\dots,\widehat{c_\nu},\dots,\widehat{w_{2p+1-i}},\dots,2p+1,
% 1',\dots,\widehat{g_\mu'},\dots,\widehat{z_1'},\dots,\widehat{z_i'},\dots,(2p+1)'].
% \]
% with
% \begin{itemize}
%     \item $\{w_1,\dots,w_{2p+1-i}\}
% = (\{h_1,\dots,h_{2p+1-i}\}\setminus\{a_\ell\})\cup\{b_\epsilon\},$
% \item $c_\nu=a_\ell$,
% \item $\{z_1',\dots,z_i'\}
% = (\{k_1',\dots,k_i'\}\setminus\{b_\epsilon'\})\cup\{a_\ell'\},$
% \item $g_\mu' = b_{\epsilon}'$.
% \end{itemize}

% We now study the signs of $\gamma_1$ and $\gamma_2$. To this end, it is convenient to arrange the $2p+1$ integers representing $F(\{h_1,\dots,h_{2p+1-i}\})$ into a $2\times (2p+1)$ table, named $\cT_{p,i}$: in the first row we write $\{\widehat{h_1},\dots,\widehat{h_{2p+1-i}}, k_1,\dots,k_i\}$ in increasing order while in the second row we place $\{h_1,\dots,h_{2p+1-i}, \widehat{k_1},\dots,\widehat{k_i}\}$ (see Table \ref{tab:sign-table}).

% \begin{table}[h]
% \centering
% \[\cT_{5,7}=
% \renewcommand{\arraystretch}{1.6}
% \begin{array}{c|c|c|c|c|c|c|c|c|c|c}
% k_1 & \widehat{h_1} & k_2 & \widehat{h_2} & k_3 & k_4 & k_5 & k_6 & \widehat{h_3} & k_7 & k_8 \\ \hline
% \widehat{k_1} & h_1 & \widehat{k_2} & h_2 & \widehat{k_3} & \widehat{k_4} & \widehat{k_5} & \widehat{k_6} & h_3 & \widehat{k_7} & \widehat{k_8}
% \end{array}
% \]
% \caption{Example for $p=5$ and $i=7$.}
% \label{tab:sign-table}
% \end{table}

% Assume that $a_\ell<b_\epsilon$. Then the sign of $\gamma_1$, denoted by $\mathrm{sign}(\gamma_1)$, is given by $(-1)$ power to
% \[
% \left(\sum_{j=1}^{2p+1-i} h_j\right)+(r+1)+(t+1).
% \]

% Reading the table $\cT_{p,i}$ where we switch the position of $a_\ell$ and $b_\epsilon$, the sign of $\gamma_2$ is $(-1)$ power to 
% \begin{align*}
% &
% \left(\sum_{\substack{j=1\\ j\neq t}}^{2p+1-i} h_j+k_r\right)
% +
% \Big(\big(\#\{k_s\in\{k_1,\dots,k_i\}\mid k_s<h_t\}+1\big)+1\Big)
% +\\[2mm]
% &+\Big(\#\{h_s\in\{h_1,\dots,h_{2p+1-i}\}\mid h_s<k_r\}+1\Big)=
% \left(\sum_{\substack{j=1\\ j\neq t}}^{2p+1-i} h_j+k_r\right)
% +
% (h_t-t+2)
% +
% (k_r-r+1)
% .
% \end{align*}

% Therefore,
% \[
% \frac{\mathrm{sign}(\gamma_1)}{\mathrm{sign}(\gamma_2)}
% =
% (-1)^{2k_r-2r-2t+1}
% =
% -1.
% \]

% Hence, $\gamma_1$ and $\gamma_2$ cancel pairwise in $\delta_{2p-i}(\Gamma)$. The case $a_\ell>b_\epsilon$ can be treated analogously.

% In conclusion we have that $\Gamma \in \ker(\delta_{2p-1})$, that is, $\ker(\delta_{2p-1})\neq (0)$, and therefore
% \[
% \widetilde{H}_{2p-1}(\Delta_W^{[p]})
% = \ker(\delta_{2p-1}) / \mathrm{Im}(\delta_{2p}) \neq (0),
% \]
% which implies that $\beta_{2p+1,4p+2}(I_{\Delta}^{[p]})\neq 0$.




\begin{Proposition}\label{Prop: other implciation}
Let $\Delta$ be a $1$-dimensional flag simplicial complex on $[n]$, and let
$2 \leq p \leq \nu(\overline{\Delta})$. Suppose that $\Delta$ does not contain an induced subgraph isomorphic to the crown graph $Cr(2p+1)$ or to $K_{2+i,\,2p-i}$ for some even integer $i$ with $0 \leq i \leq p$.  Then $I_{\Delta}^{[p]}$ has linear resolution. 
\end{Proposition}

\begin{proof}
  Suppose that $I_{\Delta}^{[p]}$ does not have a linear resolution. Then, by Corollary \ref{Coro:reg}, we have $\reg(I_{\Delta}^{[p]})=2p+1$. Hence there exists an integer $i$ such that $\beta_{i,2p+1+i}(I_{\Delta}^{[p]})\neq 0$. By Hochster's formula, there is a subset $W\subseteq [n]$ with $|W|=2p+1+i$ such that $\widetilde H_{2p-1}(\Delta_W^{[p]};K)\neq (0)$.
  
  Recall that 
\[
\tilde{H}_{2p-1}(\Delta_W^{[p]}) = \ker(\delta_{2p-1}) / \mathrm{Im}(\delta_{2p}).
\]
Since, by Proposition \ref{Thm: dimension} (1), there are no $2p$-dimensional facets, we have $C_{2p} = (0)$, and hence $\mathrm{Im}(\delta_{2p}) = (0)$. Therefore, $\ker(\delta_{2p-1}) \neq (0)$. Thus there exists a non-trivial $(2p-1)$-cycle $\Gamma$ in $\Delta_W^{[p]}$. 

We write $\Gamma$ as a formal linear combination, with coefficients in $K\setminus\{0\}$, of oriented $(2p-1)$-dimensional facets of $\Delta_W^{[p]}$. The set of facets appearing in this expression is called the support of $\Gamma$ and is denoted by $\operatorname{supp}(\Gamma)$. 

We will show that, under the assumptions of the proposition, starting from any facet in $\mathrm{supp}(\Gamma)$ forces the existence of infinitely many distinct facets in $\mathrm{supp}(\Gamma)$, a contradiction. The argument proceeds by a sequence of steps, illustrated throughout by a running example.


%Our goal is to start from a facet in $\mathrm{supp}(\Gamma)$ and show, through the steps that we describe soon, that if $\Delta$ contains no induced subgraph isomorphic to the crown graph $Cr(2p+1)$ and no induced subgraph isomorphic to $K_{2+i,\,2p-i}$ for any even integer $i$ with $0 \leq i \leq p$, then $|\mathrm{supp}(\Gamma)|$ would necessarily be infinite, yielding a contradiction. In each step, we also discuss a running example to make the argument clearer and easier to follow.

\textbf{Step 1.}
Let $F\in \operatorname{supp}(\Gamma)$. By Theorem~\ref{Thm: Facet, pure, CM}, we have $F\in \mathcal{F}(\Delta,K_{j,2p-j})$
for some odd integer $j$ with $1\leq j\leq p$. Assume that $V(\Delta_F)=\{1,\dots,j\}\sqcup\{1',\dots,(2p-j)'\}$
and $E(\Delta_F)=\bigl\{\{u,v\}:u\in [j],\ v\in[(2p-j)']\bigr\}.$
Thus, $F=[1,\dots,j,1',\dots,(2p-j)'].$ Applying the simplicial differential, we obtain
\begin{align*}
\delta_{2p-1}(F)
={}&
\sum_{h=1}^{j}
(-1)^{h+1}
[1,\dots,\widehat{h},\dots,j,1',\dots,(2p-j)']
+
\sum_{h=1}^{2p-j}
(-1)^h
[1,\dots,j,1',\dots,\widehat{h'},\dots,(2p-j)'].
\end{align*}

As a running example, let $p=2$. Then $\Delta_F$ is an induced subgraph of $\Delta$ isomorphic to $K_{1,3}$ (see Figure~\ref{fig1}), and $F=[1,1',2',3'].$ Hence
\[\delta_3(F)=
[1',2',3']
-[1,2',3']
+[1,1',3']
-[1,1',2'].\]

\begin{figure}[h]
    \centering
    \includegraphics[scale=0.5]{exa_0_final_proof.eps}
    \caption{Running example - Step 1}
    \label{fig1}
\end{figure}

\textbf{Step 2.}
Since
$
[1,\dots,\widehat{h},\dots,j,1',\dots,(2p-j)']
$
appears in $\delta_{2p-1}(F)$ and $\delta_{2p-1}(\Gamma)=0$, it must also appear in the boundary of another facet in $\operatorname{supp}(\Gamma)$. Such a facet must be obtained by adjoining a new non-primed vertex, denoted by $j+1$, such that
\begin{itemize}
    \item $\{j+1,b\}\in E(\Delta)$ for all $b\in[(2p-j)']$;
    \item $\{j+1,a\}\notin E(\Delta)$ for all $a\in[j]$, otherwise $\Delta$ would contain a triangle;
    \item $[1,\dots,\widehat{h},\dots,j,j+1,1',\dots,(2p-j)']\in\operatorname{supp}(\Gamma)$, with $1\leq h\leq j$.
\end{itemize}

Dually, the cancellation of $[1,\dots,j,1',\dots,\widehat{h'},\dots,(2p-j)']$ forces the existence of a new primed vertex $(2p-j+1)'$ satisfying
\begin{itemize}
    \item $\{a,(2p-j+1)'\}\in E(\Delta)$ for all $a\in[j]$;
    \item $\{(2p-j+1)',b\}\notin E(\Delta)$ for all $b\in[(2p-j)']$, otherwise $\Delta$ would contain a triangle;
    \item $[1,\dots,j,1',\dots,\widehat{h'},\dots,(2p-j)',(2p-j+1)']\in\operatorname{supp}(\Gamma)$, with $1\leq h\leq 2p-j$.
\end{itemize}

Now let $W'=\{1,\dots,j,j+1\}\cup\{1',\dots,(2p-j)',(2p-j+1)'\}.$ If
$\{j+1,(2p-j+1)'\}\in E(\Delta),$ then each vertex in
$\{1,\dots,j,j+1\}$ is adjacent to each vertex in
$\{1',\dots,(2p-j)',(2p-j+1)'\}.$ Hence $\Delta_{W'}$ is isomorphic to
$K_{k+2,2p-k},$ where $k=j-1$, contradicting assumption~(2). Therefore,
$\{j+1,(2p-j+1)'\}\notin E(\Delta).$

In our running example, the terms
$
[1',2',3'], [1,2',3'], [1,1',3'], [1,1',2']
$
appear in $\delta_3(\Gamma)$. Since these terms must cancel and $\Delta$ contains no triangles, there exist vertices $2$ and $4'$ such that (see Figure~\ref{fig2})
\begin{itemize}
    \item $4'$ is adjacent to $1$, but not to $1',2',3'$;
    \item $2$ is adjacent to $1',2',3'$, but not to $1$;
    \item the facets
    $
    [1,1',2',4'], [1,1',3',4'], [1,2',3',4'], [2,1',2',3']
    $
    belong to $\operatorname{supp}(\Gamma)$.
\end{itemize}
Moreover, $\{2,4'\}\notin E(\Delta)$, otherwise the induced subgraph on
$W'=\{1,2\}\cup\{1',2',3',4'\}$ would be isomorphic to $K_{2,4}$.

\begin{figure}[h]
    \centering
    \includegraphics[scale=0.5]{exa_1_final_proof.eps}
    \caption{Running example - Step 2}
    \label{fig2}
\end{figure}

\textbf{Step 3.}
By Step~2 and Theorem~\ref{Thm: Facet, pure, CM}, we have that for every $1\leq h\leq 2p-j$, the facet $[1,\dots,j,1',\dots,\widehat{h'},\dots,(2p-j)',(2p-j+1)']$ belongs to $\operatorname{supp}(\Gamma)$. Hence
$$
[1,\dots,\widehat{k},\dots,j,
1',\dots,\widehat{h'},\dots,(2p-j)',(2p-j+1)']
$$
appears in $\delta_{2p-1}(\Gamma)$ for all $1\leq k\leq j$. Since $\delta_{2p-1}(\Gamma)=0$ and $\{j+1,(2p-j+1)'\}\notin E(\Delta),$
these terms must cancel with terms coming from new facets in $\operatorname{supp}(\Gamma)$. Therefore, for every fixed $h$ with $1\leq h\leq 2p-j$, there exists a new vertex, denoted by $j+1+h$, such that
\begin{itemize}
    \item $\{j+1+h,b\}\in E(\Delta)$ for all
    $ b\in[(2p-j+1)']\setminus\{h'\};$
    
    \item $\{j+1+h,a\}\notin E(\Delta)$ for all
    $a\in[j+1+\widetilde h],$
    with $0\leq\widetilde h\leq h-1$, otherwise $\Delta$ would contain a triangle;
    
    \item
    $[1,\dots,\widehat{k},\dots,j,j+1+h, 1',\dots,\widehat{h'},\dots,(2p-j)',(2p-j+1)']
    \in\operatorname{supp}(\Gamma),$ with $1\leq k\leq j$.
\end{itemize}

Continuing our running example, the simplices $[1',2',4'], [1',3',4'], [2',3',4']$ appear in $\delta_3(\Gamma)$. Since $\{2,4'\}\notin E(\Delta),$ they must cancel with terms coming from new facets in $\operatorname{supp}(\Gamma)$. Thus there exist vertices $3,4,5$ such that (see Figure~\ref{fig3})
\begin{itemize}
    \item $3$ is adjacent to $2',3',4'$; $4$ is adjacent to $1',3',4'$; and $5$ is adjacent to $1',2',4'$;
    
    \item $3$ is not adjacent to $1,2,1'$; $4$ is not adjacent to $1,2,3,2'$; and $5$ is not adjacent to $1,2,3,4,3'$;
    
    \item $[3,2',3',4'], [4,1',2',3'], [5,1',2',4']$ belong to $\operatorname{supp}(\Gamma)$.
\end{itemize}

\begin{figure}[h]
    \centering
    \includegraphics[scale=0.5]{exa_2_final_proof.eps}
    \caption{Running example - Step 3}
    \label{fig3}
\end{figure}

\textbf{Step 4.}
By duality, for every $1\leq h\leq j$, the facet
$
[1,\dots,\widehat{h},\dots,j,j+1,1',\dots,(2p-j)']
$
belongs to $\operatorname{supp}(\Gamma)$. Hence
$$
[1,\dots,\widehat{h},\dots,j,j+1,
1',\dots,\widehat{k'},\dots,(2p-j)']
$$
appears in $\delta_{2p-1}(\Gamma)$ for all $1\leq k\leq 2p-j$.

Since $\delta_{2p-1}(\Gamma)=0$ and
$
\{j+1,(2p-j+1)'\}\notin E(\Delta),
$
these terms must cancel with terms coming from new facets in $\operatorname{supp}(\Gamma)$. Therefore, for every fixed $h$ with $1\leq h\leq j$, there exists a new primed vertex, denoted by $(2p-j+1+h)'$, such that
\begin{itemize}
    \item $\{a,(2p-j+1+h)'\}\in E(\Delta)$ for all
    $
    a\in[j+1]\setminus\{h\};
    $
    
    \item $\{(2p-j+1+h)',b\}\notin E(\Delta)$ for all
    $
    b\in[2p-j+1+\widetilde h],
    $
    with $0\leq\widetilde h\leq h-1$, otherwise $\Delta$ would contain a triangle;
    
    \item
    $
    [1,\dots,\widehat{h},\dots,j,j+1,
    1',\dots,\widehat{k'},\dots,(2p-j)',
    (2p-j+1+h)']
    \in\operatorname{supp}(\Gamma),
    $
    with $1\leq k\leq 2p-j$.
\end{itemize}

Now let
$$
W'=\{1,\dots,j,j+1,\dots,2p+1\}
\cup
\{1',\dots,(2p-j)',(2p-j+1)',\dots,(2p+1)'\}.
$$
If
$
\{j+1+h,(2p-j+1+k)'\}\in E(\Delta)
$
for all $1\leq h\leq 2p-j$ and $1\leq k\leq j$, then the construction above shows that $\Delta_{W'}$ is isomorphic to $Cr(2p+1)$, contradicting assumption~(2). Therefore, there exist $h\in[2p-j]$ and $k\in[j]$ such that $\{j+1+h,(2p-j+1+k)'\}\notin E(\Delta).$

In our running example, $[2,1',2',3']$ is the only facet in $\operatorname{supp}(\Gamma)$ of the form $[1,\dots,\widehat{h},\dots,j,j+1,1',\dots,(2p-j)'].$
Hence $[2,2',3'], [2,1',3'], [2,1',2']$ appear in $\delta_3(\Gamma)$. Since
$\{2,4'\}\notin E(\Delta),$ these terms must cancel with terms coming from new facets in $\operatorname{supp}(\Gamma)$. Thus there exists a vertex $5'$ such that (see Figure~\ref{fig4})
\begin{itemize}
    \item $5'$ is adjacent to $2$;
    \item $5'$ is not adjacent to $1,1',2',3',4'$;
    \item the facets
    $
    [2,2',3',5'], [2,1',3',5'], [2,1',2',5']
    $
    belong to $\operatorname{supp}(\Gamma)$.
\end{itemize}
Moreover, if $\{3,5'\}, \{4,5'\}, \{5,5'\}\in E(\Delta),$ then the induced subgraph on $W'=\{1,2,3,4,5\}\cup\{1',2',3',4',5'\}$ is isomorphic to $Cr(5)$, contradicting assumption~(2).  Therefore, there exist $h\in[2p-j]$ and $k\in[j]$ such that $\{j+1+h,(2p-j+1+k)'\}$ is not an edge of $\Delta$. Observe that
$
[1,\dots,\widehat{h},\dots,j,j+1,1',\dots,\widehat{k'},\dots,(2p-j)',(2p-j+1+h)']\in \mathrm{supp}(\Gamma).
$
\begin{figure}[h]
    \centering
    \includegraphics[scale=0.5]{exa_3_final_proof.eps}
    \caption{Running example - Step 5}
    \label{fig4}
\end{figure}

% {\color{red}

% \textbf{(Final) Step 5.}
% By the previous step,
% $$
% F_1:=
% [1,\dots,\widehat{k},\dots,j,j+1,
% 1',\dots,\widehat{h'},\dots,(2p-j)',
% (2p-j+1+h)']
% \in\operatorname{supp}(\Gamma).
% $$

% Moreover, by construction,
% $
% F_1\in\mathcal{F}(\Delta,K_{j,2p-j}),
% $
% so $F_1$ has the same combinatorial structure as the original facet $F$. Therefore, the argument developed in the previous steps may be iterated with $F_1$ in place of $F$.
% \begin{figure}[h]
%     \centering
%     \includegraphics[scale=0.5]{exa_4_final_proof.eps}
%     \caption{Running example - Step 7}
%     \label{fig5}
% \end{figure}

% In our running example, assume that
% $\{5,5'\}\notin E(\Delta)$, while $\{3,5'\}$ and $\{4,5'\}$ are edges, although this is not essential (see Figure~\ref{fig5}). Then
% $
% F_1=[2,1',2',5']
% \in\operatorname{supp}(\Gamma).
% $
% Repeating Step~2 for $F_1$ forces the existence of a new vertex $6$ adjacent to $1',2',5'$. The corresponding primed vertex cannot coincide with any of
% $
% 1',2',3',4',5',
% $
% otherwise $\Delta$ would contain either an induced $K_{2,4}$ or an induced $Cr(5)$. Hence a new vertex $6'$ must also be introduced.

% Repeating this procedure inductively produces infinitely many distinct facets in $\operatorname{supp}(\Gamma)$. This is impossible, since $\Delta^{[p]}$ has only finitely many facets. Therefore, our assumption was false, and $I_\Delta^{[p]}$ has a linear resolution.

% }
\textbf{(Final) Step 5.}  Observe that
$
[1,\dots,\widehat{h},\dots,j,j+1,1',\dots,\widehat{k'},\dots,(2p-j)',(2p-j+1+h)']\in \mathrm{supp}(\Gamma).
$
Moreover, from the construction developed so far, neither $\{j+1+h,(2p-j+1+k)'\}$ nor $\{h,(2p-j+1+k)'\}$ is an edge of $\Delta$. Therefore, representing
$
[1,\dots,\widehat{h},\dots,j,j+1,1',\dots,\widehat{k'},\dots,(2p-j)',(2p-j+1+h)']
$
a facet in $\mathcal{F}(\Delta,K_{j,2p-j})$, that we name $F_1$, we may then iterate the arguments given in all previous Steps, with $F_1$ replacing $F$.
\begin{figure}[h]
    \centering
    \includegraphics[scale=0.5]{exa_4_final_proof.eps}
    \caption{Running example - Step 7}
    \label{fig5}
\end{figure}
In our example, with reference to Figure \ref{fig5}, assume that $\{5,5'\}$ is not an edge of $\Delta$ (while $\{3,5'\}$ and $\{4,5'\}$ are edges, although this is not essential). Then $[2,1',2',5']$ belongs to the support of $\Gamma$, and in this case we may take $F_1=\{2,1',2',5'\}$. The process must then continue by introducing new vertices. For instance, repeating Step 2 for $[1',2',5']$ yields a new vertex $6$ (playing the role of $j+1$) such that $6$ forms edges with $1',2',5'$ and does not form an edge with $3'$, since otherwise $W=\{2,6\}\cup\{1',2',3',5'\}$ would induce $K_{2,4}$ as induced subgraph, nor with $4'$, since otherwise $W=\{1,2,3,4,6\}\cup\{1',2',3',5'\}$ would induce $Cr(5)$. Consequently, the vertex playing the role of $(2p-j+1)'$ in Step 2 cannot coincide with any of the vertices in $\{1',2',3',4',5'\}$; therefore, it must be a new vertex, denoted by $6'$. Hence the procedure described in the previous Steps must be restarted, where $F_1$ will play the role of $F$.

Repeating the above argument inductively produces infinitely many facets in $\mathrm{supp}(\Gamma)$. This is impossible, since $\Delta^{[p]}$ has only finitely many facets. The contradiction shows that our assumption was false. Therefore, $I_{\Delta}^{[p]}$ has a linear resolution, as desired.
\end{proof}


We can now combine the results of this section to complete the proof of Theorem~\ref{Thm: characterization linear resolution}.

\begin{proof}[Proof of Theorem \ref{Thm: characterization linear resolution}]
(1) $\Rightarrow$ (2). Suppose that $\Delta$ contains an induced subgraph isomorphic to the crown graph $\mathrm{Cr}(2p+1)$ or to $K_{2+i,\,2p-i}$ for some even integer $i$ with $0 \leq i \leq p$. If, for some such $i$, $\Delta$ contains $K_{2+i,\,2p-i}$, then $I_{\Delta}^{[p]}$ is not linearly related by Corollary~\ref{coro: lin related}; hence, $I_{\Delta}^{[p]}$ does not have a linear resolution, which is a contradiction. If $\Delta$ contains an induced subgraph isomorphic to the crown graph $\mathrm{Cr}(2p+1)$, then $I_{\Delta}^{[p]}$ cannot have a linear resolution by Proposition~\ref{Prop: Crown implies no linear}, again yielding a contradiction. Therefore, $I_{\Delta}^{[p]}$ must have a linear resolution.\\
(2) $\Rightarrow$ (1). This follows from Proposition~\ref{Prop: other implciation}.
\end{proof}

A consequence of Theorem \ref{Thm: characterization linear resolution} is the following.

\begin{Corollary}\label{Coro: lin res for p then lin res for p+1}
Let $\Delta$ be a $1$-dimensional flag simplicial complex on $[n]$, and let
$2 \leq p \leq \nu(\overline{\Delta})$. If $I_{\Delta}^{[p]}$ has a linear resolution, then $I_{\Delta}^{[q]}$ has a linear resolution for all $p\leq q\leq \nu(\overline{\Delta})$.
\end{Corollary}

\begin{proof}
    If $I_{\Delta}^{[p]}$ has a linear resolution, then by Theorem \ref{Thm: characterization linear resolution} $\Delta$ does not contain an induced subgraph isomorphic to $K_{2+i,\,2p-i}$ for some even integer $i$ with $0 \leq i \leq p$, or to the crown graph $Cr(2p+1)$. Hence, for any integer $q$ with $p\leq q\leq \nu(\overline{\Delta})$, $\Delta$ cannot contain an induced subgraph isomorphic to $K_{2+i,\,2q-i}$ for some even integer $i$ with $0 \leq i \leq q$, or to the crown graph $Cr(2q+1)$, so $I_{\Delta}^{[q]}$ has a linear resolution, again by Theorem \ref{Thm: characterization linear resolution}.
\end{proof}


Similarly as stated in Theorem \ref{Thm: betti 2,2p+2}, we can give a combinatorial description of the Betti number $\beta_{2p+1,\,4p+2} (I_{\Delta}^{[p]})$, which provide the failure of linearity for $I_{\Delta}^{[p]}$.

\begin{Corollary}\label{coro: betti for crown}
Let $\Delta$ be a $1$-dimensional flag simplicial complex on $[n]$, and let
$2 \leq p \leq \nu(\overline{\Delta})$. Suppose that:
\begin{enumerate}
    \item $\Delta$ does not contain an induced subgraph isomorphic to $K_{2+i,\,2p-i}$ for some even integer $i$ with $0 \leq i \leq p$.
    \item $\Delta$ contain an induced subgraph isomorphic to the crown graph $Cr(2p+1)$. 
\end{enumerate} 

Then
\[
\beta_{2p+1,\,4p+2} (I_{\Delta}^{[p]}) =
p\,\bigl|\bigl\{\text{induced subgraphs of } \Delta \text{ isomorphic to } Cr(2p+1)\bigr\}\bigr|.
\]
\end{Corollary}

\begin{proof}
 The statement follows directly from the arguments used in the proofs of Propositions~\ref{Prop: Crown implies no linear} and~\ref{Prop: other implciation}. Indeed, if $W$ is a set of $4p+2$ vertices of $\Delta$, those proofs show that the only way to obtain a $(2p-1)$-dimensional cycle in $\Delta_W^{[p]}$ is when $\Delta_W$ is isomorphic to $Cr(2p+1)$, and the corresponding $(2p-1)$-dimensional cycle is precisely the complex $\Gamma$ described in the proof of Proposition~\ref{Prop: Crown implies no linear}.
\end{proof}


We conclude this section with the following example, which shows that Condition~$(1)$ in the previous corollary is essential: if it is omitted, then the statement no longer holds.

\begin{Example}\rm
Consider the $1$-dimensional flag simplicial complex $\Delta$ illustrated in Figure~\ref{fig: last exa}.

\begin{figure}[h]
    \centering
    \includegraphics[scale=0.6]{crown_graph_exa_final.eps}
    \caption{A $1$-dimensional flag simplicial complex.}
    \label{fig: last exa}
\end{figure}

Observe that $\Delta$ contains exactly one induced subgraph isomorphic to the crown graph $Cr(5)$, namely the subgraph induced by the vertex set $\{1,2,\dots,10\}$. Moreover, $\Delta$ also contains an induced subgraph isomorphic to $K_{2,4}$, for instance the one induced by the vertex set $\{6,7,11,12,13,14\}$. On the other hand, the reader can easily verify that there are exactly $35$ induced subgraphs isomorphic to $K_{2,4}$.

Using \textit{Macaulay2}~\cite{M2}, we compute the following Betti table for $I_\Delta^{[2]}$ over a field of characteristic zero, and we can observe that $\beta_{1,6}=35$ in according to Theorem \ref{Thm: betti 2,2p+2} but $\beta_{5,10}=1817\neq 2$.

\[
\begin{array}{r|ccccccccccc}
      & 0 & 1 & 2 & 3 & 4 & 5 & 6 & 7 & 8 & 9 & 10 \\ 
\hline
\mathrm{total}: & 857 & 6603 & 23844 & 52459 & 77265 & 79201 & 57024 & 28409 & 9359 & 1840 & 164 \\
0: & . & . & . & . & . & . & . & . & . & . & . \\
1: & . & . & . & . & . & . & . & . & . & . & . \\
2: & . & . & . & . & . & . & . & . & . & . & . \\
3: & . & . & . & . & . & . & . & . & . & . & . \\
4: & 857 & 6568 & 23585 & 51619 & 75705 & 77384 & 55661 & 27763 & 9181 & 1818 & 164 \\
5: & . & 35 & 259 & 840 & 1560 & 1817 & 1363 & 646 & 178 & 22 & .
\end{array}
\]
\end{Example}

%%%%%%%%%%%%%%%%%%%%%%%%%%%%%%%%%%%%%%%%%%%%%%%%%%%%%%%%%%%%%%%%


\section{Conclusions and open questions}\label{Sec: final}

In this final section, we completely determine the shape of the Betti table of square-free powers of $1$-dimensional flag simplicial complexes. We then conclude with a subsection presenting some open problems. 

We begin by recalling some preliminary definitions. Let $\Delta$ be a simplicial complex on the vertex set $[n]$. For a vertex $v \in \Delta$, we define:
\begin{itemize}
    \item the \textit{star} of $v$ in $\Delta$, denoted by $\text{star}_\Delta(v)$, by $\{G \in \Delta : \{v\} \cup G \in \Delta\}$;
    \item the \textit{link} of $v$ in $\Delta$, denoted by $\text{link}_\Delta(v)$, by $\{G \in \Delta : \{v\} \cup G \in \Delta \text{ and } v \notin G\}$;
    \item the \textit{deletion} of $v$ from $\Delta$, denoted by $\text{del}_\Delta(v)$, by  $\{G \in \Delta : v \notin G\}$.
\end{itemize}
Furthermore, note that $\text{link}_\Delta(v) = \text{star}_\Delta(v) \cap \text{del}_\Delta(v)$.

Let $\Delta_1$ and $\Delta_2$ be two simplicial complexes on disjoint vertex sets $V$ and $W$, respectively. The \textit{join} $\Delta_1 * \Delta_2$ is the simplicial complex on the vertex set $V \cup W$ whose faces are of the form $F \cup G$, where $F \in \Delta_1$ and $G \in \Delta_2$. In particular, the \textit{cone} $\cC(\Delta)$ over a simplicial complex $\Delta$ is defined as the join $\omega * \Delta$, where $\omega$ is a newly introduced apex vertex.

If $\Delta$ is a simplicial complex, then $\tilde{H}_i(\cC(\Delta)) = 0$ for all $i \ge 0$ (see \cite[Proposition 5.2.5]{Villarreal}). As an immediate consequence, if $v$ is a vertex of $\Delta$, $\text{star}_\Delta(v)$ forms a cone under the apex $v$. It follows that the reduced homology groups of the star of any vertex vanish, i.e., $\tilde{H}_i(\text{star}_\Delta(v)) = 0$ for all $i$.

Let $\Delta$ be a simplicial complex, and let $\Delta_1$ and $\Delta_2$ be two subcomplexes such that $\Delta = \Delta_1 \cup \Delta_2$. The \textit{Mayer-Vietoris sequence} is the following long exact sequence in homology:
\[ \dots \longrightarrow \tilde{H}_i(\Delta_1 \cap \Delta_2) \longrightarrow\tilde{H}_i(\Delta_1) \oplus \tilde{H}_i(\Delta_2) \longrightarrow \tilde{H}_i(\Delta) \longrightarrow \tilde{H}_{i-1}(\Delta_1 \cap \Delta_2) \longrightarrow \dots \]

This tool allows us to prove the following. 

\begin{Lemma}\label{Lemma: resol K}
Let $\Delta$ be a $1$-dimensional flag simplicial complex on the vertex set $[n]$, and let $2 \leq p \leq \nu(\overline{\Delta})$. Suppose that $\Delta$ contains an induced subgraph isomorphic to $K_{2+i,\,2p-i}$ for some even integer $i$ satisfying $0 \leq i \leq p$. Let $V(K_{2+i,\,2p-i})$ denote the vertex set of this induced subgraph. Then, for every subset of vertices $W$ of $\Delta$ such that $V(K_{2+i,\,2p-i}) \subseteq W$, we have $\widetilde{H}_{2p-1}(\Delta_W) \neq (0)$.
\end{Lemma}

\begin{proof}
If $n=2p+2$ then the claim follows immediately from Theorem \ref{Thm: betti 2,2p+2}. Assume that $n>2p+2$, so $V(K_{2+i,\,2p-i})$ is properly contained in the set of vertices of $\Delta$. Let $W$ be a set of vertices of $\Delta$ with $V(K_{2+i,\,2p-i}) \subsetneq W$ and let $v\in V(\Delta_W)\setminus V(K_{2+i,\,2p-i})$.  From the Mayor-Vietoris sequence to $\Delta_W=\text{star}_{\Delta_W}(v)\cup \text{del}_{\Delta_W}(v)$, we have:
    
\[ \dots \longrightarrow \tilde{H}_{2p-1}(\text{link}_{\Delta_W}(v)) \longrightarrow  \tilde{H}_{2p-1}(\text{del}_{\Delta_W}(v)) \longrightarrow \tilde{H}_{2p-1}(\Delta_W) \longrightarrow \tilde{H}_{2p-2}(\text{link}_{\Delta_W}(v)) \longrightarrow \dots \]

Note that $\tilde{H}_{2p-1}(\text{link}_{\Delta_W}(v))=(0)$ because $\text{dim}((\text{link}_{\Delta_W}(v))\leq 2p$. Hence the map $$\tilde{H}_{2p-1}(\text{del}_{\Delta_W}(v)) \longrightarrow \tilde{H}_{2p-1}(\Delta_W)$$ is injective. Now, set $\Delta_1= \text{del}_{\Delta_W}(v).$ Note that $\Delta_1$ contains the graphs $K_{2+i,\,2p-i}$ as induced subgraph. Hence, if $\Delta_1=K_{2+i,\,2p-i}$, we are done, otherwise we can apply again the previous argument. Repeatedly applying the previous argument we can obtain a sequence of injective maps 
$$\tilde{H}_{2p-1}(K_{2+i,\,2p-i}) \longrightarrow \dots\longrightarrow \tilde{H}_{2p-1}(\text{del}_{\Delta_W}(v)) \longrightarrow \tilde{H}_{2p-1}(\Delta_W)$$ 
Since $\tilde{H}_{2p-1}(K_{2+i,\,2p-i}) \neq (0)$, it follows that $\tilde{H}_{2p-1}(\Delta_W) \neq (0)$.
\end{proof}

\begin{Lemma}\label{Lemma: resol Crown}
Let $\Delta$ be a $1$-dimensional flag simplicial complex on the vertex set $[n]$, and let $2 \leq p \leq \nu(\overline{\Delta})$. Suppose that $\Delta$ contains an induced subgraph isomorphic to the crown graph $Cr(2p+1)$, but does not contain an induced subgraph isomorphic to $K_{2+i,\,2p-i}$ for any even integer $i$ satisfying $0 \leq i \leq p$. Let $V(Cr(2p+1))$ denote the vertex set of this crown graph. Then, for every subset of vertices $W$ of $\Delta$ such that $V(Cr(2p+1))\subseteq W$, we have $\widetilde{H}_{2p-1}(\Delta_W) \neq (0)$.
\end{Lemma} 

  \begin{proof}
      The proof is similar to that one of Lemma \ref{Lemma: resol K}.
  \end{proof}

We are now in a position to determine the vanishing and non-vanishing behaviors of the Betti numbers in the graded free resolution.

\begin{Proposition}\label{Prop: non van betti}
Let $\Delta$ be a $1$-dimensional flag simplicial complex on the vertex set $[n]$, and let $2 \leq p \leq \nu(\overline{\Delta})$. 
\begin{enumerate}
    \item If $\Delta$ contains an induced subgraph isomorphic to $K_{2+i,\,2p-i}$ for some even integer $i$ satisfying $0 \leq i \leq p$, then $\beta_{j,j+2p+1}(I_\Delta^{[p]}) \neq 0$ for all $1 \leq j \leq n-2p-1$.
    \item If $\Delta$ contains an induced subgraph isomorphic to the crown graph $Cr(2p+1)$ but does not contain an induced subgraph isomorphic to $K_{2+i,\,2p-i}$ for any even integer $i$ satisfying $0 \leq i \leq p$, then 
    \[ \beta_{j,j+2p+1}(I_\Delta^{[p]}) = 0 \quad \text{for all } 1 \leq j \leq 2p, \]
    and 
    \[ \beta_{j,j+2p+1}(I_\Delta^{[p]}) \neq 0 \quad \text{for all } 2p+1 \leq j \leq n-2p-1. \]
\end{enumerate} 
\end{Proposition}

\begin{proof}
(1) The assertion follows directly from Lemma \ref{Lemma: resol K}.

(2) If $W \subseteq [n]$ is a subset of vertices, the arguments in Proposition \ref{Prop: other implciation} ensure that a non-trivial $(2p-1)$-dimensional homology cycle in $\Delta_W^{[p]}$ can only occur when $|W| = 4p+2$ and $\Delta_W$ is isomorphic to $Cr(2p+1)$. Thus, Hochster's formula implies that $\beta_{j,j+2p+1}(I_\Delta^{[p]}) = 0$ for all $1 \leq j \leq 2p$. The non-vanishing of the remaining Betti numbers follows from Lemma \ref{Lemma: resol Crown}.
\end{proof}

In the following discussion, we summarize the main results obtained in this work concerning the minimal graded free resolution of squarefree powers of $1$-dimensional flag simplicial complexes.  

\begin{Discussion}\label{Disc. summarize}
Let $\Delta$ be a $1$-dimensional flag simplicial complex on the vertex set $[n]$, and let $2 \leq p \leq \nu(\overline{\Delta})$. 

If $\overline{\Delta}$ possesses a unique $p$-matching, then the ideal $I_{\Delta}^{[p]}$ is principal. Assuming henceforth that $\overline{\Delta}$ has more than one $p$-matching, we have the following situations cases:

\begin{enumerate}
    \item  If $\Delta \cong K_{1,n-1}$, or if $\Delta$ does not contain an induced subgraph isomorphic to $K_{i,\,2p-i}$ for any odd integer $i$ satisfying $1 \le i \le p$, then $I_{\Delta}^{[p]}$ has a linear resolution of length $n-2p-1$.
    
    \item If $\Delta \not\cong K_{1,n-1}$ and $\Delta$ contains an induced subgraph isomorphic to $K_{i,\,2p-i}$ for some odd integer $i$ satisfying $1 \le i \le p$, then three distinct situations may arise:

\begin{itemize}
    \item[2a.] If $\Delta$ contains neither an induced subgraph isomorphic to the crown graph $Cr(2p+1)$ nor an induced subgraph isomorphic to $K_{2+i,\,2p-i}$ for any even integer $i$ satisfying $0 \leq i \leq p$, then $I_{\Delta}^{[p]}$ has a linear resolution of length $n-2p$.
    
    \item[2b.] If $\Delta$ contains an induced subgraph isomorphic to $K_{2+i,\,2p-i}$ for some even integer $i$ satisfying $0 \leq i \leq p$, then the graded Betti table is as illustrated in Figure~\ref{Fig:BettiShapes1}.
    
\begin{figure}[h]
\centering
$\begin{array}{r|cccccc}
      & 0 & 1 & \cdots & n-2p-1 & n-2p \\
\hline
\text{total:} & \beta_{0} & \beta_{1} & \cdots & \beta_{n-2p-1} & \beta_{n-2p} \\
\hline
2p:   & \beta_{0, 2p} & \beta_{1, 2p+1} & \cdots & \beta_{n-2p-1, n-1} & \beta_{n-2p, n} \\
2p+1: & 0 & \beta_{1, 2p+2} & \cdots & \beta_{n-2p-1, n} & 0 \\
\end{array}$
\caption{Betti table of $I_{\Delta}^{[p]}$ in the presence of $K_{2+i, 2p-i}$.}
\label{Fig:BettiShapes1}
\end{figure}

where \[
\beta_{1,2p+2}(I_{\Delta}^{[p]})
=
\sum_{\substack{0 \leq i \leq p \\ i \text{ even}}} \big\vert\{\text{induced subgraphs of } \Delta \text{ isomorphic to } K_{2+i,\,2p-i}\}\big\vert.% \vert N_i(\Delta)\vert.
\]

    \item[2c.] If $\Delta$ does not contain an induced subgraph isomorphic to $K_{2+i,\,2p-i}$ for any even integer $i$ satisfying $0 \leq i \leq p$, but contains an induced subgraph isomorphic to the crown graph $Cr(2p+1)$, then the Betti table is as shown in Figure~\ref{Fig:BettiShapes2}.

\begin{figure}[h]
\centering
$\begin{array}{r|ccccccccc}
      & 0 & 1 & \cdots & 2p & 2p+1 & \cdots & n-2p-1 & n-2p \\
\hline
\text{total:} & \beta_{0} & \beta_{1} & \cdots & \beta_{2p} & \beta_{2p+1} & \cdots & \beta_{n-2p-1} & \beta_{n-2p} \\
\hline
2p:   & \beta_{0, 2p} & \beta_{1, 2p+1} & \cdots & \beta_{2p, 4p} & \beta_{2p+1, 4p+1} & \cdots & \beta_{n-2p-1, n-1} & \beta_{n-2p, n} \\
2p+1: & 0 & 0 & \cdots & 0 & \beta_{2p+1, 4p+2} & \cdots & \beta_{n-2p-1, n} & 0 \\
\end{array}$
\caption{Betti table of $I_{\Delta}^{[p]}$ in the presence of $Cr(2p+1)$.}
\label{Fig:BettiShapes2}
\end{figure}
where \[
\beta_{2p+1,\,4p+2} (I_{\Delta}^{[p]}) =
p\,\bigl|\bigl\{\text{induced subgraphs of } \Delta \text{ isomorphic to } Cr(2p+1)\bigr\}\bigr|.
\]
\end{itemize}

\end{enumerate}
\end{Discussion}





\subsection{Some open questions.}\label{SubSec:open question} We conclude this work by presenting some open problems and directions for future research that naturally emerge from our results. To frame these questions, it is convenient to dually reformulate Theorem \ref{Thm: characterization linear resolution} in terms of graph theory as follows.

\begin{Theorem}\label{Thm: characterization linear resolution graph} Let $G$ be a graph such that $\overline{G}$ does not have triangles. Let $2 \leq p \leq \nu(G)$. The following conditions are equivalent: 
\begin{enumerate} 
\item $I(G)^{[p]}$ has linear resolution. 
\item $\overline{G}$ does not contain an induced subgraph isomorphic to the crown graph $Cr(2p+1)$ or to $K_{2+i,\,2p-i}$ for some even integer $i$ with $0 \leq i \leq p$. 
\end{enumerate} 
\end{Theorem}

The case in which a graph $G$ has a complement containing triangles is much more difficult to handle. Indeed, for example, if $|V(G)|\leq 6$ and $\overline{G}$ contains triangles, then the computations with \textit{Macaulay2} show that $I(G)^{[2]}$ has a linear resolution (over a field of characteristic zero) only when $\overline{G}$ coincides with one of the graphs shown in Figure \ref{fig:triangle_graphs}. This naturally arise the following problem.

\medskip
\textbf{Question 1.} Is it possible to give a combinatorial characterization of the graphs $G$, whose $\overline{G}$ contains triangles, such that $I(G)^{[p]}$ has linear resolution, for $2 \leq p \leq \nu(G)$?
\medskip

\begin{figure}[h]
\centering
\subfloat{\includegraphics[scale=0.4]{graph_triangle_1.eps}}\quad
\subfloat{\includegraphics[scale=0.4]{graph_triangle_2.eps}}\quad
\subfloat{\includegraphics[scale=0.4]{graph_triangle_3.eps}}\quad
\subfloat{\includegraphics[scale=0.4]{graph_triangle_4.eps}}\quad
\subfloat{\includegraphics[scale=0.4]{graph_triangle_5.eps}}
\caption{Some graphs.}
\label{fig:triangle_graphs}
\end{figure}


It is worth noting that in Example \ref{exa1}, the complement of $G_1$ consists of an isolated vertex and the leftmost graph in Figure \ref{fig:triangle_graphs}. However, since $\beta_{1,6}=2$, an analogue of Theorem \ref{Thm: betti 2,2p+2} does not hold.

\medskip
\textbf{Question 2.} Can we provide a combinatorial characterization of the graded Betti numbers in the lowest homological degree where $I(G)^{[p]}$ fails to be linearly related or to have a linear resolution, for all graphs $G$ whose $\overline{G}$ contains triangles, for $2 \leq p \leq \nu(G)$?

\medskip

On the other hand, we showed in Corollary \ref{Coro: lin res for p then lin res for p+1} that if a graph $G$ has triangle-free complement, then the existence of a squarefree power with a linear resolution implies that all higher squarefree powers also have linear resolutions. Moreover, computations carried out with \textit{Macaulay2} show that for all graphs $G$ with $|V(G)|\leq 8$ whose complements contain triangles (approximately $10{,}000$ graphs), if $I(G)^{[2]}$ has a linear resolution (over a field of characteristic zero), then every higher squarefree power also has a linear resolution. This leads to the following question.

\medskip
\textbf{Question 3.} Is it true that if $G$ is a graphs whose $\overline{G}$ contains triangles and $I(G)^{[p]}$ has a linear resolution for some $2 \leq p \leq \nu(G)$, then $I(G)^{[q]}$ has a linear resolution for all $p \leq q \leq \nu(G)$?

\medskip



\begin{footnotesize}
{\bf Acknowledgments.} Francesco Navarra and Ayesha Asloob Qureshi were supported by Scientific and Technological Research Council of Turkey T\"UB\.{I}TAK under the Grant No: 124F113. Francesco Navarra is a member of GNSAGA Indam and he acknowledges their support. Experiments with the computer algebra software \textit{Macaulay2} \cite{M2} have provided numerous valuable insights; the authors are grateful to the HPC group at Sabanci University for installing \textit{Macaulay2} on the Tosun cluster and for providing access to such a valuable computational resource. \\

{\bf Declaration of competing interest.}  The authors declare that they have no known competing financial interests or personal relationships that could have appeared to influence the work reported in this paper.\\

{\bf Data availability.} No data was used for the research described in the article.  
\end{footnotesize}

\bibliography{./ref.bib}{}
\bibliographystyle{alpha}	
\end{document}


